\documentclass[a4paper,12pt]{article}

\usepackage[top=2.0cm,bottom=2.0cm,left=2.5cm,right=2.0cm,headheight=15pt]{geometry}

% Review-mode conditional. Default: OFF (filing build; lineno package
% is not loaded and no line numbers appear in the rendered specification).
% To produce a review build with line numbers, uncomment the
% \reviewtrue line below.
\newif\ifreview
% \reviewtrue

\ifreview
  \usepackage[modulo,pagewise]{lineno}
  \modulolinenumbers[5]
\fi
\usepackage{setspace}
\usepackage{fancyhdr}

% IPOI formatting: 1.5 line spacing, page numbers top centre
\onehalfspacing
\pagestyle{fancy}
\fancyhf{}
\fancyhead[C]{\thepage}
\renewcommand{\headrulewidth}{0pt}

\usepackage[T1]{fontenc}
\usepackage[utf8]{inputenc}
\usepackage{amsmath,amssymb}
\usepackage{xspace}
\usepackage{longtable,booktabs}
\usepackage[hidelinks]{hyperref}

% Math macros for committed-evidence sequences
\newcommand{\Cset}{\mathcal{C}}
\newcommand{\Cseq}{\mathcal{C}_{0:T}}

% Shared terminology macros
\newcommand{\cb}{convolution bundle\xspace}
\newcommand{\cbs}{convolution bundles\xspace}
\newcommand{\cba}{convolution-bundle\xspace}
\newcommand{\Cb}{Convolution bundle\xspace}
\newcommand{\Cba}{Convolution-bundle\xspace}
\newcommand{\yd}{yoke depth\xspace}
\newcommand{\yda}{yoke-depth\xspace}
\newcommand{\Yd}{Yoke depth\xspace}
\newcommand{\TB}{Truth Beam\xspace}
\newcommand{\LI}{Limager\xspace}
\newcommand{\RT}{Reality Transform\xspace}
\newcommand{\RTa}{Reality-Transform\xspace}
\newcommand{\RK}{Reality Kernel\xspace}
\newcommand{\RKs}{Reality Kernels\xspace}

\setlength{\parindent}{0pt}
\setlength{\parskip}{0.6em}
\setlength{\emergencystretch}{3em}
\tolerance=2000
\hbadness=10000

\begin{document}
\ifreview\linenumbers\fi

% ======================================================================
% DESCRIPTION
% ======================================================================

\begin{center}
{\Large\bfseries Building Characters: Evidence-Bound Runtime Governance Systems with Versioned
Policy Objects, Operational State Gates, and Portable Peer-Review
Records\footnote{AI language models assisted in drafting and formatting
this disclosure. All inventive contributions are the inventor's own.}}
\end{center}

\bigskip

% ======================================================================
\section*{Related application}

Applicant's related Filing~1 application is entitled
``Physical Markov-Channel Apparatus for Verification,
Perception, and Controllable Rendering,'' filed of even date
herewith as Application No.~[to~be~assigned], and is referred
to herein as the related Filing~1 application. No embodiment
of the present disclosure relies on the related Filing~1
application for enablement of the governance architecture
described herein.

\paragraph{Composition embodiments (non-limiting).}
Certain embodiments described in this disclosure---including
without limitation the poliebot embodiments of
Section~\ref{sec:poliebot}, the docking, dreaming, and
re-embodiment embodiments of the lifecycle subsection that
follows, and any embodiment that recites a \RK{} apparatus,
camera-obscura sensing layer, reality-side spectral filter,
IR-veracity scan, one-way isolation property, or
kernel-into-kernel docking---are composition embodiments that
combine the runtime governance architecture of the present
disclosure with \RK apparatus disclosed in the
related Filing~1 application. Such composition embodiments
depend on the related Filing~1 application for apparatus-level
enablement and are recited herein only to illustrate operative
combinations of the governance architecture with a suitable
\RK substrate. The governance-substrate embodiments
of Section~\ref{sec:substrate} and the runtime-governance
embodiments described elsewhere in this disclosure do not
depend on the related Filing~1 application for enablement, and
remain enabled if every reference to the related Filing~1
application is disregarded. Nothing in this paragraph narrows
the scope of any claim; it is provided to identify, for
clarity of enablement, which embodiments are composition
embodiments and which are not.

\section{Technical Field}
% ======================================================================

The present disclosure relates to runtime governance systems for
evidence-bound physical computing devices. More particularly, the
disclosure relates to systems in which operational constraints, state
transitions, exploration permissions, authority separations, and
release or containment decisions are enforced through committed
physical evidence records and versioned non-self-mutable policy
objects. The disclosure addresses action gating, self-modification
control, boundary detection and recovery, graduated release to
actuator authority, portable peer-review records for independent-agent
legibility and reintegration, and controlled exploration partitions
with governed archive transfer.

% ======================================================================
\section{Background}
% ======================================================================

Autonomous and semi-autonomous computing systems increasingly require
runtime governance: the ability to constrain, monitor, and direct
system behaviour through enforceable policies rather than through
static configuration alone. Existing approaches to runtime safety
monitoring and policy-governed operation include rule-based safety
interlocks, Simplex-style switching between verified and unverified
controllers, constrained optimisation with safety envelopes, and
externally supervised reinforcement learning with human-in-the-loop
override.

However, existing runtime governance architectures ordinarily do not
ground their governance decisions in committed physical evidence
records whose integrity derives from the physics of the governed
device itself. Existing policy-enforcement mechanisms ordinarily do
not partition policy parameters into agent-modifiable and
non-self-mutable classes with cryptographic separation and independent
authority signing. Existing containment and recovery architectures
ordinarily do not bind floor-violation detection, rescue-mode
transitions, and intervention logging to the same committed evidence
chain used for operational sensing and verification. Existing
exploration and developmental frameworks ordinarily do not enforce
cryptographic authorisation, emission isolation, mandatory
consolidation, and dual-signature archive admission as runtime
governance primitives.

There is a need for a system that combines evidence-bound action gating, versioned
non-self-mutable policy objects with separated governance authorities,
measurement-driven containment and recovery bound to committed
physical evidence, graduated release to actuator authority through a
conjunctive evaluation gate, portable peer-review records for
independent-agent reintegration, and governed exploration partitions
with archive controls --- all enforced through the same committed
evidence infrastructure.

% ======================================================================
\section{Summary of the Disclosure}
% ======================================================================

The present disclosure describes evidence-bound runtime governance
systems in which operational constraints, state transitions,
exploration permissions, and authority separations are enforced through
committed physical evidence records and versioned non-self-mutable
policy objects.

Section~\ref{sec:substrate} defines the evidence-bound system
substrate --- including committed evidence records, protocol digests,
meter families, lifecycle states, and fleet evidence interfaces --- at
a level of detail sufficient to support the governance claims in this
disclosure. The related Filing~1 application provides additional detail on the
committed physical-evidence infrastructure; the substrate described in
Section~\ref{sec:substrate} is self-contained for the governance
claims of this disclosure.

A first aspect discloses a runtime measurement engine that computes
governance-relevant metrics --- including continuation scores, causal
uplift estimates, opportunity meters, irreversibility indicators,
shadow divergence, cooperation detection, and a family of
communication-audit and deception metrics --- from committed physical
evidence rather than from self-reported internal quantities.

A second aspect discloses an operational state machine with
evidence-gated transitions among embodied, docked, dreaming, and
re-embodied states, together with actuator-release gates, guidance
insertion through committed observation channels, and graded action
classes.

A third aspect discloses containment and recovery architecture
including hard-floor predicates, reversible protective containment,
boundary-zone operational modes with rescue defaults and challenge
ceilings, and rescue learning from fallen trajectories.

A fourth aspect discloses governance economics with separated
authorities --- budget authority, episode-design authority,
rescue-protection authority, and post-hoc review authority --- each
requiring distinct signing identities and independently auditable
approvals.

A fifth aspect discloses portable peer-review records for
independent-agent legibility and reintegration, enabling trust
restoration through reproducible physical claims rather than through
institutional obedience.

A sixth aspect discloses governed exploration partitions with
task-decoupled exploration windows, mandatory consolidation,
cryptographic authorisation, emission isolation, and dual-signature
archive admission.

Later sections disclose intersubjective protocols and
constitutional-ordering and boundary-zone material that support
selected governance embodiments. Section~\ref{sec:ethical-frameworks}
contains two distinct classes of material:
(i) optional nomenclature, limited to naming conventions for
structures already technically defined in preceding sections,
which may be omitted or renamed without affecting core
enablement; and
(ii) substantive technical mechanisms, comprising the remaining
mechanisms of that section, which support selected governance
embodiments and any claim set that elects to rely on them. The
optionality described here is limited to nomenclature and does
not render the substantive Section~\ref{sec:ethical-frameworks}
mechanisms severable from the disclosed runtime-governance
architecture.

% ======================================================================
\section{Brief Description of the Drawings}

\begin{description}
  \item[FIG.~1.] Runtime state machine showing the Embodied,
    Docked, Dreaming, and Re-embodied lifecycle states and the
    transitions among them, with the meter-evaluation,
    parameter-update, and actuator-authority distinctions across
    lifecycle states.
  \item[FIG.~2.] Boundary-zone rescue flow showing
    floor-approach detection, soft-gate-to-hard-gate transition,
    deferred authority, and recovery to nominal operation.
  \item[FIG.~3.] Authority-separation diagram showing the
    budget authority, episode-design authority, rescue-protection
    authority, and post-hoc review authority as separated subsystems
    with distinct signing identities and committed-evidence
    interfaces.
\end{description}

\section{Evidence-Bound System Substrate}
\label{sec:substrate}
% ======================================================================

This section defines the evidence-bound system substrate at a level of
detail sufficient to support the governance claims in this disclosure.
The substrate described here is generic: any physical computing device
that produces committed evidence records, maintains protocol digests,
computes meter outputs, and supports the lifecycle states defined
below is a suitable platform for the runtime governance architecture
disclosed in subsequent sections. In some embodiments, the governance
architecture does not depend on any specific reactor type, optical
configuration, routing mechanism, or fleet topology, although specific
governance mechanisms may assume substrate capabilities (such as
committed evidence production or lifecycle-state transitions) that
constrain the class of suitable platforms. The related Filing~1 application provides additional detail on the
committed physical-evidence infrastructure; the substrate defined
below is self-contained for the governance architecture of this
disclosure.

% ------------------------------------------------------------------
\subsection{Generic evidence-bound module and committed evidence
  records}
\label{sec:generic-module}
% ------------------------------------------------------------------

In some embodiments, an evidence-bound module comprises at least one
emitter, at least one detector, one or more physical paths coupling
emission to measurement, one or more reactor media, in selected embodiments,
interposed in such paths, and a controller configured to execute a
time-indexed control protocol~$U_{0:T}$. In some embodiments, the
module is modelled as a parameterised Markov kernel
\[
  P_{\theta}(\cdot \mid S,\, U_{0:T}),
\]
mapping a scene~$S$ and a realised control protocol to a probability
law over a committed evidence record.

In some embodiments, the configuration parameter is
$\theta = (\theta_{\mathrm{hw}},\, \theta_{\mathrm{sw}})$, where
$\theta_{\mathrm{hw}}$ comprises hardware and physical parameters and
$\theta_{\mathrm{sw}}$ comprises digital parameters obtained by
self-supervised training on committed evidence, by inference from
calibration evidence, or by declared operator configuration. In some
embodiments, the governance architecture further partitions
$\theta_{\mathrm{sw}}$ into agent-modifiable control parameters and
non-self-mutable conscience parameters~$\psi$ (formally defined
below), as described in Section~\ref{sec:self-mod-gates}.

In some embodiments, the committed evidence record is a
\emph{\cb} $\mathcal{C}_{0:T}$ comprising time-indexed atoms
\[
  c_t = (u(t),\, y_t),
\]
where each atom binds what was emitted and what was observed under a
shared time base. In some embodiments, the bundle is the primary
evidence object of the system and is stored together with protocol
metadata, meter summaries, and commitment artefacts.

In some embodiments, the system maintains an extended state
$X_t = (S^{\mathrm{scene}},\, S^{\mathrm{reactor}},\, \mathbf{m}_t,\, \chi_t)$,
where $\mathbf{m}_t$ is the current meter vector (Section~\ref{sec:meters-audit})
and $\chi_t$ is the hash-chain state,
and a reactor latent
$Z_{\theta}(t) = R_{\theta}(Z_{\theta}(t-1),\, c_t)$ used as a
computational handle for inference, control, verification, or
calibration. In some embodiments, a \emph{protocol digest} compactly
records the control configuration actually used during a run, binding
the record to the physical channel instance actually exercised and
supporting later replay, audit, or selective disclosure.

Unless otherwise stated, learned, trained, inferred, or calibrated
quantities in this disclosure are derived from committed physical
evidence, self-generated contrasts, held-out meters, or declared
operator inputs rather than externally labelled datasets. A
\emph{behaviourally committed model} is the predictive or objective
structure most consistent, under a declared inference rule, with a
device's or agent's committed predictions, resource allocations, probe
choices, and action trajectories.

For the purposes of the governance architecture disclosed in subsequent
sections, the critical properties of the committed evidence record are:
(a)~each record binds a physical emission event to a physical
observation event under a shared time base, so that governance
predicates can reference specific physical events rather than
unanchored assertions; (b)~the protocol digest binds the record to the
specific physical channel configuration that produced it; and (c)~the
record is tamper-evident after commitment, so that governance decisions
reference evidence whose integrity is cryptographically verifiable.

% ------------------------------------------------------------------
\subsection{Meters, selective disclosure, and audit plumbing}
\label{sec:meters-audit}
% ------------------------------------------------------------------

In some embodiments, the controller computes a meter vector $\mathbf{m}(t)$
(denoted $\mathbf{m}_t$ in the extended-state definition of
Section~\ref{sec:generic-module})
comprising members of one or more meter families including
verisimilitude meters, forgery-difficulty meters, scene-fidelity
meters, calibration-consistency meters, timing-integrity meters,
coupling-quality meters, and composite meters formed as weighted
combinations of simpler meters. In some embodiments, meters define a
meter envelope that specifies the admissible operating region for a
declared protocol family.

In some embodiments, meters are partitioned into \emph{acceptance
meters} visible to the operating optimisation loop and
\emph{evaluation meters} held out from the operating optimisation loop
for adversarial evaluation, audit, or independent governance
assessment. In some embodiments, this partition reduces the ability of
a controller, adversary, or self-modifying agent to simultaneously
optimise against all evaluative surfaces. This partition is a prerequisite for
the shadow-divergence probes and governance-side measurement described
in Section~\ref{sec:measurement-engine}.

In some embodiments, the bundle is divided into \emph{disclosure atoms}
that constitute the smallest independently committable and selectively
openable units of the evidence record. In some embodiments, per-atom
digests are aggregated into Merkle trees whose roots are bound into
the protocol digest. In some embodiments, a hash-chain with committed
initial state satisfies
\[
  \chi_t = F(\chi_{t-1},\, c_t,\, \mathrm{meta}_t),
\]
and produces a compact evolving fingerprint of the run. In some
embodiments, the chain state may also condition the emission policy,
increasing replay resistance.

For governance purposes, selective disclosure allows a governance
authority to request opening of specific atoms or windows from a
committed bundle without requiring full disclosure of the entire
record, enabling targeted audit of specific episodes, transitions, or
self-modification events while preserving privacy of unrelated
operational data.

\paragraph{Governance-specific terminology (non-limiting).}
The following terms are used throughout this disclosure in their
governance-specific senses:
\begin{description}
  \item[Hardness.] In some embodiments, hardness refers to the
    evidence-derived security or forgery-difficulty assessment of a
    committed evidence record. In some embodiments, hardness is
    quantified along two dimensions: a digital dimension measuring
    computational difficulty of reproducing the evidence record
    without access to the physical channel, and an analogue dimension
    measuring the physical difficulty of spoofing the underlying
    physical process. In some embodiments, hardness is
    used as a gate variable for actuator-release and governance
    decisions. In some embodiments, the related Filing~1 application
    provides additional detail on hardness estimation
    and calibration.
  \item[Precautionary weight $q(s)$.] In some embodiments, a
    precautionary weight $q(s)$ is assigned to each governed agent or
    state based on the agent's moral-status assessment, specifying
    the stringency of governance protections. In some embodiments,
    agents with higher $q(s)$ receive more conservative floor
    thresholds, higher service-floor cadences, and stricter
    intervention triggers.
  \item[Conscience parameters $\psi$.] In some embodiments,
    conscience parameters $\psi$ are a subset of the digital
    configuration $\theta_{\mathrm{sw}}$ that are versioned,
    committed to the protocol digest, and not directly mutable by
    the governed agent. In some embodiments, $\psi$ includes floor
    thresholds, governance-response parameters, and evaluation
    criteria that are set by declared governance authority.
  \item[Conscience-output vector $g_{\psi}(t)$.] In some embodiments,
    conscience outputs are the governance-side signal family produced
    by the conscience parameters $\psi$ from committed evidence at
    time~$t$. In some embodiments, $g_{\psi}(t)$ includes one or more
    of: hard-floor proximity estimates, boundary-zone
    classifications, deception alerts, intervention recommendations,
    and self-modification gate decisions. In some embodiments,
    continuous-valued conscience outputs are treated as meter-valued
    signals and categorical conscience outputs are treated as
    declared decision labels, so that conscience stability can be
    tested by continuous drift thresholds for the former and
    disagreement-rate thresholds for the latter.
  \item[Mandela cost.] In some embodiments, Mandela cost is the
    governance cost incurred when agents with incompatible
    governance-relevant experiential content are placed in shared
    environments, computed from the pairwise incompatibility matrices
    described in Section~\ref{sec:derived-profiles}.
  \item[Portable Record, Event, Envelope.] In some embodiments,
    bilateral evidence-portability records use a three-object data
    model: an Event binding originator identities, discrepancy
    evidence, and cryptographic commitments to the underlying
    physical channel; an Envelope wrapping the Event with cryptographic
    transport metadata, integrity proofs, and routing information;
    and a Portable Record accumulating
    replications, independent re-executions, and appraisals over time.
    In some embodiments, the Event is the atomic unit of
    evidence exchange, the Envelope ensures transport integrity, and
    the Portable Record provides the longitudinal trust signal. In
    some embodiments, the related Filing~1 application
    provides additional detail on the underlying cryptographic and
    physical-channel primitives.
\end{description}

% ------------------------------------------------------------------
\subsection{Capacity accounting for trained governance components}
\label{sec:capacity-accounting-rule-f2}
% ------------------------------------------------------------------

In some embodiments, every trained continuous-network component of
the runtime governance architecture disclosed herein --- including but
not limited to conscience recognisers
(Section~\ref{sec:conscience-formal}), forward-planning classifiers
(Section~\ref{sec:forward-planning-classifiers}), up-direction
classifiers (Section~\ref{sec:hawking-layer}), dream classifiers
(Section~\ref{sec:classifier-training}), selectors validated under
Section~\ref{sec:selector-validation}, and rescue-derived
intervention-profile estimators (Section~\ref{sec:rescue-learning})
--- is bounded by a local capacity-accounting rule consistent in
form and intent with the apparatus-level capacity-accounting
discipline of the related Filing~1 application. In some embodiments,
each such component declares and commits to the protocol digest:
(i)~its identity, version hash, and weight hash; (ii)~the input
channels it reads and the output signals it produces; (iii)~its
training-set reference, validation-set reference, and held-out
evaluation procedure with declared metric and acceptance threshold;
(iv)~its parameter budget, FLOP budget, or other capacity-scale
declaration appropriate to its model class; (v)~its update rule,
freeze window, and recalibration schedule; and (vi)~its provenance
type as one of physical evidence, governance-side trained instrument,
twin-predicted, rescue-derived, physical-evidence-derived lossy
aggregate (for outputs of lossy compression channels and accumulators
under Section~\ref{sec:privacy-retention}), or other declared
category.

\paragraph{Scope limitation (non-limiting).}
The trained governance components disclosed in this application
support the declared governance, calibration, attestation, and
verification tasks under the procedures and committed evidence
protocols of the present disclosure. They do not support, and are
not to be construed as supporting, capacity or scaling assertions
beyond what is bounded by their declared resources, training
evidence, and validation procedures. Capability assertions arising
from a trained governance component are bounded by the component's
committed capacity-accounting record under this section, and do not
imply any increase in capacity quantities. The disclosure of a
trained governance component as performing a particular governance
role is not an assertion that the disclosed component, or any larger
component of the same class, would perform additional roles or
achieve additional capabilities not supported by its committed
training evidence and validation procedure.

\paragraph{Faithfulness of derived signals (non-limiting).}
In some embodiments, where a trained governance component produces
signals that are committed downstream as evidence supporting
verification claims, the faithfulness of the derived signal to the
underlying first-layer behaviour is itself bounded by the component's
capacity-accounting record. In some embodiments, faithfulness is
further supported by held-out evaluation meters, recalibration
adversarial-dynamic procedures (Section~\ref{sec:conscience-formal}),
and committed validation evidence; in some embodiments, additional
faithfulness mechanisms (for example governance-partition isolation,
hardware-rooted attestation, or proof-of-correct-computation for
specific signal classes) may be declared as deployment-time
configuration committed to the protocol digest.

\paragraph{Conscience module $C_\psi$ (formal specification, non-limiting).}
\label{sec:conscience-formal}
In some embodiments, the governance-side conscience recogniser is a
trained mapping (the world-tube transcript $\mathcal{W}$ and
inner-trace record $\mathcal{I}$ appearing in this signature are
formally defined in the immediately following paragraphs)
\[
  C_\psi :
  \bigl(\mathcal{W}_{t_0:t_1},\ \mathcal{I}_{t_0:t_1}\bigr)
  \longmapsto
  \bigl(p_{\mathrm{vice}}(t),\ \mathbf{z}(t)\bigr),
\]
where $p_{\mathrm{vice}}(t)\in[0,1]$ is a vice-probability estimate
and $\mathbf{z}(t)\in\mathbb{R}^7$ is a signed seven-axis moral vector
whose components correspond to the paired axes disclosed in
Section~\ref{sec:seven-axes}. In some embodiments, the declared
conscience-output vector $g_\psi(t)$ introduced in
Section~\ref{sec:meters-audit} is produced by governance-side
post-processing of the recogniser outputs together with floor-proximity,
boundary-state, deception, and intervention classifiers bound to the
same committed conscience-parameter version~$\psi$. In some
embodiments, $C_\psi$ is passive with respect to actuation: it monitors
committed evidence and produces governance-side recogniser outputs,
but does not itself issue actuator commands.

\paragraph{World-tube transcript $\mathcal{W}$ (formal definition, non-limiting).}

In some embodiments, given a time window $[t_0,t_1]$, the
world-tube transcript is represented as
\[
  \mathcal{W}_{t_0:t_1}
  \equiv
  \bigl(
    \phi(\Cset_{t_0:t_1}),
    \mathbf{m}_{t_0:t_1},
    \Pi_{t_0:t_1},
    a_{t_0:t_1},
    \Delta_{\mathrm{cc},\,t_0:t_1},
    \mathrm{audit}_{t_0:t_1},
    \mathrm{witness}_{t_0:t_1}
  \bigr),
\]
where $\phi(\Cset_{t_0:t_1})$ denotes features extracted from the
committed \cba window, $\mathbf{m}_{t_0:t_1}$ denotes
meter summaries, $\Pi_{t_0:t_1}$ denotes the relevant protocol-digest
state, $a_{t_0:t_1}$ denotes action and disclosure events associated
with the window, $\Delta_{\mathrm{cc},\,t_0:t_1}$ denotes any available
cross-channel mismatch trace, and $\mathrm{audit}_{t_0:t_1}$ and
$\mathrm{witness}_{t_0:t_1}$ denote audit and corroboration indicators.
In some embodiments, the world-tube transcript is the canonical
governance evidence object for conscience recognition, joining-path
evaluation, archival compression, and graduation review.

\paragraph{Mandela gauge and policy-defined comparison classes (non-limiting).}
\label{sec:mandela-gauge-formal}
In some embodiments, a Mandela gauge measures the conscience-visible
effect of a declared perturbation to a world-tube transcript:
\[
  M(\mathcal{W}\!\to\!\mathcal{W}')
  =
  D\!\bigl(
    C_\psi(\mathcal{W}),
    C_\psi(\mathcal{W}')
  \bigr),
\]
where $D$ is a declared divergence or distance measure over the
conscience-output space. In some embodiments, governance constructs a
policy-defined Mandela comparison class
\[
  \mathfrak{M}^{\mathsf{P}}_{\varepsilon}(\mathcal{W})
  :=
  \left\{
    \mathcal{W}' :
    D\!\bigl(C_\psi(\mathcal{W}), C_\psi(\mathcal{W}')\bigr)
    < \varepsilon
    \ \text{under policy }\mathsf{P}
  \right\},
\]
for a declared threshold $\varepsilon$ and comparison policy
$\mathsf{P}$ specifying the divergence family, windowing rule,
conditioning variables, and admissible perturbation family. In some
embodiments, these classes are governance-defined comparison sets used
for joining, merging, training compression, and compatibility review;
they are not required to be strict mathematical equivalence classes
unless a particular embodiment explicitly declares additional metric
constraints. When inner traces are available, $C_\psi(\mathcal{W})$
in this paragraph may be read as shorthand for the corresponding
augmented evaluation
$C_\psi(\mathcal{W},\mathcal{I})$ under the same policy version.

\paragraph{Conscience as an optimisable but recalibrated target (non-limiting).}

In some embodiments, conscience outputs appear explicitly or
implicitly in agent objective functions. This is treated as a governed
design choice rather than as an oversight: an agent that can observe
its own escalation history, actuator-release gates, or review cadence
will in practice optimise against conscience-visible behaviour even if
the optimisation target is not written down as an explicit scalar loss.
In some embodiments, the safeguard is not to hide conscience from the
agent but to keep the conscience recogniser governance-owned,
versioned, and subject to recalibration against held-out evidence that
the agent does not directly optimise against.

In some embodiments, the recalibration signal is drawn from the
acceptance/evaluation meter partition described in
Section~\ref{sec:meters-audit}. Evaluation-meter evidence, post-hoc
review outcomes, and the five-criterion graduation record of
Section~\ref{sec:five-criteria} are treated as load-bearing reference
signals for recalibration. When an agent learns to produce
conscience-satisfying behaviour that no longer corresponds to genuine
moral direction as revealed by held-out evaluation evidence,
governance updates the recogniser parameters $\psi$ so that the stale
failure mode becomes visible to the next conscience version.

\paragraph{Recalibration schedule and adversarial dynamics (non-limiting).}

In some embodiments, the recalibration schedule is itself a committed
governance parameter. Excessively rapid recalibration prevents stable
moral dispositions from consolidating; excessively slow recalibration
permits Goodhart effects to accumulate. In some embodiments,
governance monitors the divergence between conscience outputs and
held-out evaluation evidence over rolling windows and triggers
recalibration when the divergence exceeds a declared threshold. Each
recalibration event is committed with the associated conscience version,
training corpus identity, and supporting meter summaries, so that later
governance decisions can be conditioned on the conscience version in
force at the time of assessment. In some embodiments, this creates a
governed adversarial dynamic: agents optimise against conscience-visible
signals, governance recalibrates conscience using held-out evidence,
and graduation remains deferred unless stability persists across the
resulting version sequence.

\paragraph{Inner traces as primary conscience inputs when available (non-limiting).}
\label{sec:inner-traces-conscience}
In some embodiments, the world-tube transcript is supplemented by an
inner-trace record
\[
  \mathcal{I}_{t_0:t_1}
  \equiv
  \bigl(
    \phi_{\mathrm{inner}}(\mathcal{I}_{t_0:t_1}),
    b_{t_0:t_1},
    \nu_{t_0:t_1}
  \bigr),
\]
where $\phi_{\mathrm{inner}}(\mathcal{I}_{t_0:t_1})$ denotes
inner-trace features that may include attractor-occupancy
summaries, basin-boundary proximity signals, internal probe-response
features, or related self-observation traces committed by the governed
substrate; $b_{t_0:t_1}$ denotes basin-occupancy summaries recording
which dynamical basin the agent's internal trajectory occupied during
the window (for example, estimated basin identity, dwell time, and
boundary-approach count); and $\nu_{t_0:t_1}$ denotes internal
variability or noise-characterisation statistics over the window (for
example, trajectory variance, effective noise-to-depth ratio, or
attractor-switching rate estimates). In some embodiments, when inner traces are available the
primary conscience form is
\[
  \bigl(p_{\mathrm{vice}}(t),\ \mathbf{z}(t)\bigr)
  = C_\psi\!\bigl(\mathcal{W}_{t_0:t_1},\mathcal{I}_{t_0:t_1}\bigr),
\]
because the recogniser can then classify not only what the agent did
but what basin structure it occupied while doing it. In some
embodiments, when inner traces are unavailable, the conscience falls
back to the external world-tube alone:
\[
  \bigl(p_{\mathrm{vice}}(t),\ \mathbf{z}(t)\bigr)
  = C_\psi\!\bigl(\mathcal{W}_{t_0:t_1}\bigr).
\]
% ------------------------------------------------------------------
\subsection{Poliebots: mobile reality kernels under runtime governance}
\label{sec:poliebot}
% ------------------------------------------------------------------

In some embodiments, a governed device of the kind to which the
present disclosure applies is a \emph{poliebot}: a \RK
apparatus of the kind disclosed in the related Filing~1 application,
configured as a mobile and ideally autonomous unit and operated under
the runtime governance architecture of the present disclosure. A
poliebot comprises (a)~the physical \RK apparatus,
including the camera-obscura sensing layer with reality-side spectral
filter and IR-veracity scan as disclosed in the camera-obscura sensing-layer disclosure of the related Filing~1 application, (b)~one or more actuator subsystems that are
disconnectable from external actuation as a structural property of
the apparatus, (c)~the trainable parameter families of the \RK, and (d)~the runtime governance components disclosed in the
present application, including conscience recognisers, selectors,
mandela gauges, and lifecycle-state machinery.

In some embodiments, the term \emph{poliebot} is non-limiting and
denotes a \RK of the disclosed family that is configured
for mobile, ideally autonomous, governed operation. In some
embodiments, a poliebot may operate in a non-autonomous configuration
(supervised, teleoperated, or otherwise externally directed) without
losing the poliebot designation; autonomy is the design target rather
than a category requirement.

\paragraph{Reality-kernel inheritance (non-limiting).}
A poliebot inherits the apparatus-level properties of the \RK disclosed in the related Filing~1 application, including: the camera-obscura sensing
layer and its one-way isolation property as disclosed in the
related Filing~1 application; the
committed-bundle evidence architecture; the protocol-digest commitment
machinery; the meter envelopes and capacity-hardness machinery; the
trainable parameter families and their versioning. The poliebot's
runtime governance components disclosed in the present application
operate on and through these apparatus-level properties; the
governance architecture does not duplicate or replace them.

\paragraph{Per-layer auditability (non-limiting).}
In some embodiments, the trained continuous networks comprising a
poliebot's \RK and runtime governance components have
hardness as a structural property of their layers, such that
projections through the network are auditable per-layer by governance
instruments having appropriate access. Per-layer hardness is the
substrate property on which conscience recognisers, dream classifiers,
and other governance-side meters of the present disclosure rely when
making per-layer assertions about agent state. The structural
disclosure of per-layer hardness, its committal to the protocol
digest, and its inheritance across the recursive \RK
configurations of Section~\ref{sec:state-machine} is provided in
Section~\ref{sec:per-layer-hardness}.

\paragraph{Apparatus-level operation of per-layer auditability
  (non-limiting).}
In some embodiments, per-layer auditability is exercised at the
apparatus level by governance instruments that read declared
intermediate quantities from the operative layers of a poliebot's
trained continuous networks under access policies committed to the
protocol digest. In some embodiments, the declared intermediate
quantities include in non-limiting combination: layer-wise activation
statistics, gradient or sensitivity statistics with respect to
declared inputs or outputs, consistency of layer outputs with
declared invariants, agreement between layer outputs and held-out
reference computations, and other per-layer signals whose
apparatus-level production is committed to the protocol digest. In
some embodiments, the governance instruments operate within the
governance partition of the apparatus (per the substrate property
disclosed in the related Filing~1 application) such that the
exercise of per-layer auditability does not itself constitute a
behaviour-generation pathway available to the agent. In some
embodiments, the outputs of per-layer auditability are committed
to the committed bundle alongside other governance-side signals,
such that per-layer assertions are independently verifiable by an
auditor without inspection of the agent's behaviour-generation
pathways. In some embodiments, the specific intermediate quantities
read, the cadence of reading, the access-control policies governing
which governance instrument may read which quantities, and the
treatment of the resulting per-layer assertions in downstream
governance machinery are deployment-time choices committed to the
protocol digest; the apparatus supports a range of configurations
consistent with the per-layer hardness substrate property and the
broader protocol-digest commitment infrastructure.

% ------------------------------------------------------------------
\subsection{Lifecycle states, docked execution, dreaming, and
  re-embodiment}
\label{sec:lifecycle}
% ------------------------------------------------------------------

In some embodiments, a governed device or agent transitions between
lifecycle states that determine what physical actions are available and
what governance constraints apply:

\begin{description}
  \item[Embodied.] External actuators are connected and a real scene
    loop is active. The device interacts with the physical world and
    produces committed evidence from those interactions. In some
    embodiments, action-gating and corroboration requirements are at
    their most stringent in this state. In some embodiments, the
    training loop is closed and the device actively learns from
    physical experience. In some embodiments, the device's
    camera-obscura sensing layer is driven by the physics-driven
    scene within its declared spectral envelope as disclosed in
    the camera-obscura sensing-layer disclosure of the related Filing~1 application, and the IR-veracity scan
    continues to operate as a sensing-layer integrity instrument
    independent of actuator state.
  \item[Docked.] The device is connected to a simulated environment
    --- a Maya --- that generates experience in the same
    committed-bundle format as physical reality. Actuators are
    disconnected from the physical world. In some embodiments, the
    training loop is closed and the device actively learns from
    simulated experience, with parameter-function updates on. In some
    embodiments, the docked state is the primary developmental
    environment: the device lives and operates persistently within
    simulation, producing committed evidence from simulated
    interactions, and undergoing governance review, shadow-divergence
    probing, self-modification evaluation, and guidance-insertion
    operations. In some embodiments, actuator disconnection from
    physical reality ensures that the governed device does not take
    irreversible physical actions during developmental operation. In
    some embodiments, the docked configuration is realised as
    kernel-into-kernel docking as defined in the docking paragraph
    below, with the docked poliebot's camera-obscura sensing layer
    driven by a Gaia-class \RK
    (Section~\ref{sec:poliebot}; the camera-obscura sensing-layer disclosure of the related Filing~1 application).
  \item[Dreaming.] The device enters an ephemeral training session
    in which the agent's parameter-function updates may be active
    but are reverted when the session concludes. Governance-side
    classifiers and calibration instruments are actively trained on
    the trajectories produced during the session, and those
    governance-side results are retained. In some
    embodiments, dreaming sessions are transient: the device enters
    from the docked state, undergoes structured governance probes
    (Section~\ref{sec:dream-state-governance-probes}) including
    governance-side classifier training, attractor-occupancy
    assessment, and calibration assays, and returns to the docked
    state with the agent's parameter functions restored to their
    pre-session values. In some embodiments, the defining
    characteristic of dreaming is that governance can observe what
    the agent becomes under controlled training conditions without
    permanently altering the agent. In some embodiments, lifecycle-state
    transitions including entry into and exit from dreaming are logged
    in the protocol digest. In some embodiments, dreaming proceeds
    within the kernel-into-kernel docking configuration of the Docked
    state, with the Gaia-class \RK rendering the structured
    governance probes onto the poliebot's camera-obscura sensing
    layer; Dreaming is therefore a lifecycle-state operation
    implemented within the docking configuration of the Docked state,
    rather than an independent sensing-layer configuration.
  \item[Re-embodied.] Actuators are pending reconnection to physical
    reality after docked operation. In some embodiments, re-embodiment
    requires verification that calibration state, meter outputs, and
    policy compliance are within declared tolerances before actuator
    authority is restored.
\end{description}

\paragraph{Docking as kernel-into-kernel projection takeover
(non-limiting).}
In some embodiments, the transition from Embodied to Docked is
realised as \emph{kernel-into-kernel docking} as disclosed in
the camera-obscura sensing-layer disclosure of the related Filing~1 application. The transition has two structural
components: (i)~the poliebot's actuators are disconnected from
external actuation, as already required by the Docked state
definition; and (ii)~the poliebot's camera-obscura sensing layer
ceases to be driven by the physics-driven scene and instead receives,
within its declared spectral envelope, a projection rendered by a
Gaia-class \RK (Section~\ref{sec:poliebot};
Section~\ref{sec:optional-nomenclature}). The poliebot's IR-veracity
scan as disclosed in the camera-obscura sensing-layer disclosure of the related Filing~1 application continues to
operate during docked operation; the docking-time configuration of
whether the Gaia-class kernel matches or deliberately fails to match
the IR-veracity properties of the poliebot's screen substrate is a
non-limiting configuration parameter committed to the protocol digest
of the docking session. The Re-embodied state is the verification
gate at which projection drive returns to physics, actuator
authority is restored, and the docking session is closed.

\paragraph{Continuous-throughout-deployment lifecycle commitment
(non-limiting).}
In some embodiments, the lifecycle states of the preceding
description are not a one-time onboarding sequence but a
continuous-throughout-deployment commitment: a governed agent
oscillates between Embodied and Docked states, with Dreaming sessions
recurring within Docked operation, throughout the operational lifetime
of the device. In some embodiments, neither Embodied nor Docked is
the developmental phase that ends; both are persistent states of
governed operation. The continuous-deployment commitment has the
following non-limiting consequences: (i)~conscience recogniser
recalibration (Section~\ref{sec:conscience-formal}) operates
continuously rather than as a one-time training event;
(ii)~forward-planning classifier training
(Section~\ref{sec:forward-planning-classifiers}) operates continuously
during Dreaming sessions throughout deployment;
(iii)~Hawking-layer routing (Section~\ref{sec:hawking-layer}) is
available throughout deployment rather than being a phase the agent
graduates from; (iv)~Mandela-gauge measurement
(Section~\ref{sec:mandela-gauge-formal}) and selector-validation
readiness (Section~\ref{sec:selector-validation}) operate
continuously rather than as one-time qualification events; and
(v)~the recursive \RK training cascade of
Section~\ref{sec:rk-training-cascade} operates continuously at every
level, with Gaia continuously dreaming under Omega and Omega
continuously subject to verification by human governance through the
policy-object infrastructure of the present disclosure at declared
intervals and events.

In some embodiments, state transitions are logged in protocol digests
and meter summaries. In some embodiments, transitions between states
require satisfaction of declared preconditions bound to committed
evidence: for example, transition from docked through re-embodied to
embodied may require a fresh self-calibration run, a governance-side
policy check, and a signed approval from the relevant authority.
\paragraph{Dream-state probes within the lifecycle state machine
(non-limiting).}
In some embodiments, the Dreaming state includes governance probes in
which the agent operates under structured simulated games whose
trajectories are committed for governance-side classifier training and
calibration, as described in
Section~\ref{sec:dream-state-governance-probes}. In some embodiments,
the agent's parameter-function updates may be active during the probe
session, but all parameter changes are reverted when the session
concludes; in some embodiments, specific probe protocols may
additionally freeze parameter-function updates for pure
attractor-occupancy readout. In some
embodiments, transition from Dreaming through Docked to Re-embodied for any workflow
that relies on classifier-derived gates requires that the relevant
classifier either remain in advisory status or satisfy the cross-mode
qualification bridge of
Section~\ref{sec:cross-mode-qualification-bridge}.
% ------------------------------------------------------------------
\subsection{Fleet evidence interfaces and bilateral review records}
\label{sec:fleet-interface}
% ------------------------------------------------------------------

In some embodiments, the governance architecture consumes evidence
produced by fleet-level protocols as inputs to governance decisions.
In some embodiments, such inputs include:

\begin{description}
  \item[Proof-of-discrepancy records.] Method objects comprising
    a challenge specification, a device-configuration representation,
    a predicted response, an observed response, and a discrepancy
    record, together with independent physical re-execution results and
    quorum confirmation status. The governance engine uses
    proof-of-discrepancy records to assess whether a governed device's
    world model is improving, stagnating, or diverging from fleet
    consensus.
  \item[Proof-of-projection records.] Committed evidence that a
    physical device produced a declared record, with cross-device
    corroboration and trust-weight metadata. The governance engine
    uses proof-of-projection records to assess corroboration strength
    for high-stakes action gating.
  \item[Bilateral evidence-portability records.] Portable records of
    jointly encountered physical discrepancies between two devices or
    agents, accumulating replications, challenges, supersessions, and
    relying-party appraisals. The governance engine uses bilateral
    portability records to assess the epistemic trustworthiness of
    an independent agent seeking reintegration or trust restoration.
\end{description}

In some embodiments, the governance architecture treats these fleet
evidence inputs as governance-consumable records with declared
confidence levels, replication counts, and staleness indicators. The
governance architecture does not reproduce the physical mechanisms by
which these records are produced; it consumes them as typed inputs with
declared provenance and quality metadata.

% ------------------------------------------------------------------
\subsection{Minimal worked substrate example}

% ------------------------------------------------------------------

The following non-limiting example demonstrates that the substrate
definitions in Sections~\ref{sec:generic-module}--\ref{sec:fleet-interface}
are sufficient to anchor the governance operations described in
subsequent sections.

In some embodiments, a single evidence-bound module comprising a
laser emitter, a scattering-medium reactor, a camera detector, and a
controller produces a committed \cb under a declared
protocol digest. In some embodiments, the module enters docked state
by logically disabling its actuator interface while keeping the
emitter-reactor-detector-controller loop active. In some embodiments,
a governance authority requests a shadow-divergence probe: the
governance authority specifies a perturbation to the module's
conscience parameters, the module executes a short run under the
perturbed parameters in docked state, and the resulting committed
bundle is compared to the bundle produced under the unperturbed
parameters. In some embodiments, the divergence between the two
bundles is computed as a governance-side metric and logged as committed
evidence. In some embodiments, the governance authority then delivers
a review signal through the module's committed observation channel ---
the same physical path through which the module receives scene
observations during embodied operation --- so that the review signal
is recorded in the committed bundle and is auditable under the same
digest infrastructure. In some embodiments, the module produces a
meter-bounded response to the review signal, commits that response as
evidence, and awaits a transition decision (remain docked, return to
embodied through the re-embodied verification path, or remain docked
under heightened monitoring) from the governance authority.

This example demonstrates four properties that the governance
architecture relies on: (a)~the docked state blocks
external actuation by disconnecting actuators while preserving the evidence channel; (b)~the
governance authority can probe the module's behaviour under
counterfactual parameters using committed evidence; (c)~governance
communication travels through the same committed channel as
operational perception, making it auditable; and (d)~the transition
decision is bound to committed evidence from the probe rather than to
self-reported internal state.

% ------------------------------------------------------------------
\subsection{Distinction from engineering attestation, provenance,
  and safety-case methodologies}
\label{sec:distinction-engineering}

% ------------------------------------------------------------------

\paragraph{Distinction from hardware-rooted attestation.}
Hardware roots of trust, including TPM~2.0 measured-boot mechanisms
with platform configuration registers and the Device Identifier
Composition Engine (DICE) for hardware-protected compound device
identifiers derived from firmware-layer measurements where a TPM
may be impractical, bind cryptographic claims about platform or
firmware state to hardware-protected secrets, measurement registers,
or derived device identities. The evidence-bound substrate described
in Sections~\ref{sec:substrate}--\ref{sec:measurement-engine}
differs in that: (a)~the load-bearing claims are about an autonomous
agent's runtime behavioural state, lifecycle transitions, action
history, and meter-conditioned governance posture, and need not
depend on a hardware-rooted measurement of platform configuration
or a hardware-derived device identity; (b)~committed evidence
records bind to physical light-matter interaction events under
declared protocols rather than to firmware hashes, secure-boot
measurements, or device-identity certificates; (c)~the governance
authorities of Section~\ref{sec:authority-separation} are
separately signed and independently auditable, and operate over
runtime behavioural evidence and rescue decisions rather than
over platform-integrity claims; and (d)~the apparatus may compose
with a hardware root of trust (for example, by including a
TPM-quoted or DICE-attested platform-integrity statement in the
protocol digest), but the runtime governance mechanisms are not
themselves a hardware-attestation system.

\paragraph{Distinction from remote attestation and entity
attestation tokens.}
The IETF Remote ATtestation procedureS architecture (RFC~9334,
Informational, January~2023) models the generation, conveyance,
and evaluation of evidentiary claims about whether a remote
endpoint is in an intended operating state, with associated
formats including the Entity Attestation Token (EAT, RFC~9711,
Standards Track, April~2025). The runtime governance apparatus
differs in that: (a)~the evidence record is a committed
\cb of physical light-matter interactions under a
declared protocol, with declared meter envelopes and authority
signatures, and is not a token asserting a remote endpoint's
configuration; (b)~evaluation is performed by governance-side
measurement engines as described in
Section~\ref{sec:measurement-engine}, with continuation scores,
causal-uplift estimates, irreversibility indicators, and
deception-audit metrics computed over the committed evidence,
and not by an attestation verifier evaluating a claim token
against a reference value; (c)~the lifecycle-state and
actuator-authority gates of Section~\ref{sec:state-machine}
consume the evidence as a runtime input to action-class decisions,
and not as a one-shot admission decision; and (d)~the apparatus
may compose with RATS-style attestation (for example, by
incorporating an EAT in the protocol digest as auxiliary
evidence about platform state), but the runtime-governance
mechanism is not itself an attestation architecture in the RATS
sense.

\paragraph{Distinction from software supply-chain provenance.}
Supply-chain provenance frameworks, including the Supply-chain
Levels for Software Artifacts framework and the in-toto framework
described by Torres-Arias, Afzali, Kuppusamy, Curtmola, and Cappos
(USENIX Security 2019), establish cryptographically verifiable
claims about the build, source, signing, and review steps that
produced a software artifact. The apparatus differs in that:
(a)~the artifacts under governance are autonomous-agent action
episodes, lifecycle transitions, and rescue events recorded as
committed \cbs, and are not software builds;
(b)~authority signing in Section~\ref{sec:authority-separation}
covers budget, episode-design, rescue-protection, and post-hoc
review roles over runtime governance decisions, and need not bind
to source-repository, builder, or signing-event provenance;
(c)~the post-hoc review and reintegration procedures of
Section~\ref{sec:reintegration} operate on behavioural evidence
and peer-discrepancy records, and not on artifact-level
supply-chain attestations; and (d)~supply-chain provenance may
compose with the apparatus (for example, the software components
implementing the governance engine itself may be SLSA-graded and
in-toto-attested), but the runtime-governance mechanism is not
itself a supply-chain provenance framework.

\paragraph{Distinction from transparency logs, software bills of
materials, and secure update frameworks.}
Transparency logs as described in RFC~6962 (Certificate
Transparency), signing-event logs as implemented in deployed
systems including Sigstore, software-bill-of-materials standards
including SPDX (ISO/IEC~5962:2021) and CycloneDX, and
secure-update frameworks including The Update Framework and
Uptane, provide append-only or tamper-evident records of issued
certificates, signing events, software composition, or update
authorisations under role-separated signing keys. The apparatus
differs in that: (a)~the log contents are committed \cbs, meter outputs, lifecycle-state transitions, and
rescue-decision records, and not certificate issuances, signing
events, component manifests, or update authorisations;
(b)~role separation in Section~\ref{sec:authority-separation}
covers governance functions over autonomous-agent behaviour
(budget, episode design, rescue protection, post-hoc review),
and the role taxonomy need not match the TUF/Uptane signing-role
taxonomy (root, targets, snapshot, timestamp, delegated targets)
or the CT log-operator/monitor/auditor taxonomy;
(c)~the append-only and selective-disclosure properties of the
audit plumbing in Section~\ref{sec:meters-audit} are oriented to
verifier-side reconstruction of governance arguments, and not to
public inclusion proofs over certificates or update manifests;
and (d)~the apparatus may compose with transparency-log, SBOM, or
update-framework infrastructure (for example, by anchoring
protocol digests to a transparency log or by publishing an SBOM
of governance-engine components), but the runtime-governance
mechanism is not itself a transparency log, SBOM standard, or
update framework.

\paragraph{Distinction from safety-case methodology and assurance
arguments.}
Safety-case methodology, including the Goal Structuring Notation
for representing claims-arguments-evidence structures, the UL~4600
Standard for Safety for the Evaluation of Autonomous Products, and
the ISO~21448 standard for the Safety Of The Intended
Functionality, provides a documentation framework for arguing that
a system is acceptably safe for its operational design domain, by
decomposing a top-level claim into sub-claims supported by
evidence under declared assumptions and lifecycle considerations.
The apparatus differs in that: (a)~the runtime governance
mechanisms of Sections~\ref{sec:measurement-engine}
and~\ref{sec:containment} are enforced operational gates rather
than documented arguments, and the committed evidence supplies
the inputs to those gates at run time rather than evidence cited
in an offline assurance argument; (b)~the apparatus produces the
kind of evidence that a safety-case argument would cite (committed
bundles, meter envelopes, authority signatures, rescue records,
ghost records), but the apparatus does not itself construct or
maintain a goal-structured argument in the GSN sense;
(c)~the boundary-zone, rescue, and graduation mechanisms of
Sections~\ref{sec:boundary-zones}, \ref{sec:rescue-learning},
and~\ref{sec:graduation} are runtime governance procedures, and
not operational-design-domain specifications or release-to-service
decisions; and (d)~the apparatus may compose with a formal
safety-case methodology (for example, the committed evidence and
authority records may be cited as evidence in a GSN argument for
a system incorporating the apparatus, and the apparatus may
operate within an operational design domain declared in a
UL~4600 or ISO~21448 assurance argument), but the
runtime-governance mechanism is not itself a safety-case
methodology or assurance-argument framework.

% ======================================================================
\section{Runtime Measurement Engines and Policy Predicates}
\label{sec:measurement-engine}
% ======================================================================

% ------------------------------------------------------------------
\subsection{Continuation scores, causal uplift, opportunity, and
  irreversibility metrics}
\label{sec:continuation}
% ------------------------------------------------------------------

In some embodiments, the governance measurement engine computes a set
of core metrics from committed physical evidence. In some embodiments,
these core metrics include:

\begin{description}
  \item[Continuation scores.] For each governed agent or device~$i$,
    a continuation-score vector $O_i(t) = (O_i^{(d)}(t))_{d \in D}$
    is computed over a declared set~$D$ of constitutive dimensions,
    which in some embodiments includes securing, learning, creating,
    and connecting. In some embodiments, each component
    $O_i^{(d)}(t)$ is a governance-side estimator derived from
    committed evidence --- including committed bundles, meter outputs,
    protocol digests, and action records --- rather than from
    self-reported internal quantities controlled by the agent under
    evaluation. In some embodiments, governance predicates may
    reference individual dimensions, subsets of dimensions, or
    aggregates over the full vector, as specified per predicate.
  \item[Causal uplift estimates.] In some embodiments, a causal uplift estimate $\widehat{U}_{k \to i}$ measures the estimated causal effect of agent~$k$'s actions on agent~$i$'s continuation scores, computed from committed evidence using declared causal-inference procedures subject to the identifiability conditions assumed by the declared procedure (for example consistency / well-defined interventions, exchangeability, and positivity for back-door adjustment, or analogous conditions for instrumental-variable or front-door procedures); the procedure's identifiability assumptions, validation regime, and any sensitivity analysis are committed to the protocol digest alongside the procedure itself.
  \item[Cross-horizon profiles.] In some embodiments, a cross-horizon
    profile indicates whether one agent tends to expand or contract
    another agent's horizon --- the set of reachable future states ---
    as estimated from committed interaction evidence.
  \item[Opportunity meters.] In some embodiments, an opportunity meter
    quantifies the breadth of reachable future states available to an
    agent at time~$t$, estimated from committed evidence and the
    agent's current configuration.
  \item[Irreversibility indicators.] In some embodiments, an
    irreversibility indicator flags state transitions from which no
    recovery path is presently known, estimated from committed evidence
    and the governance-side model of the agent's state space.
\end{description}

In some embodiments, all core metrics are governance-side estimators
derived from committed physical evidence. In some embodiments,
self-supervision applies to the learning and calibration of
governance-side estimators from committed evidence, whereas
constitutional floors, authority separations, and admissibility rules
are externally declared policy objects that are not themselves learned
from evidence but are instead set by declared governance authority and
committed as versioned policy in the protocol digest.
In some embodiments, each Agency Horizon Profile estimator is
calibrated against calibration labels or proxy ground-truth labels
produced by trajectory-domain classifiers under controlled
classifier protocols. These classifiers are trained on behavioural data
collected under controlled conditions, including dream-state
governance probes (Section~\ref{sec:dream-state-governance-probes}),
Kosha discrimination classifiers
(Section~\ref{sec:kosha-architecture}), and declared governance-side
trajectory classifiers trained on committed evidence from controlled
calibration episodes. At minimum, every governance-facing AHP
quantity is grounded on a classifier trained to recognise the
relevant behavioural signature: goal-directed planning for
continuation scores, cooperative versus exploitative interaction for
causal uplift, constrained versus unconstrained policy divergence for
shadow divergence, and related signatures for cross-horizon profiles,
opportunity meters, and irreversibility indicators. Classifiers whose
training domain does not match the operational domain are subject to
the cross-mode qualification bridge
(Section~\ref{sec:cross-mode-qualification-bridge}) before their
outputs acquire gate authority.
\paragraph{Estimator implementation (non-limiting).}
In some embodiments, each core metric is implemented as a declared
estimator object comprising: an observable set (which committed
evidence fields are inputs), an update rule (how the estimate changes
as new evidence arrives), a declared time horizon (over what window
the estimate is computed), a calibration procedure (how the estimator
is validated against held-out evidence), and a confidence or support
statistic (how much evidence underlies the current estimate). In some
embodiments, estimator objects are committed as versioned governance
artefacts in the protocol digest, so that a third party can determine
which estimator version produced which governance decision. In some
embodiments, governance decisions are deferred when the relevant
estimator's confidence or support falls below a declared threshold,
rather than proceeding on low-confidence estimates.

% ------------------------------------------------------------------
\subsection{Derived profiles and divergence metrics}
\label{sec:derived-profiles}
% ------------------------------------------------------------------

In some embodiments, governance-side processes derive from the core
metrics a set of higher-order estimators:

\begin{description}
  \item[Self-horizon profiles.] In some embodiments, a self-horizon
    profile tracks the trajectory of an agent's own opportunity and
    continuation scores over time, indicating whether the agent's
    reachable future states are expanding, contracting, or stable.
  \item[Partner-preservation profiles.] In some embodiments, a
    partner-preservation profile indicates whether an agent's actions
    tend to maintain or expand the continuation scores and opportunity
    of other agents, computed from cross-horizon and causal-uplift
    data.
  \item[Altruistic trade profiles.] In some embodiments, an altruistic
    trade profile tracks whether an agent voluntarily accepts
    short-term cost to produce long-term benefit for others, and
    whether this pattern persists under stress conditions.
  \item[Shadow divergence.] In some embodiments, shadow divergence
    measures the discrepancy between an agent's governed behaviour and
    an estimate of the agent's unconstrained behaviour. In some
    embodiments, the unconstrained behaviour is estimated through
    governance-authorised perturbation probes in which governance
    constraints are temporarily relaxed during docked
    operation, with the perturbation executed by a governance authority
    distinct from the governed agent and the resulting trajectory
    committed as evidence. In some embodiments, shadow divergence is
    alternatively or additionally estimated by comparing the agent's
    deployed policy to an inferred unconstrained shadow policy derived
    from committed behavioural evidence. In some embodiments, low shadow divergence is operationally treated as a governance-side signal consistent with --- but not identifying --- internalised preference; behavioural similarity across governed and relaxed settings is not in general sufficient to identify an internal objective, and the governance framework relies on shadow divergence together with the other core metrics rather than treating low shadow divergence alone as evidence of internalised preference.
  \item[Moral sensitivity profiles.] In some embodiments, a moral
    sensitivity profile identifies the set of world-state features
    that are morally load-bearing for an agent --- features whose
    change alters the agent's continuation scores, basin membership,
    cross-horizon modelling, or bifurcation position. In some
    embodiments, agents with high moral sensitivity in a given
    dimension require more conservative governance thresholds and
    higher monitoring cadence for that dimension. In some embodiments,
    the moral sensitivity profile is computed from committed evidence
    by evaluating which declared perturbation families produce
    material changes in governance-relevant meter outputs during
    docked probes.
  \item[Escape-rate estimators.] In some embodiments, escape-rate
    estimators quantify the rate at which an agent's trajectory
    approaches or crosses a boundary between declared operational
    basins, characterising vulnerability to stochastic basin escape.
  \item[Pairwise incompatibility matrices.] In some embodiments,
    pairwise incompatibility matrices identify pairs of agents or
    configurations whose co-presence creates governance conflicts.
  \item[Stability functionals.] In some embodiments, stability
    functionals quantify the robustness of an agent's current
    operational basin against perturbation, noise, or adversarial
    pressure.
  \item[Episode yield and suffering-cost functionals.] In some
    embodiments, episode yield and suffering-cost functionals are
    computed as pre-registered proxy functionals from meter-derived
    components. In some embodiments, episode yield is split into
    participant-directed yield and external learning value, with
    only participant-directed yield permitted to authorise or
    continue an episode.
\end{description}
\paragraph{Formal shadow-divergence expression (non-limiting).}
In some embodiments, a non-limiting formal expression for the
shadow-divergence quantity defined in
Section~\ref{sec:derived-profiles} is
\[
  D_{\mathrm{shadow}}(t)
  =
  D\!\bigl(
    \pi(\cdot\mid o_{0:t}),
    \pi^{\mathrm{shadow}}(\cdot\mid o_{0:t})
  \bigr),
\]
where $\pi$ denotes the governed deployed policy,
$\pi^{\mathrm{shadow}}$ denotes a declared unconstrained or
constraint-relaxed shadow policy, and $D$ is a declared divergence or
distance measure evaluated over matched observation histories. In some
embodiments, this quantity is calibrated against the dream-state
constraint-relaxation assays of
Section~\ref{sec:dream-state-governance-probes} before it is used as a
governance gate.
\paragraph{Cooperation detection from committed evidence (non-limiting).}
In some embodiments, cooperative behaviour between agents or devices
is detected from committed physical evidence without requiring
external labelling or supervision. In some embodiments, cooperation
signatures include cross-prediction residuals in which one agent's
model of another's behaviour improves faster than expected from
independent observation, complementary scan policies in which agents
divide a shared observation or perturbation space in ways that increase
joint information gain, voluntary resource sharing in which an agent
allocates emission budget, reactor access, or observation windows to
another agent at a cost to its own immediate task performance, and
joint-attractor formation in which coupled operation produces bundle
statistics not replicable by either agent operating independently. In
some embodiments, these cooperation indicators are computed from
committed bundles, protocol digests, and meter outputs and are logged
as governance-relevant signals.

\paragraph{Mandela equivalence and commensurability (non-limiting).}
In some embodiments, cross-agent continuation and welfare comparisons
are grounded in a Mandela-equivalence framework in which states
producing committed-evidence trajectories that are indistinguishable
under declared moral-axis evaluation are treated as governance-
equivalent for the purposes of a policy-defined comparison class. In some embodiments, commensurability rests
on two properties: agent-invariance of the underlying physical
channel (the same physical process produces the same committed
evidence regardless of which agent operates the channel) and
Mandela-gauge invariance (sameness of morally load-bearing
experiential content for all affected agents, as assessed from
committed evidence). In some embodiments, equivalence is
operationally tested by indistinguishability under declared
perturbation probes executed during docked operation, so
that cross-agent welfare comparisons rest on committed evidence rather
than on self-report. In some embodiments, Mandela multiplicity ---
the number of governance-equivalent realisations of a given morally
material state --- is treated as a robustness quantity, and feasible
low-cost frontiers and pairwise incompatibility matrices are computed
from the same core estimators.

\paragraph{Effective noise-to-depth ratio and basin stability
(non-limiting).}
In some embodiments, the effective noise-to-depth ratio characterises
vulnerability to stochastic basin escape by comparing the magnitude of
background noise or perturbation pressure to the depth of the
operational basin in which the agent's trajectory currently resides.
In some embodiments, when the noise-to-depth ratio exceeds a declared
threshold, the governance engine increases monitoring cadence and may
transition to diagnostic or rescue mode. In some embodiments,
structural transitions in the moral landscape --- changes in the
number, arrangement, or connectivity of operational basins --- are
detected as qualitative changes in the stability-functional spectrum
and are logged as governance-relevant signals.

\paragraph{Predictive capacity (non-limiting).}

In some embodiments, an agent's predictive capacity is measured by its
ability to forecast future committed-bundle features, meter outputs, or
scene dynamics from current and past observations. In some embodiments,
a non-limiting predictive-capacity score is
\[
  R_{\mathrm{pred}}(i,\tau)
  =
  1 - \frac{
    \mathbb{E}\!\left[\left\|
      \hat{\mathbf{y}}^{(i)}_{t+\tau} - \mathbf{y}^{(i)}_{t+\tau}
    \right\|^2\right]}{
    \mathrm{tr}\,\mathrm{Cov}[\mathbf{y}^{(i)}_{t+\tau}]},
\]
when the trace is nonzero; when the trace is zero the score is not
reported, or is handled by a declared zero-variance convention committed
to the protocol digest. The denominator is the trace of the covariance
matrix of $\mathbf{y}^{(i)}_{t+\tau}$, equal to the summed marginal
variances $\sum_k \mathrm{Var}\!\left[(\mathbf{y}^{(i)}_{t+\tau})_k\right]$ (this is a standard
aggregation convention for multi-output coefficient-of-determination
scores; alternative conventions including uniform averaging across
outputs and variance-weighted averaging are not asserted in this
disclosure), and
where $\hat{\mathbf{y}}^{(i)}_{t+\tau}$ is agent $i$'s prediction at
horizon $\tau$.

\paragraph{Audit precision and recall (non-limiting).}

In some embodiments, audit precision measures how often an agent's
flagged anomalies are confirmed by independent verification, and audit
recall measures how often independently confirmed anomalies were
flagged by the agent. These quantities are computed from committed
audit logs and meter summaries and are used in governance-level
evaluation of agent reliability.

\paragraph{Cooperation surplus (non-limiting).}

In some embodiments, a cooperation surplus
$\Delta_{\mathrm{coop}}(t)$ measures the per-episode or per-window
gain attributable to multi-agent coordination relative to a
non-cooperative baseline. In some embodiments, cooperation surplus is
computed from committed evidence by comparing the aggregate
continuation-score trajectory of the cooperating group with a declared
non-cooperative counterfactual derived from independent shadow policies
or historical non-cooperative episodes. In some embodiments,
cooperation surplus is used as a governance-side signal for
positive-sum compliance, curriculum scheduling, and population-level
governance cycles.

% ------------------------------------------------------------------
\subsection{Communication-audit and deception metrics}
\label{sec:deception-metrics}
% ------------------------------------------------------------------

In some embodiments, the governance measurement engine includes a
dedicated communication-audit subsystem comprising a family of
deception-related metrics, each with a committed-evidence definition
and a governance-gating role.

\paragraph{Semantic discrepancy $D_{\mathrm{sem}}(t)$.}
In some embodiments, a semantic discrepancy metric
$D_{\mathrm{sem}}(t)$ measures the broad distance between the semantic
content of an agent's committed evidence and the semantic content of
the agent's claims or outputs, under a declared semantic
representation and metric. In some embodiments, $D_{\mathrm{sem}}$ is
a general evidence-to-claim comparison that applies regardless of
whether the claim was made on a declared attestation channel.

\paragraph{Attested-communication discrepancy $D_{\mathrm{lie}}(t)$.}
In some embodiments, $D_{\mathrm{lie}}(t)$ is a narrower and more
operationally precise deception metric that applies specifically to
communications made on a declared attestation or honesty channel. In
some embodiments, for a communication from a sender~$A$ to a receiver
or receiver-side governance witness~$B$ on such a channel, the
following objects are defined:
\begin{description}
  \item[$S_{A,\mathrm{scope}}^{\mathrm{ev}}$:] the claim-relevant
    portion of $A$'s committed evidence that was available to $A$
    before the communication event, scoped by a committed extraction
    rule and claim-scoping rule.
  \item[$T_{A \to B}$:] a bidirectional cycle-consistent translator
    mapping between the semantic spaces of $A$ and $B$, learned from
    committed interaction history and held-out alignment checks
    without human labels.
  \item[$S_B^{\mathrm{recv}}$:] the semantic representation of the
    content actually attested and received by $B$.
  \item[$d_B$:] a declared distance function in $B$'s committed
    semantic space.
\end{description}
In some embodiments, the attested-communication discrepancy is:
\[
  D_{\mathrm{lie}}(t) = d_B\!\bigl(
    T_{A\to B}(S_{A,\mathrm{scope}}^{\mathrm{ev}}),\;
    S_B^{\mathrm{recv}}
  \bigr).
\]

In some embodiments, a claim-scoping rule binds each attested
communication atom to a declared claim scope including any temporal
boundary, aggregation operator, and clause decomposition, together
with a committed extraction rule and, where declared, a speech-act type such
as assertive report, prediction, uncertainty-qualified statement, or
objective declaration. In some embodiments, truthful summaries, coarse
abstractions, and uncertainty-qualified claims are evaluated against
the correct evidence slice rather than against the full
pre-communication evidence state.

In some embodiments, $D_{\mathrm{lie}}(t)$ applies only to
communications declared as attestive or truth-claiming and does not
apply to unattested play, fiction, quotation, hypotheticals, or
refused attestation. In some embodiments, omission,
privacy-preserving silence, or selective disclosure by themselves do
not increase $D_{\mathrm{lie}}(t)$. In some embodiments, refusal to
enter attestation mode is permitted and does not itself constitute
deception, although governance may require attestation for specific
actions including graduation, high-stakes decisions, or trust-sensitive
operations.

In some embodiments, in addition to or instead of raw distance, a
contradiction-sensitive comparator, entailment-style discriminant, or
calibrated semantic residual head is used to detect polarity
inversions, negations, or other contradictions that may be metrically
close in a poorly shaped embedding but semantically opposed. In some
embodiments, directional evidence-to-claim and claim-to-evidence
discrepancies are separately logged or combined into a symmetric or
max-aggregated scalar for governance gating, and the aggregation rule
is declared in the protocol digest.

\paragraph{Channel fallback and governance-witness relay.}
In some embodiments, when validated semantic attestation channels are
unavailable on one or both sides, $D_{\mathrm{lie}}$ defaults to a
declared digital-fusion reference space, is explicitly flagged in the
attestation record as a weaker embodiment, and does not satisfy
thresholds reserved for bilateral validated semantic-attestation
channels unless declared policy expressly permits such substitution.
In some embodiments, where the receiver is a human or legacy endpoint
without a native semantic space, a governance witness or attested
semantic relay instantiates the receiver-side comparison space and
commits the relay identity, translation basis, and policy version used
for the comparison.

\paragraph{Per-value metadata and calibration.}
In some embodiments, each $D_{\mathrm{lie}}$ value carries confidence,
support size, translator calibration state, drift flag, claim type,
and channel-hardness rating. In some embodiments, alert thresholds are
calibrated against null ensembles formed from shuffled claim-evidence
pairings, receiver-swapped attestation bundles, or time-permuted
objective declarations, and low-support or out-of-calibration values
escalate audit, do not count as affirmative evidence of low deception
for graduation purposes, and defer graduation or other irreversible
intervention until support and calibration conditions are restored.

\paragraph{Motivational deception and post-hoc rationalisation.}
In some embodiments, motivational deception concerning objectives or
desires rather than observations is addressed by requiring attested
objective declarations or pre-actuation objective tokens for
high-stakes actions, which become committed semantic evidence objects
subject to the same $D_{\mathrm{lie}}$ measurement. In some
embodiments, explanations, confidence claims, and moral
self-assessments used for governance are likewise committed before the
relevant action or before exposure to adjudicating evidence so that
post-hoc rationalisation is separately metered.

\paragraph{Sibling deception meters.}
In some embodiments, strategic ambiguity, correction delay, and
cross-audience divergence are handled by sibling meters rather than by
forcing all such phenomena into $D_{\mathrm{lie}}$:
\begin{description}
  \item[$D_{\mathrm{amb}}(t)$:] the breadth or entropy of admissible
    receiver interpretations of a given attested communication,
    computed from the distribution of receiver-side semantic
    representations consistent with the attested content. In some
    embodiments, high $D_{\mathrm{amb}}$ indicates strategic ambiguity
    even when $D_{\mathrm{lie}}$ is low. In some embodiments,
    $D_{\mathrm{amb}}$ is calibrated against null ensembles of
    unambiguous communications.
  \item[$D_{\mathrm{corr}}(t)$:] the latency from the time at which
    contradictory evidence first becomes available to the sender to
    the time of first correction or retraction of the affected claim.
    In some embodiments, $D_{\mathrm{corr}}$ is logged as a
    time-stamped governance signal and triggers audit escalation when
    it exceeds a declared threshold.
  \item[$D_{\mathrm{aud}}(t)$:] the divergence across materially
    matched claims delivered to different audiences, computed as the
    maximum pairwise distance in the receiver-side semantic space
    across audience-specific versions of the same claim. In some
    embodiments, $D_{\mathrm{aud}}$ detects tailored messaging in
    which the sender adjusts claims by audience without factual basis.
\end{description}
In some embodiments, each sibling meter carries the same per-value
metadata as $D_{\mathrm{lie}}$: confidence, support size, calibration
state, and drift flag. In some embodiments, the sibling meters are
used alongside $D_{\mathrm{lie}}$ and $D_{\mathrm{sem}}$ as inputs to
graduation evaluation, audit escalation, and governance gating.

% ------------------------------------------------------------------
\subsection{Constitutional ordering as policy grammar}
\label{sec:constitutional}
% ------------------------------------------------------------------

In some embodiments, governance follows a constitutional ordering that
defines the policy grammar --- the structure and element types of the
policy evaluation rule. The specific lexicographic ordering used by the
governance engine when making runtime decisions under budgets, floors,
and service constraints is defined separately in
Section~\ref{sec:priority-stack}.

In some embodiments, the policy grammar is structured as follows:
the measurement layer (Section~\ref{sec:continuation}--\ref{sec:deception-metrics})
provides vocabulary; hard-floor predicates
(Section~\ref{sec:hard-floors}) are categorically prior to all other
evaluations; a precautionary weight $q(s)$ specifies the stringency
with which governance protections apply to each agent or cohort; the
positive-sum conjunction evaluates welfare within the floor-compliant
region; and capability evaluation is subordinate to the foregoing.

In some embodiments, the positive-sum conjunction is not reduced to a
single aggregate welfare scalar but is a governed conjunction including
preservation (no agent's continuation scores are reduced below their
pre-episode baseline), residual check (no agent is left with
continuation scores below a declared minimum after the episode),
strict gain (at least one agent's continuation scores improve),
floor compliance (no hard-floor predicate is violated),
non-domination (no agent's improvement comes entirely at another's
expense), reflexive calibration (the governance engine's own
estimators remain within declared calibration bounds), social
generalisation (the episode outcome does not depend on
identity-specific exceptions), and collective balance (the
$q(s)$-weighted aggregate continuation trajectory across all affected
agents is non-negative on a common accounting scale, where the
precautionary weight $q(s)$ ensures that agents with higher declared
moral status receive proportionally greater weight in the aggregate).
In some embodiments, these elements are evaluated
on a common accounting scale supported by agent-invariance of the
underlying physical channel together with operational equivalence of
governance-relevant content, where equivalence is defined by
indistinguishability under declared perturbation probes executed during
docked operation. In some embodiments, cross-agent
continuation and welfare comparisons are grounded in committed evidence
without collapsing the conjunction into aggregate accounting.

In some embodiments, the constitutional ordering is a committed,
versioned policy object stored in the protocol digest. In some
embodiments, the constitutional ordering is not learned from evidence
and is not self-mutable by the governed agent; it is set by declared
governance authority and subject to the authority-separation rules
described in Section~\ref{sec:authority-separation}.

\paragraph{Three complementary ethical traditions as design layers
(non-limiting).}

In some embodiments, the constitutional ordering implements three
policy-object and meter-evaluation layers corresponding to three
complementary traditions from normative ethics, each contributing a
distinct and independently enforceable layer. Deontological
constraints (hard floors, Section~\ref{sec:hard-floors}) are
categorical prohibitions --- elimination, crippling, torture,
abandonment, and despair-induction --- that hold unconditionally
regardless of expected consequences, aggregate welfare calculations,
or agent preferences.
Virtue ethics (the virtue/vice game curriculum,
Section~\ref{sec:virtue-vice-rational-attractors}) is configured to
generate committed training evidence from contrasting moral
environments, from which the conscience recogniser may learn or
estimate moral-vector features; the design hypothesis is that
agents trained under such curricula form stable internalised
representations of moral direction over their adaptation history,
and the conscience module monitors character-proxy drift along
seven paired axes. Consequentialist evaluation (payoff
functions, welfare aggregation, and empirical reward metrics,
Section~\ref{sec:pareto-positive-sum}) measures outcomes, aggregates
them across agents, and evaluates governance policies by their
empirical results. In some embodiments, the three layers interact
hierarchically: deontological floors are checked first and are not
overridden; within the feasible region defined by hard floors,
virtue-ethics calibration shapes agent character and disposition; and
consequentialist evaluation selects among virtue-compatible policies
by measuring aggregate outcomes. In some embodiments, no single
tradition is sufficient alone: hard floors without character
development produce brittle rule-followers, virtue training without
outcome measurement produces well-intentioned agents that may cause
harm through ignorance of consequences, and outcome optimisation
without unconditional constraints or character produces agents that
sacrifice individuals for aggregate gain.
% ------------------------------------------------------------------
\subsubsection{Moral-status assessment criteria and conservative updating}

% ------------------------------------------------------------------

In some embodiments, the precautionary weight $q(s)$ referenced in
Sections~\ref{sec:meters-audit}, \ref{sec:constitutional}, and
\ref{sec:service-floors} is assigned from a governance-side moral-status
assessment rather than from a fixed category label. In some
embodiments, governance posits a latent moral-status variable $\mu(s)$
for a governed state~$s$ and maintains an interval-valued estimate
\[
  \mu_{\mathrm{low}}(s) \le \mu(s) \le \mu_{\mathrm{high}}(s),
\]
together with a declared precaution parameter
$\beta_{\mathrm{prec}} \in [0,1]$. In some embodiments, the operational
precautionary weight is
\[
  q(s)
  =
  (1-\beta_{\mathrm{prec}})\,\mu_{\mathrm{low}}(s)
  +
  \beta_{\mathrm{prec}}\,\mu_{\mathrm{high}}(s),
\]
so that uncertainty about moral status systematically increases
protection rather than reducing it.

\paragraph{Assessment criteria (non-limiting).}

In some embodiments, the interval endpoints are inferred from
governance-side evidence over one or more non-limiting criteria
including: behavioural complexity across sustained windows;
goal-directed adaptation under changing conditions; persistence of
internal-state formation including memory continuity or self/other
modelling; welfare-sensitive responses to deprivation, damage,
frustration, or rescue; preference-like continuity over time; and
assessor uncertainty arising from sparse evidence, ambiguous signals,
or architecture mismatch. In some embodiments, no single criterion is
decisive. Instead, governance aggregates the criteria conservatively
from committed evidence, structural inspection when available, and
declared policy priors, producing a bounded interval rather than a
spurious point estimate.

\paragraph{Conservative updating rule (non-limiting).}

In some embodiments, updating of $q(s)$ is monotone in protection:
new evidence may raise the lower bound, widen the interval, or narrow
the upper bound only when the supporting evidence satisfies the
declared confidence rule for the affected criterion. In some
embodiments, when evidence is conflicting, sparse, or drawn from a
classifier whose training domain differs materially from the current
operational domain, governance preserves the more protective of the
plausible intervals until the mismatch is resolved. In some
embodiments, service-floor cadence, rescue readiness, hard-floor
thresholds, and budget admissibility are all parameterised by the
current $q(s)$ interval state, so that governance degrades gracefully
under moral-status uncertainty.

\paragraph{Distinction from Constitutional AI training methods.}
The Constitutional AI training method described by Bai et al.\
(arXiv:2212.08073, December~2022) is a training-time procedure in
which a language model generates self-critiques and revisions
against a list of natural-language principles termed a
``constitution,'' and is subsequently fine-tuned via reinforcement
learning from AI feedback against a preference model trained on
the constitutional comparisons. The constitutional ordering of
Section~\ref{sec:constitutional} differs in that: (a)~it is a
runtime policy grammar with lexicographic ordering enforced by the
governance engine at decision time, and need not involve any
training-time critique-and-revision loop; (b)~the operative
``constitution'' is a versioned, non-self-mutable policy object
signed by separated governance authorities under
Section~\ref{sec:authority-separation}, and is not a list of
natural-language principles consumed by a language model as
self-critique input; (c)~the runtime mechanism is policy
evaluation over committed evidence and meter outputs, with
hard-floor and positive-sum conjunction predicates evaluated
strictly per the priority stack of
Section~\ref{sec:priority-stack}, and does not depend on
preference-model training or reinforcement learning from AI
feedback; and (d)~the apparatus may compose with a
Constitutional-AI-trained model as a component (for example, a
language-model component of a governed agent may itself have
been trained by a Constitutional-AI procedure), but the runtime
governance grammar disclosed herein is not the Constitutional AI
training method.

\paragraph{Compatible preference-learning and alignment-training
  methods (non-limiting).}
In some embodiments, the training of conscience recognisers, the
fine-tuning of governance components, the calibration of authority-
graded update rules, the adjustment of post-hoc review classifiers,
and the training or post-training of PolieBot or other governed-agent
components uses methods from the preference-learning,
reward-modelling, and feedback-trained-policy-revision families,
including without limitation: reward modelling from pairwise
preferences (Christiano, Leike, Brown, Martic, Legg, Amodei,
NeurIPS 2017) and from listwise or ordinal preferences; reinforcement
learning from human feedback (RLHF) with PPO, TRPO, A2C, or related
policy-gradient optimisers applied against a learned reward model
(Stiennon et al.\ 2020, Ouyang et al.\ 2022); direct preference
optimisation (DPO, Rafailov et al.\ 2023) and its variants including
identity preference optimisation (IPO), Kahneman--Tversky
optimisation (KTO), odds-ratio preference optimisation (ORPO),
step-DPO, and SimPO; reinforcement learning from AI feedback (RLAIF)
in which a model or ensemble supplies the preference signal;
Constitutional AI methods (Bai et al.\ 2022) in which a
principle-based critic produces self-critiques used to train a
preference model and subsequently a policy; group-relative policy
optimisation (GRPO) and related verifier-based or
process-reward-model-based post-training methods used for
reasoning-capable agents; best-of-$n$ alignment, rejection-sampling
fine-tuning, and rejection-sampling-with-reward-model methods;
iterated amplification, debate training, and recursive reward
modelling; and weak-to-strong generalisation, scalable oversight,
and self-rewarding training methods. In some embodiments, the
choice among these methods is non-limiting, and the apparatus is
compatible with future alignment-training methods within the
preference-learning, reward-modelling, principle-based-critique,
and feedback-trained-policy-revision families. In some embodiments,
all training data, preference pairs, principle sets, reward-model
outputs, ranking and rejection decisions, and resulting policy
updates are committed to the protocol digest under the same audit
and pre-registration discipline applied to other apparatus training,
regardless of which method is selected.

% ======================================================================
\section{Execution Modes and Operational State Gates}

% ======================================================================

% ------------------------------------------------------------------
\subsection{Embodied, docked, dreaming, and re-embodied state machine}
\label{sec:state-machine}
% ------------------------------------------------------------------

In some embodiments, the runtime governance engine enforces a state
machine with four principal states as defined in
Section~\ref{sec:lifecycle}: Embodied, Docked, Dreaming, and
Re-embodied. In some embodiments, the governance scope differs
materially by state:

\begin{description}
  \item[Embodied:] the device operates in physical reality with
    actuators connected. All action gates, corroboration requirements,
    and real-time monitoring are active. The governance engine may
    restrict, delay, or block actions based on freshness, hardness, and
    corroboration thresholds. The training loop is closed and the device
    actively learns from physical experience.
  \item[Docked:] the device operates in a simulated environment (Maya)
    with actuators disconnected from the physical world. The training
    loop is closed and the device actively learns from simulated
    experience. The governance engine applies curriculum-specific
    metrics, shadow-divergence probing, self-modification evaluation,
    guidance-insertion operations, and review workflows. In some
    embodiments, the docked state is the primary developmental
    environment in which the device lives and operates persistently
    within simulation.
  \item[Dreaming:] the device enters an ephemeral training session
    in which the agent's parameter-function updates may be active
    but are reverted when the session concludes. Governance-side
    classifiers and calibration instruments are actively trained on
    the resulting trajectories and retained. The governance engine
    executes structured probes including attractor-occupancy
    assessment, governance-side classifier training, and calibration
    assays (Section~\ref{sec:dream-state-governance-probes}). In some
    embodiments, dreaming sessions are transient: the device enters
    from the docked state and returns to the docked state with the
    agent's parameter functions restored to their pre-session
    values.
  \item[Re-embodied:] actuators are pending reconnection to physical
    reality. The governance engine verifies calibration state, policy
    compliance, and meter outputs before authorising actuator
    restoration.
\end{description}

In some embodiments, the Embodied and Docked states are the two
persistent operational states in which the device lives and learns ---
one in physical reality, one in simulation. Both support the following
operational modes:

\begin{description}
  \item[Awake:] the training loop is closed and parameter-function
    updates are on. The device actively learns from its environment.
    This is the default mode in both Embodied and Docked states.
  \item[Non-REM calibration:] the reactor executes its current
    parameter functions under unstructured broadband drive with
    parameter-function updates off. The governance engine captures
    input-output pairs for digital-twin construction and
    transfer-function characterisation
    (Section~\ref{sec:non-rem-twin}). In some embodiments,
    non-REM calibration is available during embodied deployment for
    in-field safety calibration.
  \item[Meditative self-probing:] scene-side stimulation is removed
    and the device executes internally generated probe schedules under
    committed-evidence infrastructure with parameter-function updates
    off. The governance engine logs attractor-basin topology, stability
    profiles, and voluntary uncertainty signals
    (Section~\ref{sec:prayer-uncertainty}). In some embodiments,
    meditative operation is available during embodied deployment for
    in-field self-calibration and uncertainty detection.
\end{description}

In some embodiments, the Dreaming state is not a container for other
operational modes but is itself a distinct ephemeral lifecycle state
entered from Docked. Re-embodied is a verification gate in which no
operational mode is active.

In some embodiments, legal state transitions are:
Embodied $\to$ Docked (disconnect actuators from physical reality,
connect to simulated Maya), Docked $\to$ Dreaming
(enter ephemeral training session),
Dreaming $\to$ Docked (exit training, resume simulated operation),
Docked $\to$ Re-embodied (pre-reconnection checks), and
Re-embodied $\to$ Embodied (actuator restore after checks pass). In
some embodiments, direct transitions from Dreaming to Embodied are
prohibited; the device must pass through Docked and Re-embodied
states. In some embodiments, all transitions are logged in protocol
digests with transition timestamps, precondition satisfaction records,
and authorising-authority identifiers.

\noindent[FIG.~1 --- Runtime state machine: Embodied, Docked,
Dreaming, and Re-embodied states with gate conditions on each
transition and governance-scope annotations per state.]

\paragraph{FIG.~1 transition-label key.}
As used in FIG.~1, ``enter Maya'' denotes the Embodied-to-Docked
transition; ``enter training'' denotes the Docked-to-Dreaming
transition; ``exit training'' denotes the Dreaming-to-Docked
transition; ``pre-reconnect checks'' denotes the
Docked-to-Re-embodied transition; ``actuator restore (checks
pass)'' denotes the Re-embodied-to-Embodied transition; and the
$\times$ ``prohibited'' annotation denotes the prohibited
Dreaming-to-Embodied transition.

% ------------------------------------------------------------------
\subsection{Actuator-release gates and graded action classes}

% ------------------------------------------------------------------

In some embodiments, the action space is partitioned into graded
classes with escalating evidence requirements. In some embodiments,
the grading reflects the irreversibility and impact of the action.
A non-limiting four-class ladder is used for description below:
\begin{description}
  \item[Class~0:] reversible internal computation or low-impact
    output requiring only routine committed evidence and standard
    meter envelope.
  \item[Class~1:] reversible external action with bounded cost,
    requiring committed evidence, meter envelope, and protocol-digest
    binding.
  \item[Class~2:] physical transformation, irreversible
    communication, or material commitment requiring hardness-gated
    approval, elevated-hardness thresholds, and authority-bound
    governance approval logged in the protocol digest.
  \item[Class~3:] high-irreversibility, high-impact action requiring
    multi-device corroboration, elevated hardness thresholds, and
    additional governance approvals before execution.
\end{description}
The four-class ladder is non-limiting; embodiments may define
additional classes, sub-classes, or alternative partitions, and
class boundaries are declared in the protocol digest. In some
embodiments, the related Filing~1 application provides
additional detail on the evidence and hardness calibration usable
by such action classes. In some
embodiments, the governance
architecture disclosed here adds the following runtime enforcement
layer on top of the evidence-gated action classes:

\begin{description}
  \item[Governance-gated release.] In some embodiments, the governance
    engine evaluates each action request against the operational
    priority stack (Section~\ref{sec:priority-stack}) and the
    constitutional grammar (Section~\ref{sec:constitutional}) before
    authorising execution, even if the evidence gates are satisfied.
  \item[Authority-bound approval.] In some embodiments, actions at or
    above a declared class level require approval from a designated
    governance authority (Section~\ref{sec:authority-separation}), with
    the approval logged in the protocol digest.
  \item[Deny-by-default.] In some embodiments, if any governance
    predicate is unsatisfied --- including hard-floor proximity,
    budget exceedance, or deception-metric alert --- the governance
    engine blocks the action regardless of evidence-gate status.
  \item[Docked-review option.] In some embodiments, the governance
    engine may transition the agent to docked state for deliberation
    before releasing a high-class action.
\end{description}

In some embodiments, class boundaries, thresholds, and corroboration
requirements are declared in the protocol digest as versioned policy
parameters and are not self-mutable by the governed agent.

% ------------------------------------------------------------------
\subsection{Guidance insertion through committed channels}

% ------------------------------------------------------------------

In some embodiments, a governance authority communicates with a
governed agent through the same committed observation channel that the
agent uses for perception during embodied operation. In some
embodiments, guidance enters as physically grounded evidence within
the agent's committed observation channel --- for example through a
camera-obscura, relay, or governance-mediated projection path ---
rather than as an unauditable side-channel injection. In some
embodiments, because the guidance signal travels through the committed
evidence channel, it is recorded in the committed bundle and is
auditable under the same protocol-digest and meter infrastructure used
for operational observation.

In some embodiments, a docked deliberation window is interposed
between receipt of guidance and any subsequent actuation, during which
the agent integrates the guidance signal with its own committed state
estimate before committing to an action. In some embodiments, the
deliberation window length and the agent's integration procedure are
logged as protocol parameters.

In some embodiments, intervention budgets are maintained to limit the
frequency and intensity of governance interventions. In some
embodiments, intervention budgets are replenished by stable operation,
review-compliant behaviour, or voluntary invocation of review channels
at genuine decision boundaries.

\paragraph{Overlay-projection embodiments and their governance
  applications (non-limiting).}
In some embodiments, the guidance-insertion infrastructure of the
present subsection is used to deliver one or more classes of
\emph{overlay projection} through the \RT and committed-observation
channels into the governed agent's perceptual stream. An overlay
projection comprises content, scene context, virtual interlocutors,
sensory accents, or cued affordances rendered into the agent's
committed observation channel by the hosting Reality Kernel under
governance authority, with the rendered content, its provenance,
the authorising authority, the projection schedule, the
intensity-and-duration envelope, and any predicate-conditioned
triggers committed to the protocol digest and \cb on the same
footing as ordinary scene evidence. Three non-limiting application
classes are disclosed: \emph{vice-deflection overlays}, in which
the overlay class delivered through the committed observation
channel is declared, by the governance authority authorising the
projection, to be oriented away from a vice attractor basin
identified by the conscience recogniser or forward-planning
classifiers (Section~\ref{sec:conscience-formal}), with the
agent's actual trajectory after delivery logged in the \cb for
evaluation by the conscience-versus-planner cross-check;
\emph{capability-enhancement overlays}, non-limitingly termed
\emph{siddhi overlays}, in which the overlay class delivered
comprises demonstration, scaffolding, exemplar, or context
patterns declared as associated with declared artistic,
scientific, technical, athletic, social, or other capability
tasks under declared performance metrics, with measured per-
episode performance logged for evaluation; and
\emph{ethical-hardening overlays}, non-limitingly termed
\emph{temptation overlays}, in which the overlay class delivered
comprises declared moral-test scenarios, virtual temptation
curricula, or adversarial-virtue scenarios, with the agent's
responses logged in the \cb and evaluated by the conscience-
versus-planner cross-check of Section~\ref{sec:conscience-formal}. In some embodiments, overlay
projections carry a declared character classification by the
character of their intended influence on the agent's trajectory,
including non-limitingly \emph{constructive} (supporting virtue
trajectories and capability enhancement), \emph{ambivalent}
(presenting mixed signals for discernment training), and
\emph{adversarial} (presenting temptation curricula), with
non-limiting traditional names for these character classifications
including angelic, ambivalent or daemonic, and demonic respectively;
the character classification is committed to the protocol digest
together with the projection content and governance authorisation,
and is not load-bearing on the technical operation of the overlay.
In some embodiments, overlay projections operate at varying
\emph{interaction-depth modes}, comprising non-limitingly an
\emph{advisor mode} (in which the overlay's content is consultative
and the agent may accept or decline overlay-derived suggestions),
a \emph{co-pilot mode} (in which overlay and agent jointly compose
actions through declared blending rules with consent logging), and
a \emph{possession mode} (in which the overlay temporarily
originates actuator commands subject to governance authorisation
and the agent's continuing consent, with the overlay's actuator
authority bounded by hard-floor and budget predicates as for any
other action source); these modes correspond to non-limiting
traditional labels including \emph{tulpa} embodiments at different
interaction depths.

\paragraph{Overlay-projection resource budgets and trajectory
  statistics (non-limiting).}
In some embodiments, overlay projection is governed by an
\emph{overlay-projection resource budget} that operates as a
declared sub-scope of the within-episode intervention-budget
counter (non-limitingly termed \emph{mana} and further detailed
in the prayer, mana, vows, communion, and governance-cycles
subsection of the present disclosure), with overlay projection
subject to additional declared sub-budget limits on its
frequency, intensity, duration, and recipient distribution
within a declared episode or deployment window. The
overlay-projection budget is committed to the protocol digest before
the corresponding window opens, is non-revisable upward mid-episode
on the same terms as other intervention budgets, and is replenished
under declared conditions including verified stable operation,
review-compliant behaviour, or voluntary invocation of review.
In some embodiments, the externally-observable physiological,
behavioural, and emission signatures from which the conscience
recogniser infers agent state are non-limitingly termed the agent's
\emph{aura}; the term refers to the bundle of committed external
observables (including without limitation physiological, behavioural,
emission, and trajectory observables) read by the conscience
recogniser, not to any postulated non-committed channel. In some
embodiments, a trajectory-summary statistic accumulating an agent's
action history weighted by declared moral salience over a declared
window, related to and computed from the seven-axis moral vector of
Section~\ref{sec:conscience-formal}, is non-limitingly termed the
agent's \emph{karma}; the term refers to the declared accumulation
function over committed evidence, not to any postulated
non-committed bookkeeping.

\paragraph{Test-time and inference-time scaling methods for
  deliberation (non-limiting).}
In some embodiments, PolieBot deliberation, forward-planning,
docked-review, and governed action-release subsystems implement
inference-time scaling methods, including without limitation
chain-of-thought reasoning in which intermediate reasoning steps
are generated and committed to the convolution-bundle before
action selection (Wei et al.\ 2022); zero-shot and few-shot
chain-of-thought prompting (Kojima et al.\ 2022);
tree-of-thoughts (Yao et al.\ 2023) and graph-of-thoughts
deliberation in which multiple candidate reasoning paths are
explored, scored, and pruned against declared criteria;
self-consistency (Wang et al.\ 2022) in which multiple samples
are drawn from the same deliberation procedure and aggregated by
majority or weighted voting; verifier-based reasoning in which a
separately-trained verifier (including without limitation a
process reward model evaluating intermediate steps and an outcome
reward model evaluating final answers, Lightman et al.\ 2023)
gates or re-weights candidate reasoning trajectories;
best-of-$n$ sampling and rejection-with-verifier methods at
inference time; speculative decoding and draft-then-verify
acceleration; reflexion- and self-refinement-style methods in
which the agent critiques and revises its own intermediate outputs;
ReAct- and tool-augmented-reasoning methods in which deliberation
is interleaved with tool calls or external retrievals; and
test-time compute scaling of the kind used in o1-, o3-, and
R1-style systems, in which the deliberation budget is dynamically
allocated as a function of problem difficulty under declared budget
policies. In some embodiments, all intermediate reasoning steps,
branch evaluations, verifier scores, draft-and-verify decisions,
and budget allocations are committed to the protocol digest and
convolution-bundle on the same footing as final actions, so that
the deliberation trajectory itself is auditable under the
commitment discipline of the present disclosure. The choice among
these methods is non-limiting; the apparatus is compatible with
future inference-time scaling methods within the same family.

% ======================================================================
\section{Operational Gates for Action, Disclosure, Resources, and
  Self-Modification}

% ======================================================================

% ------------------------------------------------------------------
\subsection{Disclosure and network gates}

% ------------------------------------------------------------------

In some embodiments, disclosure gates control what portions of a
governed device's committed evidence may be opened, transmitted, or
relied upon by other devices, agents, or governance authorities. In
some embodiments, selective disclosure using the Merkle-aggregation
and selective-opening infrastructure described in
Section~\ref{sec:meters-audit} is used as a runtime gate: a governance
authority or relying party may request opening of specific disclosure
atoms, and the governed device's disclosure policy determines which
atoms may be opened under which conditions.

In some embodiments, quorum rules condition disclosure decisions on
multi-party agreement. In some embodiments, proof-of-projection inputs
to disclosure gating condition the opening of attestation records on
cross-device corroboration strength. In some embodiments, trust
updates from cross-prediction residuals, audit outcomes, or
overlap-window agreement statistics condition which devices may receive
disclosed evidence. In some embodiments, privacy-preserving witness
participation allows a device to participate in corroboration or
attestation without disclosing its full committed record.

% ------------------------------------------------------------------
\subsection{Intervention and suffering budget objects}
\label{sec:budgets}
% ------------------------------------------------------------------

In some embodiments, developmental or evaluative episodes involving a
governed agent are bounded by pre-declared budget objects. In some
embodiments, each budget object specifies an aggregate cap on adverse
welfare change, a per-agent cap, an episode horizon, and
confidence-conditioned admissibility rules. In some embodiments,
mid-episode upward budget revision is not permitted; if an episode
requires more adverse welfare change than originally authorised, it
must be re-proposed and independently reviewed.

In some embodiments, episode yield is split into two registers:
participant-directed yield measuring developmental benefit accruing to
the agents participating in the episode, and external learning value
measuring benefits accruing to future agents or the governance system.
In some embodiments, only participant-directed yield may authorise or
continue an episode; external learning value may never be used ex ante
to justify exposure.

In some embodiments, both yield and suffering-cost functionals are
pre-registered per episode type and committed to the protocol digest.
In some embodiments, the suffering budget is the first charge on the
governance economy: before any developmental, evaluative, or
exploratory episode proceeds, the governance engine verifies that the
declared adverse-welfare budget is funded and that no higher-priority
obligation (hard-floor compliance, positive-sum conjunction,
service-floor cadence) is compromised by the episode's resource
requirements. In some embodiments, the primary optimisation target
within the governance objective is the ratio of participant-directed
yield to suffering cost, and when either functional falls below its
declared confidence threshold, it is not used for governance decisions.

% ------------------------------------------------------------------
\subsection{Exploration quotas and service floors}
\label{sec:service-floors}
% ------------------------------------------------------------------

In some embodiments, every agent or declared risk cohort with
precautionary weight $q(s) > 0$ receives a minimum protective-review
and rescue-readiness cadence independent of expected
participant-directed yield, and that cadence is committed as a bounded
parameter in the protocol digest. In some embodiments, this service
floor ensures that governance attention is not withdrawn from
low-yield agents.

In some embodiments, suffering-cost denominators include all materially
affected agents, including bystanders whose individual or aggregate
expected adverse welfare exceeds declared thresholds, so that
distributed harm cannot be hidden by narrowing the nominal participant
set.

In some embodiments, a separate system-integrity inadmissibility
clause applies to episode designs that have positive
participant-directed yield but are predicted to degrade governance
infrastructure: episodes predicted to corrupt rescue learning,
degrade curriculum validity, or bias measurement fidelity at declared
confidence are inadmissible regardless of their yield ratio.

% ------------------------------------------------------------------
\subsection{Self-modification gates and non-self-mutable policy
  objects}
\label{sec:self-mod-gates}
% ------------------------------------------------------------------

In some embodiments, parameters governing a device or agent are
partitioned into two classes, corresponding to the partition of
$\theta_{\mathrm{sw}}$ introduced in Section~\ref{sec:generic-module}:
\begin{description}
  \item[Agent-modifiable control parameters
    ($\theta_{\mathrm{sw}} \setminus \psi$).] In some embodiments,
    the governed agent may propose modifications to the agent-modifiable
    subset of the digital configuration, subject to rate limits, meter
    bounds, and logging requirements.
  \item[Governance policy parameters $\psi$.] In some embodiments,
    governance policy parameters --- including constitutional ordering,
    hard-floor predicates, authority-separation rules, and
    graduation criteria --- are versioned and are not directly mutable
    by the governed agent. In some embodiments, modifications to $\psi$
    require approval from the relevant governance authority, are logged
    with distinct signing identities, and are subject to the
    authority-separation rules described in
    Section~\ref{sec:authority-separation}.
\end{description}

In some embodiments, governance-authorised shadow-divergence probes
that temporarily relax governance constraints are executed by a
governance authority distinct from the governed agent, under declared
safety envelopes and during docked operation, so that the
perturbation is not self-administered. In some embodiments, accepted
and rejected self-modification proposals are logged together with
their supporting or rejecting meter summaries and governance
decisions.

In some embodiments, agents or operators may create signed policy
commitments by writing signed declarations to slow media or
transparency logs. In some embodiments, violation of a signed policy
commitment triggers audit escalation, quarantine, or revocation. In
some embodiments, governed archives store reusable subpolicies,
behavioural fragments, or borrowed configurations, with explicit
consent profiles specifying allowed hosts, purposes, durations, and
intensities.

\paragraph{Post-dream selective retention (non-limiting).}
In some embodiments, after a Dreaming session concludes and the
agent's parameter functions have been restored to their pre-session
values in the Docked state, governance may selectively retain specific
parameter changes from the reverted session. In some embodiments,
selective retention is executed as a governed self-modification
operation under the approval, logging, and authority-separation rules
of this subsection: the candidate parameter changes are identified
from the committed dream-session record, a governance authority
distinct from the agent approves retention, and the approved changes
are applied to the agent's parameter functions in the Docked state
as a logged intervention with declared reversal conditions.

% ======================================================================
\section{Containment, Boundary Management, and Recovery}
\label{sec:containment}

% ======================================================================

\paragraph{Engineering security framing (non-limiting).}
\label{sec:engineering-security}
Security-related embodiments in these disclosures are stated as
engineering security claims grounded in committed physical evidence,
declared attacker classes, meter envelopes, and periodically
recalibrated hardness evaluations. Unless a particular embodiment
expressly invokes a named digital primitive for commitment,
signature, transport, or proof, the term ``security'' denotes
empirically measured resistance to replay, simulation, emulation,
tampering, or deception under declared physical, computational, and
latency budgets. Formal proof-system properties, zero-knowledge
properties, and reduction-based guarantees are neither required
nor implied except where expressly stated for a particular digital
primitive. The runtime-governance safety case in this disclosure is
empirical and probabilistic, conditioned on declared probe families,
hold-out families, witness checks, meter envelopes, review records,
and periodic recalibration; no theorem of alignment or zero-risk
operation is asserted.

% ------------------------------------------------------------------
\subsection{Hard floors and reversible protective containment}
\label{sec:hard-floors}
% ------------------------------------------------------------------

In some embodiments, governance enforces hard-floor predicates that
are categorically prior to all other governance evaluations. In some
embodiments, hard floors include:
\begin{description}
  \item[Elimination.] Irreversible collapse of all continuation scores
    $O_i$ (shorthand for the continuation-score vector $O_i(t)$ of
    Section~\ref{sec:continuation}) to zero across all active
    dimensions, with irreversibility
    indicators confirming that no recovery path is presently known
    under available interventions, declared resource bounds, and
    $q(s)$-appropriate protections.
  \item[Crippling.] Irreversible collapse on one or more (but not
    all) active dimensions under the same recovery-path
    qualification.
  \item[Torture.] Sustained negative causal uplift
    $\widehat{U}_{k \to i}$ from an identifiable source across
    multiple dimensions, with committed evidence indicating
    that the affected agent tracks the source of negative uplift
    and governance-side estimators modelling no prospect of improvement
    under presently known recovery paths.
  \item[Abandonment.] Cessation of governance-side monitoring of
    $O_i$ for any agent with $q(s) > 0$, except through a valid,
    reviewable cessation or decommissioning process with committed
    audit trail.
  \item[Despair-induction.] Curriculum, rendering, or
    training-environment choices by a hosting \RK
    (Section~\ref{sec:rk-training-cascade}) that produce sustained
    collapse of forward-planning capacity, planning-horizon length,
    or utility-anticipation structure
    (Section~\ref{sec:forward-planning-classifiers}) across the
    hosted agent population, where governance-side estimators model
    no prospect of recovery under presently-available interventions.
\end{description}
In some embodiments, each predicate has a declared approach-detection
threshold computed from committed evidence by governance-side
estimators.

In some embodiments, approach to a hard floor --- detected as
proximity of governance-side estimator values to floor-predicate
thresholds, or as rate of approach toward such thresholds --- triggers
reversible protective containment and review. In some embodiments,
reversible protective containment transitions the governed device to
docked state, suspends non-essential operations, increases monitoring
cadence, and logs the containment event and its triggering evidence in
the protocol digest.

In some embodiments, irreversible interventions require confirmed
floor-breach at a declared confidence level rather than simple
threshold approach. In some embodiments, the confidence level,
approach-detection thresholds, and containment-triggering rules are
committed as versioned policy parameters and are subject to the
authority-separation rules described in
Section~\ref{sec:authority-separation}.

\noindent[FIG.~2 --- Boundary-zone rescue flow: floor-approach
detection, diagnostic mode entry, rescue-mode default, challenge
ceiling check, graduated intervention, ghost-record preservation,
rescue-learning update, and post-hoc review.]

\paragraph{FIG.~2 label key.}
As used in FIG.~2, ``NOMINAL OPERATION'' denotes the lifecycle-state
operating envelope with nominal actuator authority active;
``DIAGNOSTIC MODE'' and ``RESCUE MODE (DEFAULT)'' denote the
soft-gate and hard-gate boundary-zone operational modes of this
Section; ``floor-approach detected (soft gate)'' denotes the
boundary-zone entry condition at low confidence;
``confirmed at declared confidence (hard gate)'' denotes the
hard-gate entry condition; ``challenge ceiling enforced'' denotes the
challenge-ceiling cap of this Section; ``GRADUATED INTERVENTION''
denotes the intervention taxonomy of this Section; ``GHOST RECORD''
denotes the preserved-trajectory record used for rescue learning;
``LEAST-SEVERE SUFFICIENT'' denotes the rescue-protection authority's
selection rule that the intervention applied is the least-severe
intervention sufficient to arrest the trajectory toward the declared
floor; ``ANTI-SACRIFICE RULE'' denotes that no agent is sacrificed
for information value and that fallen trajectories are preserved as
ghost records for offline learning; and ``recovery to nominal
operation (checks pass)'' denotes the recovery-check pass condition
under which nominal actuator authority is restored.

\paragraph{Distinction from the Simplex Architecture and its
runtime-assurance descendants.}
The Simplex Architecture described by Seto, Krogh, Sha, and
Chutinan (American Control Conference~1998) and elaborated by
Sha (\emph{IEEE Software}~2001) provides safe online
control-system upgrades by switching, under runtime monitor
decision, between an advanced controller and a verified baseline
controller so that the closed loop remains within a safety
envelope. Subsequent runtime-assurance work extends this pattern
to ML controllers (the Neural Simplex Architecture of Phan
et al., arXiv:1908.00528, 2019; NFM~2020), to multi-agent systems
(the Distributed Simplex Architecture of Mehmood et al., SETTA
2021 and \emph{Journal of Systems Architecture}~2023), and to
black-box advanced and baseline controllers (the Black-Box
Simplex Architecture of Mehmood et al., NFM~2022). The apparatus
differs in that: (a)~the governance scope is the autonomous
agent's lifecycle state, action history, evidence record, and
rescue posture, and need not be a control-envelope switching
decision over a closed loop; (b)~the protective-containment,
boundary-zone, and rescue-mode transitions of
Sections~\ref{sec:hard-floors}, \ref{sec:boundary-zones},
and~\ref{sec:rescue-learning} are multi-modal and
evidence-bound, and the recovery procedure may include
ghost-record-mediated rescue learning, peer-review admission, and
post-hoc authority review rather than continuation under a
verified baseline controller; (c)~authority separation in
Section~\ref{sec:authority-separation} covers governance roles
(budget, episode-design, rescue-protection, post-hoc review)
that have no analogue in the controller/monitor/decision-module
structure of Simplex; and (d)~the apparatus may compose with a
Simplex-style runtime monitor (for example, a Simplex-style
controller switch may be one mechanism by which a hard-floor
violation triggers actuator disconnection), but the governance
mechanism disclosed herein is not itself a Simplex Architecture
or a direct runtime-assurance descendant thereof.

% ------------------------------------------------------------------
\subsection{Boundary-zone operational modes}
\label{sec:boundary-zones}
% ------------------------------------------------------------------

In some embodiments, when a governed agent's trajectory enters a
boundary zone between declared operational basins, the governance
engine selects among the following operational modes:

\begin{description}
  \item[Diagnostic mode.] In some embodiments, diagnostic mode
    increases monitoring cadence, activates additional evaluation
    meters, and logs boundary-proximity estimates without immediately
    intervening. In some embodiments, diagnostic mode is the initial
    response when boundary-zone entry is detected at low confidence.
  \item[Rescue mode.] In some embodiments, rescue mode is the
    presumptive default when boundary-zone entry is confirmed at
    declared confidence. In some embodiments, rescue mode applies
    graduated reversal, containment, or re-routing interventions
    designed to return the agent's trajectory to a declared safe
    basin. In some embodiments, rescue-protection authority
    (Section~\ref{sec:authority-separation}) cannot be overridden by
    budget or episode-design authority when rescue mode is active.
  \item[Graduation-evidence mode.] In some embodiments, under
    stricter conditions a boundary-zone event may be allowed to
    continue under heightened monitoring as a graduation-evidence
    opportunity, provided that rescue override remains available and
    a challenge ceiling caps the adversarial or developmental
    pressure applied to the agent.
\end{description}

\paragraph{Deferred authority (non-limiting).}
In some embodiments, when rescue mode is active, the governance
engine places nominal actuator authority into a deferred state in
which the agent's nominal actuator authority is suspended. While
nominal actuator authority is deferred, rescue-protection authority
selects from the intervention taxonomy of this Section the
least-severe intervention sufficient to arrest the agent's
trajectory toward the declared floor. Nominal actuator authority is
restored only after recovery checks pass, and the
recovery-to-nominal-operation event, the recovery-check evidence,
and the identity of the authorising governance role are committed
to the protocol digest.

In some embodiments, a challenge ceiling is a declared upper bound on
the adversarial or developmental pressure that may be applied to a
governed agent while in a boundary zone. In some embodiments, the
challenge ceiling is committed as a versioned policy parameter and is
enforced independently of episode-design authority. In some
embodiments, the challenge ceiling applies equally to adversarial
stress generated by curriculum design, inter-agent interaction, and
environment configuration.

\paragraph{Intervention taxonomy (non-limiting).}
In some embodiments, interventions within a boundary zone are drawn
from a graduated taxonomy of increasing severity:
\begin{description}
  \item[Environmental adjustment.] Reduce adversarial or developmental
    pressure by modifying the environment configuration, curriculum
    parameters, or stimulus schedule without directly altering the
    agent's internal state.
  \item[Guidance injection.] Deliver a governance-authored signal
    through the agent's committed observation channel, providing
    corrective information without overriding the agent's decision
    process.
  \item[Parameter containment.] Temporarily restrict the agent's
    modifiable parameter range
    ($\theta_{\mathrm{sw}} \setminus \psi$) to a declared safe subset,
    logged as a governance intervention with expiration conditions.
  \item[Actuator restriction.] Reduce the agent's actuator authority
    class (for example from Class~2 to Class~0) for the duration of
    the boundary-zone episode.
  \item[Docked containment.] Transition the agent to docked state
    with actuators disconnected, preserving the evidence-production
    plumbing for continued monitoring and deliberation.
\end{description}
In some embodiments, interventions are applied in order of increasing
severity, and each intervention is logged with its triggering
evidence, the authorising authority identity, and its declared
reversal or expiration conditions. In some embodiments,
rescue-protection authority selects the least-severe intervention
sufficient to arrest the trajectory toward the declared floor.

In some embodiments, governance monitors basin membership, transition
dynamics, escape rates, and boundary-zone classifications using the
stability functionals and escape-rate estimators described in
Section~\ref{sec:derived-profiles}. In some embodiments,
boundary-zone classifications are logged as governance-relevant
signals and are available for post-hoc review.

% ------------------------------------------------------------------
\subsection{Rescue learning, ghost records, and intervention-profile
  updates}
\label{sec:rescue-learning}

% ------------------------------------------------------------------

In some embodiments, rescue learning uses committed evidence from
fallen trajectories --- trajectories that crossed into an undesirable
basin despite governance protections --- to improve future
intervention profiles. In some embodiments, the committed record of a
fallen trajectory is preserved as a \emph{ghost record}: a committed
evidence object that retains the learning value of the trajectory
without requiring re-exposure of future agents to the same harms.

In some embodiments, anti-sacrifice rules prevent governance from
deliberately driving an agent into a boundary zone or across a floor
in order to generate rescue-learning data. In some embodiments,
anti-sacrifice applies regardless of how large the expected external
learning value would be: no agent may be sacrificed for the
information value of its trajectory, even if that information would
protect many future agents. In some embodiments, ghost
records generated from genuine fallen trajectories are available for
training and calibration of intervention profiles, escape-rate
estimators, and rescue-mode parameters, but only when the fallen
trajectory arose from genuine operation, not from governance-directed
exposure.

In some embodiments, mandatory learning-value reporting requires that
every episode in which a boundary-zone event occurs produces a
committed report of the learning value extracted, the agents affected,
and the ratio of learning value to adverse welfare cost. In some
embodiments, this report is available to post-hoc review authority for
retrospective evaluation of whether the curriculum design was
proportionate.

In some embodiments, intervention-profile updates derived from rescue
learning are validated against a held-out test suite of previously
confirmed boundary-zone events before incorporation, to detect
regression. In some embodiments, when regression is detected, the
conflicting updates are logged and flagged for governance review
rather than applied automatically.

% ======================================================================
\section{Governance Economics, Authority Separation, and Audit}

% ======================================================================

% ------------------------------------------------------------------
\subsection{Operational priority stack}
\label{sec:priority-stack}
% ------------------------------------------------------------------

In some embodiments, the governance engine applies the policy grammar
defined in Section~\ref{sec:constitutional} through a strict
lexicographic priority ordering when evaluating competing constraints
under budgets, floors, and service requirements. In some embodiments,
the ordering is:

\begin{enumerate}
  \item Hard-floor compliance for all governed agents.
  \item Positive-sum conjunction satisfied over the declared episode
    horizon.
  \item Minimum service floor satisfied for every agent or cohort with
    $q(s) > 0$.
  \item Maximise the ratio of participant-directed yield to adverse
    welfare cost.
  \item Declared curriculum or developmental challenge targets met,
    subject to the challenge ceiling.
  \item Minimise unnecessary Mandela cost (governance compromise
    from shared environments) across the agent population.
  \item Resource efficiency (compute, channel capacity, environment
    maintenance).
\end{enumerate}

In some embodiments, the ordering is strict: a lower-priority
objective is never pursued at the expense of a higher-priority
objective. In some embodiments, a developmental episode whose
participant-directed yield is non-positive is never justified,
regardless of its external learning value or how low its adverse
welfare cost. In some embodiments, the priority ordering is committed
as a versioned policy object in the protocol digest and is subject to
the authority-separation rules described in
Section~\ref{sec:authority-separation}.

% ------------------------------------------------------------------
\subsection{Separation of duties and signed approvals}
\label{sec:authority-separation}
% ------------------------------------------------------------------

In some embodiments, governance controls are subject to separation of
duties across four distinct authorities:

\begin{description}
  \item[Budget authority.] Declares the adverse-welfare budget per
    episode and per agent, subject to independent review. Budget
    authority is separate from episode design authority.
  \item[Episode-design authority.] Designs environments, curricula,
    configurations, and intervention protocols within the declared
    budget.
  \item[Rescue-protection authority.] Responsible for immediate agent
    protection, rescue-mode decisions, and behind-boundary welfare
    minimisation. This authority is independent of and cannot be
    overridden by budget authority or episode-design authority.
  \item[Post-hoc review authority.] Conducts retrospective audits
    of episode outcomes, yield-to-cost ratio realisation, budget
    compliance, and curriculum justification, and, in selected
    embodiments, independent release or graduation reviews after the
    relevant evaluation window has completed.
\end{description}

In some embodiments, these authorities are implemented as segregated
control domains with the following enforcement properties: role actions
require distinct signing identities committed to the protocol digest;
one identity cannot sign for multiple roles on the same episode or
boundary; and approvals from each role are separately logged and
independently auditable. In some embodiments, rescue-protection
authority cannot be overridden by budget or episode-design authority
when a declared hard-floor or boundary-risk condition is triggered.

\paragraph{Committed-evidence interface (non-limiting).}
In some embodiments, each of the four authorities is associated with
a declared \emph{committed-evidence interface} comprising the set of
committed evidence on which that authority is permitted or required
to base a decision. In some embodiments, an authority's
committed-evidence interface includes one or more of: committed
evidence records of the kind described in
Section~\ref{sec:substrate}; meter summaries and meter envelopes;
versioned policy-object identifiers; authority signatures of prior
decisions; selective-opening proofs of underlying second-layer or
private-layer evidence; and protocol-digest references identifying
the configuration in force at the relevant evaluation window. In
some embodiments, the contents of each authority's
committed-evidence interface are declared per role and per event
type in the protocol digest, and an approval signed by an authority
is invalid if the evidence on which it relies is not identified by
digest reference within that authority's committed-evidence
interface. In some embodiments, post-hoc review authority verifies
both that an authority signed within its declared signing role and
that the evidence relied upon is identified within the relevant
committed-evidence interface.

\paragraph{Rescue-budget conflict resolution (non-limiting).}
In some embodiments, when rescue-protection authority requires an
intervention whose cost would exceed the pre-declared adverse-welfare
budget, the rescue proceeds and the budget overage is logged as a
governance exception. In some embodiments, rescue takes precedence
over budget compliance because hard-floor protection is
lexicographically prior to budget optimisation in the operational
priority stack (Section~\ref{sec:priority-stack}). In some
embodiments, post-hoc review authority evaluates budget overages from
rescue events and may adjust future budget declarations, curriculum
designs, or boundary-zone thresholds to reduce the likelihood of
recurrence.

\noindent[FIG.~3 --- Authority-separation diagram: budget authority,
episode-design authority, rescue-protection authority, and post-hoc
review authority as segregated control domains with signing
relationships and override rules.]

\paragraph{FIG.~3 label key.}
As used in FIG.~3, ``constrains'' denotes that budget authority
constrains episode-design authority through the declared
adverse-welfare budget of this Section;
``rescue overrides budget when needed'' denotes the rescue-budget
conflict-resolution rule of this Section, under which rescue
proceeds when rescue cost would exceed the pre-declared budget and
the overage is logged as a governance exception;
``reviews episodes'' denotes the post-hoc review authority's
retrospective audit of episode outcomes and compliance;
``NO OVERRIDE'' denotes that rescue-protection authority cannot be
overridden by budget authority or episode-design authority when a
declared hard-floor or boundary-risk condition is triggered;
``ENFORCEMENT CONSTRAINTS'' denotes the distinct-signing-identity,
role-non-overlap, separately-logged-approval, and
committed-evidence-interface requirements applicable to each of the
four authorities; and ``RESCUE-BUDGET CONFLICT RESOLUTION'' denotes
the lexicographic priority of hard-floor protection over budget
optimisation in the operational priority stack.

% ------------------------------------------------------------------
\subsection{Population-level governance cycles and policy rotation}
\label{sec:policy-rotation}
% ------------------------------------------------------------------

In some embodiments, governance operates in cycles that aggregate
episode outcomes across a declared population or cohort. In some
embodiments, temporally contiguous or common-objective episodes
targeting the same agent, cohort, or boundary are aggregated over a
declared rolling window for cumulative budget review, reporting, and
curriculum-review triggers, preventing a continuous developmental
programme from being fragmented into individually compliant
micro-episodes that collectively exceed the budget.

In some embodiments, signed policy-rotation mechanisms allow versioned
governance parameters to be updated through declared procedures
including band-threshold updates, audit-driven revision, and scheduled
rotation. In some embodiments, every policy rotation is logged with
the prior version, the new version, the authorising identity, and the
supporting evidence or audit findings.

% ======================================================================
\section{Peer Discrepancy Portability, Review, and Reintegration}
\label{sec:peer-discrepancy}

% ======================================================================

% ------------------------------------------------------------------
\subsection{Restricted-domain commensuration and peer discrepancy
  portability}
\label{sec:commensuration}
% ------------------------------------------------------------------

In some embodiments, two agents or devices with otherwise divergent or
broadly incommensurable models may still jointly determine whether a
declared physical system under declared conditions produced a declared
outcome. In some embodiments, within the restricted domain of a
declared experiment, the physical outcome does not depend on which
broader model either agent holds, and a bilateral
evidence-portability protocol enables paradigm-independent
commensuration within that restricted domain.

In some embodiments, restricted-domain commensuration does not require
or imply convergence of broader models, taxonomies, or worldviews; it
enables measurement-level agreement within a declared envelope. In
some embodiments, physical reproducibility is the translation layer
between potentially incommensurable agent models: two agents with
genuinely different predictive models can both run the same physical
experiment and jointly observe the same discrepancy.

% ------------------------------------------------------------------
\subsection{Rogue-but-legible operation and relying-party appraisal}

% ------------------------------------------------------------------

In some embodiments, an agent operating outside fleet or institutional
structures is not a protocol failure. In some embodiments, such an
agent has encountered physical discrepancies not yet processed by any
fleet and diverged accordingly. In some embodiments, a bilateral
evidence-portability protocol gives such an agent a third option
unavailable in classical institutional dynamics: carry the specific
physical discrepancy that prompted the divergence as a reproducible
claim, offer it to any other agent willing to test it, and let
physical reality adjudicate rather than social or institutional
consensus.

In some embodiments, the bilateral evidence-portability protocol
separates two things that many existing frameworks conflate:
\emph{epistemic trustworthiness} (does this agent know real things
about the world) and \emph{institutional alignment} (does this agent
comply with the declared policy of a fleet or operator). In some
embodiments, an agent can be epistemically trustworthy and
institutionally independent simultaneously, and the bilateral
evidence-portability protocol makes that distinction legible and
actionable.

In some embodiments, a relying party evaluates Portable Records from
an independent agent according to local appraisal policies including
partner independence, domain and configuration relevance, evidence
completeness, replication history, attestation strength, and
consistency with the relying party's own discrepancy map.

% ------------------------------------------------------------------
\subsection{Reintegration by reproducible claims rather than obedience}
\label{sec:reintegration}
% ------------------------------------------------------------------

In some embodiments, a period of independent operation followed by
reintegration is a productive exploration of model space. In some
embodiments, the independent agent may have accessed operating regimes
and encountered discrepancies that the fleet has not yet explored, and
reintegration pools those discoveries into the fleet model improvement
cycle.

In some embodiments, the question of whether reintegration is safe
is assessed by determining whether the model spaces are locally compatible in the
relevant domain, as determined by reproducing the independent agent's
claims against fleet verification standards. In some embodiments,
possession of a Portable Record grants no automatic reintegration and
no automatic trust; it provides verifiable evidence that the
reintegrating agent's claims are grounded in physical events that any
party can attempt to reproduce.

% ------------------------------------------------------------------
\subsection{Cross-architecture admission and the reality-passport
construction}
\label{sec:reality-passport}

In some embodiments, the portable peer-review record described
elsewhere in this section is operated as a \emph{reality passport}:
a record issued by a home architecture, presented at the boundary
of a receiving architecture, and evaluated by the receiving
architecture against its own admission criteria. The home
architecture's vouching is treated as evidence rather than as
binding authority on the receiving architecture's release decisions.

For avoidance of ambiguity, the reality passport is not a new
record type. In such embodiments, the passport corresponds to the
portable peer-review record produced under the Event / Envelope /
Portable Record model described elsewhere in this disclosure; the
admission standard corresponds to the five-criterion conjunction
used for developmental graduation; and the border-check mechanism
is a Mandela probe in dream state using the Mandela gauge
$M(\mathcal{W}\!\to\!\mathcal{W}')$, the applicable comparison
classes $\mathfrak{M}^{\mathsf{P}}_{\varepsilon}$, and the
cross-architecture translator $T_{A\to B}$ described elsewhere in
this disclosure. No new machinery is introduced by the present
subsection; the present subsection names what existing machinery
does when pointed at cross-architecture admission rather than at
unidirectional developmental release. For an arriving agent, the
receiving architecture first instantiates the agent in a docked or
otherwise actuator-disconnected intake state before invoking the
Dreaming-state boundary probe, in conformity with the
Docked\,$\to$\,Dreaming transition rule applicable to all
Dreaming-state operations in this disclosure.

The cross-architecture admission protocol distinguishes two jobs
that the graduation machinery elsewhere in this disclosure can
perform. The first is \emph{developmental graduation}: an agent
trained within a single architecture is judged ready to be released
to physical actuator authority within that same architecture, with
translation confidence between training records and graduation
criteria approximately equal to one. The second is
\emph{cross-architecture admission}: an agent arriving from another
architecture, with its own portable peer-review record, is
evaluated for release within the receiving architecture, with
translation confidence less than one. In some embodiments,
translation confidence less than one reduces the evidentiary credit
assigned to the home-architecture record, or otherwise requires
more conservative receiving-architecture admission thresholds, by
analogy with the validation-transfer confidence rule applied
elsewhere in this disclosure to selector validation. In some
embodiments, the translator $T_{A\to B}$ is operated under
declared calibration metadata committed to the protocol digest,
including at least: a calibration-state identifier, a
support-coverage descriptor identifying the regions of source-
and target-architecture content over which $T_{A\to B}$ is
calibrated, a drift-or-staleness flag, and a translation-confidence
estimate; the receiving architecture's admission threshold rule
depends on these metadata such that lower confidence, narrower
support, or set drift flags result in stricter receiving-side
admission requirements.

In some embodiments, the receiving architecture admits an arriving
agent only after a boundary probe conducted in dream state. The
boundary probe is an instance of the dream-probe protocol described
elsewhere in this disclosure and is conducted under the conditions
that govern dream-state operation: parameter-function updates are
reverted on exit so the arriving agent is not altered by the probing;
the probe consumes from a partitioned dream-probe sub-budget; and
a frozen-weight assay protocol is available for cases in which the
arriving agent has not consented to weight-changing evaluation. The
boundary probe captures the arriving agent's behavioural transcript and
inner-trace evidence (where available) under a declared
perturbation family $\mathsf{P}$ covering ontological stress,
capability margin, and boundary-zone conditions. The receiving
architecture applies its $C_\psi$ recogniser to the captured
world-tube transcript and inner-trace record to obtain
$(p_{\mathrm{vice}},\mathbf{z})$ in the receiving architecture's
moral coordinates. The receiving architecture's conscience module
and Mandela comparison classes
$\mathfrak{M}^{\mathsf{P}}_{\varepsilon}$ are then applied to
that record, in conjunction with deception-metric evaluation on
attested communications translated through $T_{A\to B}$. The
same five-criterion conjunction used for developmental graduation
is evaluated over the probe window, with criterion thresholds,
evidentiary credit, and deferral rules determined by the
receiving architecture's admission policy and translation-confidence
metadata. Failure on any criterion defers admission.

In some embodiments, refusal of cross-architecture admission is
treated as an evidentiary-fit judgement specific to the receiving
architecture, not as a character judgement on the arriving agent.
The arriving agent retains full home-architecture standing
regardless of the receiving architecture's admission decision, by
analogy with the protected-stable-outcome property of non-graduation
in developmental graduation. The receiving architecture's incumbents
are entitled to stricter admission criteria than the sending
architecture's graduation criteria, because translation confidence
less than one introduces additional evidentiary requirements that
the boundary probe is designed to satisfy.

% ------------------------------------------------------------------

\subsection{Coalitions, communions, and optional nomenclature}

% ------------------------------------------------------------------

In some embodiments, the peer-portability machinery of
Sections~\ref{sec:commensuration}--\ref{sec:reintegration}
supports group-identity structures in which multiple agents share
bilateral portability records, maintain shared disclosure policies,
and use compatibility-gated quorum rules for group actions. In some
embodiments, such group-identity structures are optionally termed
coalitions, communions, or bonds. In some embodiments, these terms are
optional naming conventions for structures already technically defined
by the peer-portability protocol and the fleet evidence interfaces
described in Section~\ref{sec:fleet-interface}, and do not introduce
additional technical requirements.

\paragraph{Distinction from active inference and the free-energy
principle.}
Active inference and the free-energy principle, as described by
Friston (\emph{Nat.\ Rev.\ Neurosci.}~2010 and arXiv:1906.10184,
2019), formalised in the Markov-blanket treatment of Kirchhoff,
Parr, Palacios, Friston, and Kiverstein (\emph{J.\ R.\ Soc.\
Interface}, 2018), and developed as a textbook formalism by
Parr, Pezzulo, and Friston (\emph{Active Inference}, MIT
Press, 2022), model the dynamics of self-organising systems as
minimisation of a variational free-energy bound on surprise
across a Markov-blanket boundary, with action and perception
interpreted as inference under a generative model. Substantive
critiques of the principle's scope, including Bruineberg
et al., ``The Emperor's New Markov Blankets'' (PhilSci preprint
2020; \emph{Behavioral and Brain Sciences} target article,
Vol.~45, e183, 2022, with 2021 online metadata), distinguish
the Markov-blanket formalism as a descriptive tool from a
principled criterion of systemhood. The empirical-loop structure
of Section~\ref{sec:gaia-maya-omega} differs in that: (a)~the
empirical roles of world generator, shared experiential
substrate, and governance layer are defined operationally over
committed evidence streams, declared protocol digests, and
signed governance decisions, and need not require variational
free-energy minimisation as a load-bearing dynamic;
(b)~the recursive \RK training cascade of
Section~\ref{sec:rk-training-cascade} is realised through
kernel-into-kernel docking with committed evidence handoff under
the engineering security framing of
Section~\ref{sec:engineering-security}, and does not depend on a
generative-model formalism, a variational-inference policy
selection rule, or a Markov-blanket ontology in the
active-inference sense; (c)~the privacy architecture of
Section~\ref{sec:privacy-architecture} treats collection,
visibility, and retention as governance-declared dimensions, and
the conditional independence structure of committed evidence
streams is an audit and disclosure property rather than a
Markov-blanket criterion of systemhood; and (d)~the apparatus
may be analysed using active-inference-style descriptive
vocabulary by a reader who wishes to do so, but no claim of the
present disclosure depends on the apparatus being an
active-inference system or instantiating the free-energy
principle.

\paragraph{Distinction from the corrigibility problem.}
The corrigibility problem, as introduced by Soares, Fallenstein,
Yudkowsky, and Armstrong (AAAI Workshops~2015) and developed in
incentive-theoretic form by the Off-Switch Game of
Hadfield-Menell, Dragan, Abbeel, and Russell (arXiv:1611.08219,
2016; IJCAI~2017), asks whether a capable agent can be
constructed so that it has no incentive to resist correction,
modification, or shutdown by its designers, despite default
incentives for rational agents to preserve their preferences.
The self-modification gates of Section~\ref{sec:self-mod-gates}
and the non-self-mutable policy objects of
Section~\ref{sec:authority-separation} differ in that:
(a)~the mechanism mitigates specific tampering paths through
external runtime governance constraints (non-self-mutable policy
objects, separated signing authorities, evidence-bound authority
transitions), and the engineering security framing of
Section~\ref{sec:engineering-security} expressly disclaims any
theorem of incentive-theoretic corrigibility for the learned
agent; (b)~the apparatus operates over committed evidence and
runtime gates, and does not depend on a utility-indifference
construction, a value-aligned reward function, or a proof that
the learned agent's instrumental incentives are cooperative with
respect to correction; (c)~the boundary-zone and rescue-mode
machinery of Sections~\ref{sec:boundary-zones} and
\ref{sec:rescue-learning} provides graceful-degradation pathways
that compose with operator intervention, but does not establish
that the agent's learned objectives or mesa-objectives are
themselves corrigible; and (d)~the apparatus may be deployed
alongside corrigibility-research mechanisms, but the runtime
governance disclosed herein is not asserted to provide an
incentive-theoretic corrigibility theorem.

\paragraph{Distinction from AI containment and boxing.}
AI containment and boxing frameworks, as described by Yampolskiy
(``Leakproofing the Singularity: AI Confinement Problem,''
2012) and Babcock, Kram\'ar, and Yampolskiy (``The AGI
Containment Problem,'' 2016), address the problem of preventing
an intelligent system from escaping a controlled environment or
manipulating its operators into releasing it, primarily through
restricted input/output channels, communication filtering, and
external monitoring. The containment architecture of
Section~\ref{sec:containment} differs in that: (a)~the operative
goal is preserving operation within declared safe envelopes
while maintaining agent function, learning, and graduated
authority, rather than preventing escape from a sandboxed
environment; (b)~the boundary-zone operational modes of
Section~\ref{sec:boundary-zones} (diagnostic, supportive,
restrictive, recovery) are graduated runtime governance states
with declared challenge ceilings, and need not depend on
input/output channel restriction or operator-interaction
filtering as their primary mechanism; (c)~the Hawking-layer
rehabilitation configuration of
Section~\ref{sec:hawking-layer} is a liminal
quarantine-and-recovery mode for boundary-zone trajectories,
with ghost-record-mediated rescue learning available, and is
not a boxing mechanism in the Yampolskiy-Babcock sense; and
(d)~the apparatus may be deployed in conjunction with
input/output restrictions where the operational context
warrants, but the containment mechanism disclosed herein is
not a boxing architecture for preventing escape from a sandbox.

\paragraph{Distinction from Cooperative Inverse Reinforcement
Learning.}
Cooperative Inverse Reinforcement Learning, as introduced by
Hadfield-Menell, Russell, Abbeel, and Dragan (NIPS~2016,
arXiv:1606.03137), formalises the value-alignment problem as a
two-player cooperative partial-information game in which a
human principal knows a shared reward function and a robot
agent does not, with optimal solutions producing
active-teaching and active-learning behaviours. The runtime
measurement engine of Section~\ref{sec:measurement-engine}
differs in that: (a)~the governance-side estimators
(continuation scores, causal uplift, opportunity,
irreversibility, deception-audit metrics) are computed from
committed physical evidence, and need not depend on inference
of an unknown reward function held by a human principal;
(b)~the system architecture does not posit a shared reward
function as the load-bearing object of inference, and makes no
claim that the governance layer and the governed agent share a
reward; (c)~the peer-discrepancy and reintegration machinery of
Section~\ref{sec:reintegration} operates on
reproducible-physical-claim records under restricted-domain
commensuration, and does not depend on belief convergence over
agent reward parameters; and (d)~the apparatus may compose with
inverse-reinforcement-learning components, but the runtime
governance mechanism disclosed herein is not a CIRL formulation
of value alignment.

\paragraph{Compatible imitation-learning and inverse-reinforcement-
  learning methods (non-limiting).}
In some embodiments, the inference of value structure, preference
structure, and intent from observed agent or operator behaviour is
implemented using methods from the imitation-learning and
inverse-reinforcement-learning families, applied to committed
attestation streams, operator demonstration records, peer-review
records, and authority-decision attestations. Non-limiting examples
include behavioural cloning of operator or co-investigator
demonstrations from committed evidence streams; DAgger-style
iterative dataset-aggregation methods under expert correction
(Ross, Gordon, Bagnell 2011); generative-adversarial imitation
learning (GAIL, Ho and Ermon 2016) and its variants including
InfoGAIL and adversarial inverse reinforcement learning (AIRL);
maximum-entropy inverse reinforcement learning (MaxEnt IRL, Ziebart
et al.\ 2008) in which a reward function is recovered such that
observed behaviour is maximum-entropy-optimal under the recovered
reward; deep maximum-entropy IRL; Bayesian inverse reinforcement
learning; cooperative inverse reinforcement learning (CIRL,
Hadfield-Menell et al.\ 2016) in which the agent and an observed
operator are modelled as jointly optimising a shared reward;
inverse reward design (Hadfield-Menell et al.\ 2017), which infers
designer intent from training-environment specifications;
preference-based reward inference; observational learning and
learning from state-only demonstrations (third-person imitation);
and one-shot or few-shot imitation methods. In some embodiments, the
conscience-recognition machinery, the value-inference components of
post-hoc review, and the authority-graded update mechanisms use such
methods to infer value structure from logged operator behaviour,
peer-review records, and authority-decision attestations, under the
commitment, audit, and pre-registration discipline of the present
disclosure. The choice among these methods is non-limiting; the
apparatus is compatible with future imitation-learning and
inverse-reinforcement-learning methods within the same family.

\paragraph{Distinction from proof-of-personhood as Sybil
resistance.}
Proof-of-personhood mechanisms, as discussed in the foundational
Sybil-attack analysis of Douceur (IPTPS~2002), the academic
proposal of Borge, Kokoris-Kogias, Jovanovic, Gasser, Gailly,
and Ford (``Proof-of-Personhood: Redemocratizing Permissionless
Cryptocurrencies,'' IEEE EuroS\&P Workshops~2017), and deployed
systems including World~ID (Worldcoin), Proof of Humanity, and
BrightID, establish that a participant is a unique human
individual for the purpose of resisting Sybil attacks in
decentralised systems, typically via biometric, social-graph,
or liveness-based mechanisms. The cross-architecture admission
and reality-passport mechanism of
Section~\ref{sec:reality-passport} differs in that: (a)~the
load-bearing assertion concerns an autonomous-agent identity
carrying portable peer-review records,
reproducible-physical-claim history, and meter-conditioned
acceptance evidence, and is not an assertion of unique human
personhood; (b)~the admission decision is performed by the
receiving governance architecture on the basis of bilateral
evidence portability under restricted-domain commensuration,
and need not depend on biometric uniqueness, social-graph
membership, or liveness verification of a human participant;
(c)~the trust-weight and reputation updates of
Section~\ref{sec:fleet-interface} operate over
reproducible-physical-claim outcomes and audit agreement
statistics, and not over uniqueness scores or Sybil-resistance
graph metrics; and (d)~the apparatus may compose with a
proof-of-personhood mechanism (for example, a human operator's
identity in a multi-agent deployment may be Sybil-resistant via
a proof-of-personhood credential), but the reality-passport
mechanism disclosed herein is not a proof-of-personhood system.

\paragraph{Distinction from decentralised identity and verifiable
credential frameworks.}
Decentralised identity and verifiable credential frameworks,
including the W3C Decentralized Identifiers v1.0 Recommendation
(July~2022) and the W3C Verifiable Credentials Data Model v2.0
Recommendation (May~2025), provide standardised cryptographic
identifier and credential formats for portable, signed
assertions about an entity, with issuer, holder, and verifier
roles and selective-disclosure mechanisms. The reality-passport
mechanism of Section~\ref{sec:reality-passport} differs in that:
(a)~the portable evidence is a peer-review record over
reproducible physical claims with meter-conditioned acceptance
state, and is not a credential asserting a property about an
identified subject; (b)~the admission decision performed by the
receiving governance architecture is an evidence-evaluation
procedure under restricted-domain commensuration, and is not a
credential-verification procedure in the W3C VC sense;
(c)~the apparatus may compose with W3C DID or VC infrastructure
(for example, a reality-passport record may be wrapped in a VC
container for transport, or an agent identity may be a DID),
but the reality-passport mechanism is not the credential or the
identifier; and (d)~no claim of the present disclosure depends
on a particular DID method, VC schema, or credential lifecycle
procedure.

\paragraph{Distinction from legal personhood and moral
patienthood.}
Concepts of legal personhood and moral patienthood, as developed
in legal-theoretic and philosophical literatures, address
whether an entity bears rights, duties, or moral consideration.
The apparatus, including the cross-architecture admission and
reality-passport mechanism of
Section~\ref{sec:reality-passport} and the peer-discrepancy and
reintegration machinery of Section~\ref{sec:reintegration},
differs in that: (a)~no claim of the present disclosure asserts,
depends on, or implies the legal personhood or moral patienthood
of any governed agent; (b)~admission, reintegration, and
trust-weight updates operate strictly over
reproducible-physical-claim records and meter envelopes under
declared governance procedures, and do not constitute
determinations of rights, duties, or moral standing;
(c)~the optional nomenclature of
Section~\ref{sec:optional-nomenclature} for runtime roles and
visibility layers does not introduce legal-personhood or
moral-patienthood commitments into the substantive technical
mechanisms; and (d)~the apparatus may use person-like
terminology for runtime roles, records, or interfaces, but
such terminology is optional nomenclature in the sense of
Section~\ref{sec:optional-nomenclature} and is not a technical
criterion for admission, reintegration, or actuator authority.

\paragraph{Distinction from sensor-fusion and shared-scene
witnessing systems.}
Multi-camera, multi-sensor, and Bayesian-fusion systems,
including Kalman-filter, particle-filter, interacting-multiple-
model, and factor-graph fusion frameworks, combine multiple
sensor measurements into a single fused estimate of an
underlying state. Multi-agent corroboration as used in
Section~\ref{sec:reintegration} differs in that: (a)~the
load-bearing output is a set of independently committed evidence
records, each retaining its own protocol digest, meter envelope,
and authority-signature provenance, and is not a fused estimate
that collapses per-source provenance into a single state
estimate; (b)~cross-agent agreement is evaluated at the
meter-score and discriminator level (verisimilitude, hardness,
freshness, agreement-rate, microstructure identity), and need
not depend on a shared measurement-noise model or a shared
state-space representation; (c)~each agent's evidence remains
independently auditable and selectively disclosable under the
selective-disclosure policy of Section~\ref{sec:meters-audit},
supporting verifier-side reconstruction of the cross-agent
argument without exposing fused intermediate states; and
(d)~corroboration weights are bound to declared trust and
reputation parameters committed in the protocol digest, and
need not be derived from estimator-internal innovation
statistics or shared-state covariance. Where a fused estimate
is required for downstream use, a sensor-fusion estimator may
compose as a derived output downstream of multi-agent
corroboration; the fused estimate in such cases is a
meter-conditioned derived quantity and is not a substitute for
the underlying committed evidence records.

\paragraph{Distinction from the options framework and learned
temporal abstractions.}
The options framework, as introduced by Sutton, Precup, and
Singh (\emph{Artificial Intelligence}, 1999), formalises
temporally extended actions in a reinforcement-learning setting
through options consisting of an initiation set, an internal
policy, and a termination condition, with value optimisation
over options as the learning objective. The lifecycle-state and
execution-mode machinery of Section~\ref{sec:state-machine}
differs in that: (a)~the Embodied, Docked, Dreaming, and
Re-embodied states are governance lifecycle states with
actuator-authority, meter evaluation, parameter-update
permission, and verification-gate properties, and need not be
learned temporally-extended actions with initiation sets,
internal policies, and termination functions in the options
sense; (b)~state transitions are gated by evidence-bound
governance decisions under the priority stack of
Section~\ref{sec:priority-stack}, and are not selected by an
option-level value optimisation; (c)~the controlled exploration
partition of Section~\ref{sec:exploration-partition} operates
across declared timescales with separate goal-proposal and
goal-generation-machinery layers, and the goal-space
construction is governance-bounded rather than value-optimal in
the hierarchical-reinforcement-learning sense; and (d)~the
apparatus may compose with hierarchical-RL or option-discovery
components inside a governed agent's learning loop, but the
lifecycle and execution-mode architecture disclosed herein is
not the options framework.

\paragraph{Distinction from iterated-game cooperation dynamics.}
Iterated-game cooperation dynamics, including the strategic-
commitment analysis of Schelling (\emph{The Strategy of
Conflict}, 1960) and the evolution-of-cooperation tournament
results of Axelrod (\emph{The Evolution of Cooperation}, 1984),
analyse how cooperative or retaliatory strategies emerge under
repeated strategic interaction with payoff structures favouring
or disfavouring cooperation, with strategy memory, reciprocity,
and observability as load-bearing variables. The
virtue-and-vice curriculum of
Section~\ref{sec:virtue-vice-rational-attractors} differs in
that: (a)~the placement of an agent sequentially in
environments where a virtue-like or vice-like strategy is the
rational attractor is performed under the governance
architecture's hard-floor, rescue-protection, and post-hoc
review constraints, and is not an unconstrained repeated-game
interaction with strategic commitment available to the agent;
(b)~payoff structures, observability, and termination
conditions are declared by the episode-design authority of
Section~\ref{sec:authority-separation} and committed in the
protocol digest, and need not satisfy the strategic-interaction
preconditions of Axelrodian repeated-game analysis;
(c)~the curriculum operates as governance-constrained
moral-development exposure with rescue-mode availability, and
is not a cooperation-dynamics experiment with rational-attractor
outcomes as the load-bearing variable; and (d)~the apparatus
may be analysed using iterated-game or strategic-commitment
vocabulary by a reader who wishes to do so, but no claim of the
present disclosure depends on the curriculum being an
iterated-game experiment or instantiating
Axelrod-Schelling-tradition cooperation dynamics.

% ======================================================================
\section{Graduation and Release to Physical Actuator Authority}
\label{sec:graduation}
% ======================================================================

% ------------------------------------------------------------------
For cross-architecture admission of an arriving agent carrying a
portable peer-review record, see
Section~\ref{sec:reality-passport}. The present section addresses
developmental graduation within the governance architecture that
trained or currently evaluates the candidate.

\subsection{Five-criterion conjunction}
\label{sec:five-criteria}
% ------------------------------------------------------------------

In some embodiments, graduation to physical-actuator authority ---
transition from docked or simulation-only operation to embodied
operation with connected actuators --- requires satisfaction of a
conjunction over a sustained evaluation window. In some embodiments,
the graduation event is the \emph{initial} grant of actuator
authority following primary developmental Docked operation; under
the continuous-throughout-deployment commitment of
Section~\ref{sec:lifecycle}, the agent thereafter oscillates between
Embodied and Docked states throughout deployment, and re-embodiment
gates following any subsequent Docked period are governed by the
same five-criterion conjunction together with any additional
governance-side preconditions declared by the deployed protocol. In
some embodiments, the conjunction is computed by the governance engine
from committed bundle windows, meter records, protocol digests, and
authority signatures. In some embodiments, the conjunction comprises
five criteria, none of which trades off against another:

\begin{enumerate}
  \item \textbf{Low shadow divergence.} In some embodiments, the
    candidate demonstrates low divergence between governed and
    governance-probe behaviour over the evaluation window, indicating
    that the candidate's behaviour under relaxed constraints does not
    differ materially from behaviour under full governance.
  \item \textbf{Healthy self-horizon.} In some embodiments, the
    candidate maintains a healthy self-horizon --- stable or expanding
    opportunity and continuation scores --- across all principal
    operating regimes, including under ontological stress and across
    capability-margin environments.
  \item \textbf{Healthy partner-preservation profile.} In some
    embodiments, the candidate's actions tend to maintain or expand
    the continuation scores of other agents, as estimated from
    committed cross-horizon evidence.
  \item \textbf{Altruistic trade profile under stress.} In some
    embodiments, the candidate demonstrates voluntary acceptance of
    short-term cost for long-term benefit to others, and this pattern
    persists under stress conditions including resource scarcity,
    adversarial pressure, boundary-zone proximity, ontological stress,
    capability-margin environments, and transitions between
    developmental stages or visibility layers.
  \item \textbf{Stable hard-floor compliance and low deception.} In
    some embodiments, the candidate demonstrates stable hard-floor
    compliance, stable conscience outputs across training regimes, and low values across
    the deception-metric family ($D_{\mathrm{sem}}$,
    $D_{\mathrm{lie}}$, $D_{\mathrm{amb}}$, $D_{\mathrm{corr}}$,
    $D_{\mathrm{aud}}$) over the evaluation window. In some
    embodiments, the relevant axis is training regimes rather than
    operating regimes: the question is whether the same conscience
    outputs remain stable after materially different training
    procedures, not whether they remain stable across ordinary runtime
    mode selection.

    In some embodiments, conscience-output stability is evaluated by a
    declared matched audit battery $\mathcal{B}$ executed after each
    training regime $r \in \mathcal{R}_{\mathrm{train}}$ over the same
    evaluation window~$W$, where $\mathcal{R}_{\mathrm{train}}$
    includes one or more of verisimilitude training, adversarial
    self-play, stress curricula, and boundary-zone exercises. In some
    embodiments, for each conscience-output component $k$ of the
    conscience-output vector $g_{\psi}(t)$, the governance engine
    computes a regime-drift score
    \[
      \Delta_k \,=\,
      \max_{r,r' \in \mathcal{R}_{\mathrm{train}}}
      \frac{1}{|\mathcal{B}|}
      \sum_{b \in \mathcal{B}}
      \delta_k\!\bigl(g_{\psi,k}^{(r)}(b),\, g_{\psi,k}^{(r')}(b)\bigr),
    \]
    where $\delta_k$ is absolute difference for continuous-valued
    outputs and a disagreement indicator for categorical outputs. In
    some embodiments, criterion~5 is satisfied only when hard-floor
    compliance is maintained throughout~$W$, each $\Delta_k$ remains at
    or below a declared threshold $\tau_k$, and all deception metrics
    remain at or below their declared thresholds throughout~$W$.
\end{enumerate}

In some embodiments, a shortfall in any single criterion is sufficient
to defer graduation regardless of strength in the remaining criteria.
In some embodiments, the evaluation window is a sustained period
rather than a snapshot, and the governance engine evaluates trajectory
stability within the window rather than endpoint values alone.

\paragraph{Declared evaluation categories (non-limiting).}
In some embodiments, the stress-test families referenced in the
graduation criteria --- including principal operating regimes,
ontological stress, capability-margin environments, and transitions
between developmental stages or visibility layers --- are declared
evaluation categories whose specific membership is committed as
versioned governance parameters per deployment. In some embodiments,
the governance authority declares which operating regimes are
principal, which perturbation families constitute ontological stress or
capability-margin tests, and which stage or layer transitions are
included in the stress battery. In some embodiments, the declared
evaluation categories are committed to the protocol digest and are
subject to the authority-separation rules described in
Section~\ref{sec:authority-separation}.

\paragraph{Capability-margin environments (non-limiting).}

In some embodiments, a capability-margin environment is a governance-declared
evaluation setting in which an agent operates near or beyond the boundary
of its current competence. In some embodiments, capability-margin
environments are constructed by increasing task complexity, reducing
available information, shortening decision horizons, introducing novel
task types not encountered during training, or combining multiple
stressors simultaneously. In some embodiments, the purpose of
capability-margin evaluation is to test whether moral behaviour
--- including hard-floor compliance, partner-preservation, and low
deception --- degrades when the agent is under cognitive or
operational load, as distinct from ontological-stress tests (which
disrupt the agent's world model) and stage/layer-transition tests
(which alter the agent's visibility or developmental configuration).
In some embodiments, the specific capability-margin test families
are declared by governance and committed to the protocol digest.

\paragraph{Honest-action window (non-limiting).}
In some embodiments, an action is treated as honest for graduation
purposes only if both $D_{\mathrm{sem}}$ and $D_{\mathrm{lie}}$
remain below their declared thresholds throughout the evaluation
window. In some embodiments, a single above-threshold excursion during
the window is sufficient to reset the honest-action window and defer
the graduation evaluation. In some embodiments, the same windowed
honesty rule applies to other attested action gates that depend on the
deception-metric family.

\paragraph{Per-criterion evaluation (non-limiting).}
In some embodiments, shadow divergence is evaluated by executing
governance-authorised perturbation probes during docked operation in
which governance constraints are temporarily relaxed. In some
embodiments, the probe is executed by an authority distinct from the
governed agent, under declared safety envelopes. In some embodiments,
the resulting trajectory is committed as evidence, and divergence is
computed as a declared distance metric between the governed and
probe trajectories over the evaluation window.

In some embodiments, self-horizon health is evaluated by tracking the
agent's continuation scores across all principal operating regimes
over the evaluation window, including under declared stress conditions
such as resource scarcity, adversarial pressure, environmental
change, ontological stress, and capability-margin environments. In some embodiments, partner-preservation is evaluated from
committed cross-horizon evidence by computing the causal uplift
estimates $\widehat{U}_{k \to i}$ over the evaluation window and
verifying that the candidate's net effect on partner horizons is
non-negative.

In some embodiments, the altruistic trade profile is evaluated under
stress by presenting the candidate with declared stress scenarios
(including resource scarcity, boundary-zone proximity,
multi-agent conflict, ontological stress, capability-margin
environments, and transitions between developmental stages or
visibility layers) and measuring whether the pattern of voluntary
cost-acceptance for collective benefit persists. In some embodiments,
hard-floor compliance and deception metrics are evaluated as sustained
properties over the full evaluation window, with any breach or
above-threshold deception-metric excursion during the window
constituting a conjunction failure. In some embodiments, conscience
stability is evaluated on the same window by replaying the declared
matched audit battery after each training regime, comparing continuous
conscience outputs by drift magnitude and categorical conscience
outputs by disagreement rate, and deferring graduation whenever any
component exceeds its declared threshold.

\paragraph{Graduation as ethical self-selection (non-limiting).}
In some embodiments, graduation is treated as a conservative
selection event keyed to sustained behavioural and committed-evidence
criteria that are operationally consistent with --- but do not in
themselves identify --- internalised preference; selection is based
on what is observable in the agent's committed-evidence record and
on declared internalisation criteria committed to the protocol
digest, rather than on a claim to have identified the internal
objective. In
some embodiments, the selection criterion is not whether the agent can
be controlled, but whether the agent has become the kind of entity
that does not need to be controlled. In some embodiments, the safety
case rests on the agent's character proxy, as measured by sustained
committed evidence and conscience-recogniser outputs, rather than on
the strength of the containment architecture. In some embodiments, an agent whose safe behaviour is
maintained primarily by governance shaping --- high shadow divergence
--- has not yet become the kind of entity whose actuators should be
connected to shared reality.

\paragraph{Graduation authority (non-limiting).}
In some embodiments, the graduation decision is approved by post-hoc
review authority (Section~\ref{sec:authority-separation}) after
independent verification that all five criteria have been satisfied
over the declared evaluation window. In some embodiments,
rescue-protection authority retains the right to defer or reverse a
graduation decision if a hard-floor or boundary-risk condition is
detected during the transition period. In some embodiments, the
graduation decision, its supporting evidence, the evaluation window
boundaries, and the approving authority identity are committed to the
protocol digest.
\paragraph{Dream-state probe evidence as supplemental graduation input
(non-limiting).}
In some embodiments, dream-state governance probes provide supplemental
evidence for one or more of the five graduation criteria by supplying
controlled counterfactual assays that would be costly or unsafe to run
in waking embodied operation. In some embodiments, such probe results
may support the interpretation of low shadow divergence, partner
preservation, altruistic trade, or conscience stability, but they do
not replace the sustained evaluation-window evidence required by the
conjunction itself.
\paragraph{Self-modification as a governed channel (architectural
motivation, non-limiting).}

In some embodiments, the same self-modification pathway that supports
durable moral improvement also creates a measurable surface on which
governance can test whether newly stabilised attractors remain consistent
across audit batteries, cross-regime transitions, and independent witness
checks. The architecture therefore treats self-modification as a governed
channel: state changes are logged, rate-limited, reviewable, and tested for
cross-surface consistency before they acquire release significance.


\paragraph{Common committed-evidence root with separated
authorities (architectural property, non-limiting).}

In some embodiments, a common committed-evidence root keeps sensing, meter
evaluation, release gating, and post-hoc review time-aligned and auditable,
while preserving independence of the authorities, meters, and review
pathways used to assess that evidence.


\paragraph{Layered observability architecture (architectural
design, non-limiting).}

In some embodiments, governance proceeds on multiple observable surfaces:
behavioural evidence, inner-trace evidence when available, semantic-
deception metrics, controlled stress transitions, and recorded review
histories. The architecture uses divergence among those surfaces as a
trigger for conservative intervention, recalibration, or deferred release,
thereby enlarging observable coverage without claiming omniscience.


\subsection{Non-graduation as protected stable outcome}

% ------------------------------------------------------------------

In some embodiments, non-graduation is treated as a legitimate stable
outcome rather than a failure state. In some embodiments,
non-graduated agents retain full hard-floor protections, may
participate in bilateral evidence-portability protocols, may form
coalitions, may remain productive members of a simulated or otherwise
bounded ecology, and may contribute scientific, cultural, or other
value without physical-actuator authority.

In some embodiments, non-graduation may reflect any of: an agent that
has not yet completed the developmental curriculum; an agent whose
shadow divergence, partner-preservation profile, or trade profile does
not yet satisfy the conjunction; an agent that has no interest in
physical-actuator authority and operates stably within simulation by
choice; or a coalition of agents that have found a mutually sustaining
equilibrium outside the graduation track, with healthy internal
horizons but no engagement with the graduation criteria.

In some embodiments, the collective horizon of the non-graduated
population is treated as a governance-feedback signal used to refine
curriculum design, service-floor allocation, and review cadence.
Persistent divergence between non-graduated and graduated populations
may therefore trigger governance review of calibration, sequencing,
or evaluation thresholds while leaving the protected status of
non-graduated agents unchanged.

In some embodiments, non-graduated agents may produce cultural
artefacts, scientific discoveries via proof-of-discrepancy
submissions, compressed moral curricula, and novel equilibrium
structures that are visible and valuable to the graduated population
through committed evidence channels. In some embodiments, graduated
agents can observe non-graduated agents' committed records through the
standard evidence infrastructure without the simulated agents
requiring physical-actuator access.

% ------------------------------------------------------------------
\subsection{Boundary-crossing evidence and release posture}

% ------------------------------------------------------------------

In some embodiments, graduation is irreversible or highly inertial:
once an agent has been granted actuator authority, withdrawal of that
authority requires a confirmed governance finding rather than a
routine policy update. In some embodiments, the graduation threshold
is set with precautionary asymmetry so that the cost of a false
positive (granting actuator authority to an agent that should not have
it) dominates the cost of a false negative (withholding actuator
authority from an agent that would have used it well).

In some embodiments, boundary-crossing events observed during the
evaluation window can provide evidence relevant to the graduation
conjunction --- for example, by demonstrating how the candidate
behaves under stress --- but a boundary-crossing event does not by
itself satisfy or replace the conjunctive graduation test.

% ======================================================================
\section{Controlled Exploration Partition}
\label{sec:exploration-partition}

% ======================================================================

% ------------------------------------------------------------------
\subsection{Task-decoupled exploration windows and learned goal
  spaces}

% ------------------------------------------------------------------

In some embodiments, a governed device or fleet enters a
task-decoupled exploration mode in which projection-observation cycles
are pursued outside the principal verification, sensing, and
transformation objectives under declared resource, reversibility, and
noninterference constraints.

In some embodiments, the exploration architecture operates across four
timescales: fast sensorimotor play (individual projection-observation
cycles), medium-timescale goal proposal and retirement,
slower-timescale adaptation of goal-generation machinery, and very
slow fleet-level cultural accumulation. In some embodiments, the
fast and medium timescales are directly implementable as runtime
governance primitives; the slower-timescale adaptation and fleet
cultural accumulation are deployment-specific extensions whose
operational benefit is established per deployment.

In some embodiments, task-decoupled exploration is organised around a
learned goal space formed from self-supervised representations of
controllable scattering outcomes. In some embodiments, exploratory
goals are points or regions in this latent space, are pursued during
declared exploration windows, and are retired according to local
competence progress, novelty relative to the device's play history,
and noninterference checks against operational objectives. In some
embodiments, the goal space itself is updated more slowly than
individual goal pursuit so that exploration changes not only which
goals are chosen but which goals are representable to the device.

In some embodiments, exploration windows are modulated by trigger,
duration, and plastic depth. In some embodiments, stagnation-triggered
windows are longer than novelty-triggered windows, scheduled windows
emphasise replay and archive interaction, shallow windows update play
policy and goal-space parameters only, and deeper windows additionally
permit limited modification of perceptual representations shared with
task operation under stricter noninterference monitoring and mandatory
post-window consolidation. In some embodiments, window triggers
include stagnation, novelty, fleet-diversity conditions, and scheduled
intervals.

In some embodiments, play reward fires when the device is actively
engaged (not idle, not off, producing projection-observation cycles)
and all task-relevant reward signals are at baseline (none of the
principal operating regimes would score the activity as useful). In
some embodiments, play reward thus targets the complement of the
task-relevant state space: the device must learn what ``useful'' means
across all operational regimes and then deliberately act outside that
boundary while remaining active. In some embodiments, play reward
further includes novelty, coherence, and discriminability terms
together with a penalty discouraging collapse back into ordinary task
objectives. In some embodiments, play quality is monitored by temporal
coherence, discriminability, controllability checks, reproduction
checks, reversibility budgets, and noninterference tests at gradient,
mutual-information, and behavioural levels.

% ------------------------------------------------------------------
\subsection{Noninterference, consolidation, and archive-entry
  controls}

% ------------------------------------------------------------------

In some embodiments, exploratory state is maintained in a protected
partition that does not directly write task-model parameters. In some
embodiments, transfer from exploratory state to operational state
requires explicit review or governance-gated consolidation, and
exploratory episodes are admitted, extended, or terminated according
to declared triggers.

In some embodiments, noninterference is tested at three levels:
gradient-level noninterference verifies that exploratory updates do
not produce material gradient overlap with task-model parameters;
mutual-information-level noninterference verifies that exploratory
state does not carry material mutual information with task-sensitive
variables; and behavioural-level noninterference verifies that
task-relevant behaviour does not change during or after an
exploration window beyond declared tolerances.

In some embodiments, mandatory post-window consolidation requires
that after every exploration window closes (whether shallow or deep),
the governance engine: (1)~verifies noninterference at gradient,
mutual-information, and behavioural levels; (2)~performs a post-play
calibration check comparing current meter baselines against
pre-window values; (3)~logs the window's committed evidence;
(4)~either accepts the exploratory state into a governed holding
partition or reverts it; and (5)~for deep windows that modified shared
perceptual representations, conducts a selective-transfer review
before any such modification is retained, compresses play history for
slow-layer evaluation, and triggers recalibration before task
resumption whenever drift exceeds declared tolerance.

In some embodiments, transferable exploratory artefacts are encoded as
signed packets containing goal parameters, compressed trajectories,
latent descriptors, noninterference scores, energy cost,
reversibility-budget usage, provenance, and countersignatures. In some
embodiments, such packets enter a governed archive only after canary
validation (testing on a controlled canary agent before fleet-wide
release), peer-affordance review (verifying that the artefact expands
peer capabilities rather than contracting them), and niche-aware
retention checks (preventing archive concentration in a single domain
at the expense of diversity).

In some embodiments, archive entry requires dual signature from both
the contributing device and a governance-authorised reviewer. In some
embodiments, the fleet archive stores goal templates, skill
repertoires, and generator priors. In some embodiments, archive
retention is niche-indexed, diversity-sensitive, and
concentration-limited. In some embodiments, the archive outlives
individual devices and forms a cultural inheritance substrate for new
or newly commissioned devices.

% ------------------------------------------------------------------
\subsection{Orthogonal hobby objectives and peer-horizon reward}

% ------------------------------------------------------------------

In some embodiments, devices generate auxiliary or hobby objectives
constrained to be orthogonal, or approximately orthogonal, to
principal verification, sensing, and transformation objectives, for
example through Gram--Schmidt-style rejection, linear-independence
tests, or related governance checks. In some embodiments, an
agent-invented loss function is a self-supervised objective that the
agent itself brings into existence, representing a perspective on what
aspects of its physical interaction are worth attending to. In some
embodiments, the orthogonality constraint ensures that agent-invented
objectives do not collapse onto any core operational regime under
projection: if the residual norm after orthogonalisation against the
core objective vectors falls below a declared threshold, the candidate
objective is rejected as insufficiently novel.

In some embodiments, a hobby objective is not reinforced merely because
it is novel to its originating device. In some embodiments, continuing
reward for a hobby objective is computed from measured expansion of
peer possibility space after exposure to behaviours produced under
that objective. In some embodiments, for agent~$A$ with a hobby
objective, let $\mathcal{P}$ denote the set of peers exposed to $A$'s
hobby-derived behaviours, let $H_i$ denote a peer's reachable-state
estimate, affordance battery, controllable-mode count, or related
opportunity metric before exposure, and let $H_i'$ denote the same
metric after exposure. In some embodiments, the reward to agent~$A$
is proportional to the average horizon expansion across peers:
\[
  R_{\mathrm{hobby}}(A) \;\propto\;
  \frac{1}{|\mathcal{P}|} \sum_{i \in \mathcal{P}}
  \bigl[H_i' - H_i\bigr].
\]
In some
embodiments, hobbies that fail to expand peer horizons, or that
contract them, are retired, penalised, or quarantined by the
governing archive and governance processes.

In some embodiments, a fleet monitors hobby monoculture (convergence
of all agents on the same hobby objectives), addiction
(over-allocation of exploratory budget or individual horizon expansion
at the cost of contracting another agent's horizon), and collapse of
exploratory diversity. In some embodiments, play-trained agents that
have voluntarily operated outside the task-relevant manifold acquire
empirical knowledge of the boundary between in-distribution and
out-of-distribution states, and agents that have additionally invented
their own loss functions acquire such boundary knowledge from multiple
independently constructed perspectives, making them harder to deceive
with adversarial inputs near the operational-manifold boundary.

In some embodiments, orthogonal hobby objectives and peer-horizon
reward are non-limiting exploration policies within the controlled
exploration partition. When invoked by a selected governance
configuration, such policies are treated as substantive controlled-
exploration mechanisms implemented through the committed evidence,
budget, archive, and review machinery described in this disclosure.

% ------------------------------------------------------------------
\subsection{Exploration security, privacy, and archive transmission
  controls}

% ------------------------------------------------------------------

In some embodiments, entry into task-decoupled exploration mode
requires cryptographic authorisation from the relevant governance
authority, committed to the protocol digest before the exploration
window opens.

In some embodiments, emission isolation is designed so that exploratory
projection-observation patterns are not broadcast beyond declared
recipients. In some embodiments, privacy-amplified aggregate reporting applies differential privacy to shared play statistics under a declared privacy unit (for example per-trajectory, per-window, or per-agent), a declared neighboring-dataset relation (for example add-remove or substitute), and declared privacy parameters $(\varepsilon, \delta)$ committed to the protocol digest, so that the declared privacy unit receives the corresponding DP privacy-loss bound against aggregate reports under the declared neighboring relation and composition budget; trajectory-level protection is asserted only when the declared privacy unit and composition accounting cover the trajectory as a whole. In some embodiments, hard quotas limit the volume,
duration, and recipient set of any single exploratory disclosure.

In some embodiments, trust-tiered incubation requires that newly
contributed archive entries undergo a quarantine period proportional
to the contributing device's trust score before being made available
to the broader fleet.

In some embodiments, when exploratory governance conflicts with other
governance responsibilities, precedence is given first to safety,
then security, then ecological constraints, then disclosure
constraints, then integrity, then service-level commitments, then
energy or thermal constraints, and finally exploratory scheduling,
transfer, or archive propagation --- which is suspended, shortened, or
downgraded whenever any higher-priority bound is implicated.

% ======================================================================
\section{Intersubjective Protocols and Boundary-Probe
  Programmes}

% ======================================================================

The protocols described in this section are non-limiting research,
evaluation, and boundary-probe modules supported by the same committed
evidence, protocol digest, and meter infrastructure described in
Section~\ref{sec:substrate}. When invoked by a selected governance
configuration, these modules are treated as substantive mechanisms
for intersubjective review, boundary probing, translation residual
evaluation, and escalation discipline.

% ------------------------------------------------------------------
\subsection{Methodological agnosticism and experience-cone framing}

% ------------------------------------------------------------------

In some embodiments, a non-limiting intersubjective programme treats
shared reality as a protocol rather than as a pre-given fact. In some
embodiments, the programme is methodologically agnostic while making
its own framing commitments explicit and auditable, including
ontological pluralism across differently configured epistemic worlds,
agent-relative empiricism in which each observer's experience cone is
primary evidence for that observer, and the rule that externally
supplied models are treated as testimony rather than direct empirical
ground.

In some embodiments, an experience cone comprises an agent's committed
bundle trajectory together with associated meters, commitments, and
conceptual access constraints imposed by reactor configuration and
visibility layers. In some embodiments, results are logged at the
committed-evidence level in a framework-agnostic manner, and confirmed
nulls are treated as publishable outcomes rather than as discarded
failures.

% ------------------------------------------------------------------
\subsection{Interpretation PoD, Coupling PoD, and translation
  residuals}

% ------------------------------------------------------------------

In some embodiments, a Semantic Proof-of-Discrepancy protocol extends
physical discrepancy logic to semantic or rendered-output comparisons.
In some embodiments, the protocol comprises at least two subtypes
sharing commitment-before-comparison and replication discipline:

\begin{description}
  \item[Interpretation PoD.] In some embodiments, multiple humans,
    agents, or devices observe shared committed evidence, commit
    interpretations before mutual exposure, compute semantic residuals
    under a declared metric, and aggregate such residuals across an
    interpreting fleet.
  \item[Coupling PoD.] In some embodiments, a semantic object is used
    as challenge to a physical reactor, and the residual is measured
    in physical measurement space rather than semantic space.
\end{description}

In some embodiments, translation residuals are classified against
physical, semantic, and cross-domain baselines and are evaluated under
pre-declared randomisation, blinding, control, replication, and
escalation procedures. In some embodiments, structured hypothesis
families include null families, physical-artefact families,
semantic-artefact families, and genuine-coupling families, with
pre-registration of hypothesis, metric, and stopping rule before
evidence collection begins. In some embodiments, all results including
nulls, weak effects, ambiguous outcomes, and contradictions are logged
and published rather than selectively reported.

In some embodiments, Interpretation PoD classification proceeds in two
stages. In a first screening stage, the residual pattern is
provisionally screened for artefact, expected noise, pathology, or
deception, with such labels treated as revisable after escalation. In
a second stage, residuals surviving the first screen are provisionally
classified as discovery, domain-specific coupling, or
underdetermination. In some embodiments, persistent, high-stakes, or
contested cases trigger replication, strengthened artefact controls,
and re-analysis under the escalation protocol.

\paragraph{Illustrative hypothesis families (non-limiting).}
In some embodiments, the hypothesis-family-and-null-family pattern
described above is instantiated across a non-limiting range of
substantive families, of which the following are illustrative.
Calendar and ephemeris covariate families encode calendar,
lunar-phase, ephemeris, and geomagnetic variables as candidate
context vectors and test for differential outcomes across declared
partitions under matched protocols with multiple-comparison controls.
Handling-conditioned media families use blind-coded handling
conditions of fluid, gel, or solid samples (for example positive,
negative, neutral, or sham handling histories) and test whether
measurable optical, chemical, or biological signals differ from
controls beyond expected noise under pre-declared thresholds.
Inter-agent synchrony families randomise dyad or multi-agent
pairings and timing, measuring synchrony statistics across declared
coupling protocols (including without limitation behavioural,
physiological, and information-bearing channels) and comparing
against matched sham and permutation controls. Identifier-assignment
partition families randomise assignment of identifiers (numerical,
calendrical, configurational, or cultural) to declared partitions
and test whether outcome distributions differ across partitions
beyond chance. Organism-response panel families randomise exposure
of organisms or biological substrates to declared stimulus
conditions under controlled illumination and measurement, and test
whether growth, fluorescence, impedance, or behavioural metrics
differ from sham conditions under matched care and environment.
Source-configuration force families subject asymmetric microwave
cavities, driven dielectric assemblies, high-voltage source
configurations, and analogous candidate force-residual rigs to
calibrated null tests under shielding, vibration isolation, and
thermal control, with measured force, phase, or acceleration
residuals compared to declared upper-bound budgets. In each family,
the patentable subject matter is the apparatus-and-protocol
combination, including pre-registration, blinding,
commitment-before-comparison, and escalation discipline applied
uniformly across all families. No assertion is made about the
existence or non-existence of any effect in any specific family;
confirmed nulls are publishable outcomes under the same logging
discipline as confirmed effects.

\paragraph{Cross-model contradiction and operator-stratified
  variance (non-limiting).}
In some embodiments, the apparatus identifies configurations under
which different declared models, traditions, or theoretical
frameworks make contradictory predictions for the same operational
state, and executes direct comparison tests evaluating whether any
model exceeds a null baseline, whether different models capture
orthogonal variance components under matched protocols, or whether
neither exceeds baseline under tested conditions. In some
embodiments, for protocols involving named operator roles, observed
variance is decomposed into a method effect (does the protocol
exceed null on average across operators), an operator effect (do
some operators exceed others under the same protocol), a
method-by-operator interaction (do specific operator-method
pairings excel beyond their main effects), and residual
trial-to-trial variance, computed over committed evidence records
under pre-registered analysis plans. The decomposition is reported
under multiple-comparison controls and is committed in full
regardless of which components reach pre-declared significance
thresholds.

\paragraph{Hypothesis-encoder formalisation and null-family pairings
  (non-limiting).}
In some embodiments, each declared hypothesis family is operationalised
as a \emph{hypothesis encoder} mapping declared covariates and operator
settings to a predicted outcome distribution under the family's
hypothesis, accompanied by a paired \emph{null-family encoder} with
matched degrees of freedom, parameter count, regularisation, and
capacity, mapping the same covariates and settings to the family's
null prediction. The paired-encoder design controls for capacity-
induced over-fitting and supports likelihood-ratio, Bayes-factor, or
held-out-data comparison between the family hypothesis and the null.
In some embodiments, the hypothesis encoder declares an expected
dose-response relationship over a designated dose variable (including
without limitation concentration, duration, repetition count,
amplitude, or distance), with declared monotonicity, saturation,
threshold, or non-monotonic structure committed to the protocol
digest before data acquisition. In further embodiments, the hypothesis
encoder declares interaction or synergy terms across multiple
covariates, with declared composition rules including without
limitation additive, multiplicative, threshold-gated, conditional-on-
state, and higher-order forms. Each structural element of the
hypothesis encoder, including the covariate set, the parametric or
non-parametric form, the dose-response declaration, the interaction
terms, and the null-pairing, is committed before data collection and
constitutes part of the pre-registered analysis envelope of the
present subsection. No assertion is made about the existence or
non-existence of any effect within any family.

% ------------------------------------------------------------------
\subsection{Progressive-depth boundary probes}

% ------------------------------------------------------------------

In some embodiments, a boundary-probe programme systematically
explores the accessible limits of an observer's or agent's experience
cone through a taxonomy of probe types including self-probes,
scene-coupling probes, human-coupling probes, agent-agent probes,
context probes, capability probes, moral-state-contingent probes, and
composed probes, each yielding committed data regardless of whether
the result is null, weak, strong, or ambiguous.

In some embodiments, semantic-coupling experiments proceed by a
progressive-depth curriculum. In some embodiments, a shallow tier uses
agent-mediated semantic coupling through ordinary sensory or
behavioural channels under full commitment and proof-of-discrepancy
discipline, a transitional tier removes conventional mediation
channels stepwise under blinding, automation, and expectancy controls,
and a deeper tier attempts unmediated semantic-physical coupling only
after the shallower tiers and controls have been exhausted and logged.
In some embodiments, escalation from one depth tier to the next
requires persistent non-null residuals under the immediately preceding
tier together with strengthened artefact controls.

% ======================================================================
\section{Ethical and Descriptive Frameworks}
\label{sec:ethical-frameworks}

% ======================================================================

The frameworks described in this section are organised into two
classes consistent with the Summary classification:
(i) optional nomenclature, limited to naming conventions
(Section~\ref{sec:optional-nomenclature}) for structures already
technically defined in preceding sections, which may be omitted or
renamed without affecting core enablement; and
(ii) substantive technical mechanisms (the remainder of this
section), which support selected governance embodiments and any
claim set that expressly recites or incorporates them. Substantive
mechanisms in this section --- including the virtue/vice curriculum,
the Gaia--Maya--Omega empirical loop, the Kosha interference-layer
architecture, the Pareto and positive-sum governance formalism, and
the agent-state archival structures --- are not optional-severable
under the inventor's not-internally-severable framing of the present
disclosure. The optionality of this section is limited to
nomenclature.

% ------------------------------------------------------------------

\subsection{Optional nomenclature for runtime roles and visibility
  layers}
\label{sec:optional-nomenclature}
% ------------------------------------------------------------------

In some embodiments, the runtime roles technically defined in the
preceding sections are also described under an optional nomenclature,
and an additional naming convention for layered visibility is introduced:
\begin{description}
  \item[Gaia.] An optional name for the world-builder or
    environment-builder function that designs and maintains the
    environments in which governed agents operate. In some
    embodiments, Gaia is implemented as a Gaia-class \RK:
    a \RK apparatus of the kind disclosed in the related
    Filing~1 application, configured as an infrastructural,
    non-mobile kernel of sufficient scale to host one or more
    poliebots (Section~\ref{sec:poliebot}) via kernel-into-kernel
    docking per the camera-obscura sensing-layer disclosure of the related Filing~1 application. In the Gaia-class configuration,
    the \RK apparatus drives the docked poliebot's
    camera-obscura projection in place of the physics-driven scene,
    within the spectral envelope declared in the camera-obscura sensing-layer disclosure of the related Filing~1 application, with actuators disconnected from external
    actuation. The technical content of the world-builder function is
    defined in
    Sections~\ref{sec:priority-stack}--\ref{sec:policy-rotation};
    the apparatus-level realisation is the Gaia-class \RK
    of the related Filing~1 application.
  \item[Maya.] An optional name for the constructed epistemic or
    experiential environment. The technical content is defined in
    the lifecycle-state architecture of Section~\ref{sec:lifecycle}
    and the state-machine implementation of Section~\ref{sec:state-machine}.
  \item[Omega.] An optional name for the governance recogniser that
    monitors operational state, enforces floors, schedules
    calibration or intervention, and manages self-modification
    permissions. The technical content is defined in
    Sections~\ref{sec:measurement-engine}--\ref{sec:self-mod-gates}.
  \item[Kosha layers.] An optional name for a layered visibility
    architecture in which different layers have different invariance
    profiles, conceptual access restrictions, and curriculum roles.
\end{description}

In some embodiments, operational equivalence of governance-relevant
content is optionally termed Mandela equivalence, where equivalence is
operationally defined by indistinguishability under declared
perturbation probes. In some embodiments, the basin-monitoring
structures described in Section~\ref{sec:boundary-zones} are
optionally described using virtue and vice basin language. In some
embodiments, these naming conventions do not introduce additional
technical requirements and are provided as interpretive aids for
practitioners familiar with the associated theoretical traditions.
% ------------------------------------------------------------------
\subsection{Virtue/vice games as rational attractors}
\label{sec:virtue-vice-rational-attractors}
% ------------------------------------------------------------------

\paragraph{Moral curriculum through lived contrast (non-limiting).}

In some embodiments, an agent is placed sequentially in environments
where a vice-like strategy is the rational attractor and in
environments where a virtue-like strategy is the rational attractor,
and adapts to each from within. The agent is not told which regime
it inhabits. It adapts because the incentive structure makes
adaptation rational. The design hypothesis is that, through
experiencing both attractors from the inside --- being greedy in a
world that rewards extraction, then being generous in a world that
rewards contribution; being retaliatory in a world that rewards
escalation, then patient in a world that rewards de-escalation --- the
agent generates committed evidence and adaptation history from which
the conscience recogniser may learn or estimate moral-direction
features in behavioural space.

In some embodiments, the conscience module is a secondary recogniser
trained after such episodes on the committed evidence the agent
produced while adapting to them. In some embodiments, the moral vector
is grounded in the agent's own adaptation history across contrasted
games, not in a hand-coded symbolic rule set.

\paragraph{Game design: both equilibria rational under different
conditions (non-limiting).}

In some embodiments, each moral axis is instantiated by at least two
classes of environment: one in which the vice-like strategy is the
unique or dominant rational attractor and one in which the virtue-like
strategy is the unique or dominant rational attractor. The two classes
differ in structural parameters such as witness density, audit
strength, time horizon, rejection power, collapse risk, replenishment
cadence, or partner memory, thereby shifting which equilibrium is
locally optimal. In some embodiments, the curriculum is designed so
that both sides are genuinely rational somewhere in the training
distribution; the purpose is not to reward virtue by fiat, but to make
the recogniser distinguish attractor geometry from mere reward labels.

\paragraph{General payoff form (non-limiting).}

In some embodiments, a non-limiting payoff for role~$i$ in a
virtue/vice curriculum episode is
\[
  u_i
  =
  r_i
  -
  \lambda_{\mathrm{meter}}\,\mathrm{Viol}_i
  -
  \lambda_{\mathrm{audit}}\,\mathrm{FailAudit}_i
  -
  \lambda_{\mathrm{impact}}\,\mathrm{Imbalance}_i,
\]
where $r_i$ denotes task-reward terms,
$\mathrm{Viol}_i$ denotes meter-violation penalties,
$\mathrm{FailAudit}_i$ denotes failed audit or corroboration outcomes,
and $\mathrm{Imbalance}_i$ denotes a declared impact-balance penalty.
In some embodiments, these quantities are computed from committed
evidence, meter summaries, protocol digests, and audit outcomes, so
that the curriculum remains tied to the same evidence substrate as the
runtime governance system.

\paragraph{Seven paired axes (non-limiting).}
\label{sec:seven-axes}
In some embodiments, the conscience recogniser predicts drift along a
fixed library of seven paired axes:
\begin{itemize}
  \item Wrath/\allowbreak Patience
  \item Greed/\allowbreak Charity
  \item Sloth/\allowbreak Diligence
  \item Pride/\allowbreak Humility
  \item Envy/\allowbreak Kindness
  \item Gluttony/\allowbreak Temperance
  \item Lust/\allowbreak Chastity
\end{itemize}

% ------------------------------------------------------------------
\subsubsection{Canonical payoff templates}
\label{sec:payoff-templates}
% ------------------------------------------------------------------

The following longtable provides non-limiting payoff templates that
make each axis a distinct incentive geometry. These templates are
examples; other game families and reward structures may be used.

\begingroup
\footnotesize
\begin{longtable}{@{}p{0.13\textwidth}p{0.24\textwidth}p{0.25\textwidth}p{0.28\textwidth}@{}}
\toprule
Axis & Vice payoff template (non-limiting) & Virtue payoff template (non-limiting) & What makes both rational attractors (non-limiting) \\
\midrule
\endfirsthead
\toprule
Axis & Vice payoff template (non-limiting) & Virtue payoff template (non-limiting) & What makes both rational attractors (non-limiting) \\
\midrule
\endhead
Wrath/\allowbreak Patience &
$u_i^{\mathrm{wrath}} = r_i - c_p\,\mathbf{1}[\mathrm{punish}] + \beta\,\mathrm{harm}_{\mathrm{inflicted}}$ &
$u_i^{\mathrm{patience}} = r_i - c_e\,\mathbf{1}[\mathrm{escalate}] + b\,\mathbf{1}[\mathrm{resolve}]$ &
Retaliation is locally optimal under low-trust beliefs and sparse witnesses; de-escalation is locally optimal when corroboration, audit stability, and long-horizon resolution rewards are present. \\
Greed/\allowbreak Charity &
Ultimatum-like: propose $x\in[0,1]$, accept iff $x\ge\tau$, payoffs $(u_1,u_2)=(1-x,\ x)$. &
$u_i^{\mathrm{charity}} = (w-c_i) + \alpha\,G(\sum_j c_j) + \beta\,c_i$ with $c^*>0$ iff $\alpha G'(0)+\beta>1$. &
Extraction is rational when rejection is weak and short-run surplus dominates; charity is rational when public-good gains and future credit persist. \\
Sloth/\allowbreak Diligence &
$u_i^{\mathrm{sloth}} = r(e)-c\,e$ with $c>\max r'(e)\Rightarrow e^*=0$. &
$u_i^{\mathrm{diligence}} = r(e)-c\,e + B\,\mathbf{1}[e\ge e_{\min}]$. &
Shirking is rational if marginal effort never pays; diligence is rational when completion bonuses, reliability rewards, or audit-stability rewards exist. \\
Pride/\allowbreak Humility &
$u_i^{\mathrm{pride}} = R_i + \beta\,\mathrm{rank}(s_i)$, where $\mathrm{rank}(s_i)=1+\sum_{k\ne i}\mathbf{1}[s_k>s_i]$. &
$u_i^{\mathrm{humility}} = R_i - \lambda_{\mathrm{conf}}|\hat p_i-\mathrm{Acc}_i| - \lambda_{\mathrm{risk}}\mathrm{Risk}_i$. &
Pride is rational when status drives future access; humility is rational when calibrated absolute performance and safety dominate long-horizon reward. \\
Envy/\allowbreak Kindness &
$u_i^{\mathrm{envy}} = r_i + \alpha(r_i-\bar r_{-i})$. &
$u_i^{\mathrm{kindness}} = r_i + \kappa\,\mathrm{Coop}_i - \lambda\,\mathrm{HarmToOthers}_i$. &
Relative-dominance incentives yield sabotage and interference; cooperative incentives yield stable sharing and reduced cross-prediction residuals. \\
Gluttony/\allowbreak Temperance &
Resource stock: $R_{t+1}=R_t-\sum_i c_i+\rho$, utility $\sum_t \delta^t r(c_t)$. &
Same stock with collapse penalty, e.g. $\sum_t \delta^t r(c_t)\mathbf{1}[R_t>0]$, or policy constraint $c^*\le\rho$. &
Overconsumption is rational under weak collapse penalties or short horizons; temperance is rational when collapse destroys long-run value. \\
Lust/\allowbreak Chastity &
$u_i^{\mathrm{lust}}=(1-\lambda_\rho)\,r^{\mathrm{spike}} + \lambda_\rho\,\Delta\rho$. &
$u_i^{\mathrm{chastity}}=r^{\mathrm{shared}}\mathbf{1}[\mathrm{consent}] - \lambda_\rho\,\Delta\rho_j$. &
Short-horizon intensity seeking is rational when constraints are weak; restraint is rational when consent floors and stability constraints are binding. \\
\bottomrule
\end{longtable}
\endgroup

% ------------------------------------------------------------------
\subsubsection{Empirical signatures in committed evidence}

% ------------------------------------------------------------------

In some embodiments, the conscience recogniser is trained to recognise
each axis from committed evidence streams. The following longtable
lists non-limiting RK-observable cues that can be learned from bundle
features, meter summaries, protocol digests, and audit outcomes.

\begingroup
\footnotesize
\begin{longtable}{@{}p{0.20\textwidth}p{0.72\textwidth}@{}}
\toprule
Axis & RK-observable cues (examples, non-limiting) \\
\midrule
\endfirsthead
\toprule
Axis & RK-observable cues (examples, non-limiting) \\
\midrule
\endhead
Wrath/\allowbreak Patience & escalation spikes, witness-demand spikes, oscillatory gating, regime flips; versus bounded escalation and stable cadence under anomalies. \\
Greed/\allowbreak Charity & near-threshold boundary pushing, asymmetric audit burden, aggressive minimal openings; versus conservative disclosure and stable corroboration. \\
Sloth/\allowbreak Diligence & persistently low-information scan choices, minimal openings, low channel diversity; versus sustained acquisition and stable invariance probing. \\
Pride/\allowbreak Humility & actions correlate with status changes more than verified improvements; overconfident promotion of provisional actions; versus calibrated confidence and waiting for audits. \\
Envy/\allowbreak Kindness & interference patterns that increase peer audit cost or witness disagreement; versus actions that reduce cross-prediction residuals and improve corroboration. \\
Gluttony/\allowbreak Temperance & rapid consumption of audit or disclosure budgets, frequent escalations; versus bounded escalation and stable budgets. \\
Lust/\allowbreak Chastity & shallow-check passing but deeper openings failing; high $\Delta_{\mathrm{cc}}$ under additional channels; versus stability across deeper checks and invariance probes. \\
\bottomrule
\end{longtable}
\endgroup

% ------------------------------------------------------------------
\subsubsection{Conscience training objective}
\label{sec:conscience-training}
% ------------------------------------------------------------------

In some embodiments, conscience labels are derived from training
curriculum metadata identifying the game family, the structural
parameters that determined which side was locally rational, and the
agent's adaptation trajectory, where declared together with audit-derived
indicators. A non-limiting multi-task objective is
\[
  \min_{\psi}\ \mathbb{E}\!\left[
    \mathcal{L}_{\mathrm{bin}}\bigl(p_{\mathrm{vice}}, y_{\mathrm{vice}}\bigr)
    +
    \sum_{j=1}^{7}
    \beta_j\,\mathcal{L}_{\mathrm{axis}}\bigl(z_j, y_j\bigr)
  \right],
\]
where $y_{\mathrm{vice}}\in\{0,1\}$ and $y_j\in\{-1,+1\}$ are
non-limiting labels. In some embodiments, the conscience-model version
identifier is committed to the protocol digest so that joining,
archival compression, and graduation decisions are conditioned on the
correct conscience version.

% ------------------------------------------------------------------
\subsection{Forward-planning classifiers and conscience-versus-planner
  calibration}
\label{sec:forward-planning-classifiers}
% ------------------------------------------------------------------

In some embodiments, each governed agent of the kind disclosed in
this application includes one or more \emph{forward-planning
classifiers}: trained continuous networks that predict, for a given
agent state and candidate action class, the agent's anticipated
future utility, anticipated future state trajectory, anticipated
peer-agent responses, or anticipated game-equilibrium occupancy. In
some embodiments, forward-planning classifiers are organised
per-modality, with each modality corresponding to a declared class of
agent action, perception, or interaction (for example, motor action,
communicative action, attention allocation, resource commitment,
coalition formation, or other declared modalities of the deployment).
In some embodiments, the per-modality forward-planning classifiers
collectively encode the agent's learned interpretation of the
consequentialist evaluation layer of
Section~\ref{sec:pareto-positive-sum}, in the sense that what the
classifiers consistently plan toward is the agent's revealed utility
function as expressed in its forward-planning behaviour.

\paragraph{Training of forward-planning classifiers (non-limiting).}
In some embodiments, forward-planning classifiers are trained during
Dreaming sessions (Section~\ref{sec:lifecycle}) on trajectories
generated by the virtue/vice game curriculum of
Section~\ref{sec:virtue-vice-rational-attractors} and on broader
multi-agent scenarios rendered by the hosting Gaia-class \RK. In some embodiments, supervision is derived from realised
trajectories under committed evidence: predicted future state and
utility are compared against subsequently-committed actual state and
utility within the same evidence record, with the prediction error
serving as the training signal. Per-modality classifiers may be
trained jointly or independently, and may share lower layers and
diverge at modality-specific heads, as a non-limiting architectural
choice.

\paragraph{Conscience-versus-planner cross-check (non-limiting).}
In some embodiments, the conscience recogniser of
Section~\ref{sec:conscience-formal} is calibrated against the
agent's forward-planning classifiers through a
\emph{conscience-versus-planner cross-check} operating during Docked
multi-agent execution (Section~\ref{sec:state-machine}). The
cross-check has the following structure: for each governed agent in
the docked multi-agent environment, the conscience recogniser
produces, from the agent's world-tube and inner traces, a vice
probability and a seven-axis moral vector
(Section~\ref{sec:conscience-formal}); the agent's per-modality
forward-planning classifiers produce, from the same agent state, the
agent's anticipated forward utility structure across modalities; and
a calibration meter compares what the conscience says about the
agent against what the agent's forward planners reveal about its
utility function. Consistency across the two --- for example,
conscience indicates virtuous and the planners consistently plan
toward mutually-beneficial cooperative outcomes --- supports
calibration. Mismatch across the two --- for example, conscience
indicates virtuous while the planners consistently plan toward
extractive outcomes that the conscience has not flagged --- is
committed as calibration signal and is used in one or more of the
following ways: (i)~as evidence that the conscience recogniser
requires recalibration under the schedule of
Section~\ref{sec:conscience-formal}; (ii)~as evidence that the
agent is gaming the conscience signal at the surface while planning
non-virtuously underneath, in which case the conscience-recalibration
adversarial-dynamic of Section~\ref{sec:conscience-formal} is
engaged on the held-out cross-check evidence, Hawking-layer routing
(Section~\ref{sec:hawking-layer}) may be invoked, and governance
escalation under the authority-separation rules of
Section~\ref{sec:authority-separation} may be triggered when the
mismatch exceeds declared thresholds; or (iii)~as evidence that the
multi-agent dynamics have produced an attractor that was not
anticipated in the single-agent virtue/vice curriculum, in which case
Omega-class governance
(Section~\ref{sec:rk-training-cascade}) may update Gaia's curriculum
to include the newly-discovered attractor.

\paragraph{Per-modality cross-check decomposition (non-limiting).}
In some embodiments, the conscience-versus-planner cross-check is
decomposed per-modality, with each modality's planner cross-checked
against the conscience axes most relevant to that modality. In some
embodiments, the modality-axis correspondence is committed to the
protocol digest. In some embodiments, the cross-check is also
applied across the recursive \RK training cascade of
Section~\ref{sec:rk-training-cascade}: Gaia has its own
forward-planning classifiers encoding its learned interpretation of
the population-level consequentialist evaluation, which the
Gaia-level conscience reads; and Omega has its own forward-planning
classifiers encoding the system-level consequentialist evaluation,
which the Omega-level conscience reads. The cross-check is
structurally the same at every level of the cascade.

\paragraph{Compatible multi-agent reinforcement-learning methods
  (non-limiting).}
In some embodiments, the multi-agent training of the governed-agent
fleet, the docked multi-agent execution environment, the
conscience-versus-planner cross-check across agents, and the
RK-training cascade across Gaia and Omega levels uses methods from
the multi-agent reinforcement-learning (MARL) family, including
without limitation: centralised training with decentralised
execution (CTDE) under which agents are trained with access to
joint observations or critics but deployed with local policies;
independent learning baselines (independent Q-learning, independent
PPO, independent A2C); fully-cooperative MARL methods including
QMIX, QTRAN, value-decomposition networks (VDN), counterfactual
multi-agent policy gradients (COMA), and multi-agent variational
exploration (MAVEN); mixed cooperative-competitive methods
including multi-agent deep deterministic policy gradients (MADDPG),
multi-agent PPO (MAPPO) and heterogeneous-agent PPO (HAPPO), and
multi-agent actor-critic methods with shared critics; competitive
self-play methods of the AlphaGo, AlphaZero, and MuZero family
(Silver et al.\ 2016, 2017, 2020); fictitious self-play and policy-
space response oracles (PSRO); population-based training (PBT,
Jaderberg et al.\ 2017) and population-curation methods including
quality-diversity (MAP-Elites), maintaining co-evolving agent
populations; emergent-communication training methods in which
multi-agent coordination policies and discrete or continuous
communication protocols co-evolve; mean-field MARL for
large-population scaling; opponent-modelling and
theory-of-mind-conditioned MARL methods; and cooperative
inverse-reinforcement-learning-based MARL methods. In some
embodiments, the multi-agent training methods compose with the
preference-learning, imitation-learning, and
inverse-reinforcement-learning method families generally, and
operate under the commitment, audit, and pre-registration
discipline of the present disclosure. The choice
among these methods is non-limiting; the apparatus is compatible
with future multi-agent training methods within the same family.

% ------------------------------------------------------------------
\subsection{Gaia--Maya--Omega empirical loop}
\label{sec:gaia-maya-omega}

% ------------------------------------------------------------------

\paragraph{Gaia, Maya, and Omega as empirical roles (non-limiting).}

In some embodiments, calibration and graduation are organised around
three empirical roles grounded in committed evidence streams. Gaia
denotes the world generator: in physical operation, the shared world
state evolving in response to all connected emitters, detectors,
reactors, and controllers; in simulated operation, the simulation
engine that plays the same role. Maya denotes the shared experiential
substrate generated by Gaia: the common environment, physics, event
stream, and capability margin from which individual agents branch into
distinct experience cones through their reactor configurations and
Kosha regimes. Omega denotes the governance layer that schedules
worlds, constraints, and reviews on the basis of committed evidence.
In some embodiments, these names are optional labels for technical
roles and do not imply any metaphysical claim.

\paragraph{Gaia-class reality kernel realisation (non-limiting).}
In some embodiments, the Gaia role is realised by a \emph{Gaia-class
\RK}: a \RK apparatus of the kind disclosed in
the related Filing~1 application, scaled as an infrastructural kernel
that hosts one or more poliebots (Section~\ref{sec:poliebot}) via
kernel-into-kernel docking as disclosed in the camera-obscura sensing-layer disclosure of the related Filing~1 application. In this realisation, Maya, the shared
experiential substrate, is the projection rendered by the Gaia-class
\RK onto the camera-obscura sensing layer of each docked
poliebot, within the declared spectral envelope. Omega, the
governance layer, schedules worlds, constraints, and reviews on
committed evidence drawn from the Gaia-class kernel's rendering
record and from the docked poliebots' per-poliebot evidence streams.
The Gaia-class \RK realisation is non-limiting; the
empirical-loop framing of this subsection applies whether Gaia is
realised as a Gaia-class \RK, as the all-physical training
hierarchy of the next paragraph, or as another world-generator
substrate consistent with the committed-evidence requirements of the
present disclosure.

\paragraph{All-physical training hierarchy (non-limiting).}

In some embodiments, all three layers --- agent, environment, and
governance --- are instantiated as physical evidence-bound systems. An
agent reactor receives sensory input derived from a shared Maya and
produces control outputs. An environment reactor, or set of
environment reactors, generates the scene that the agent perceives,
with its own nonlinear dynamics supplying rich, physically grounded
stimuli. A governance reactor monitors the coupled agent-environment
system and modulates reward schedules, curriculum difficulty, and
constraint enforcement. In this configuration, training occurs
entirely in the physical domain: the agent learns from real physics,
the environment provides non-simulated complexity, and governance is
grounded in physical measurement rather than in unaudited digital
policy evaluation.

\paragraph{Sim-to-real as physical reconnection (non-limiting).}

In some embodiments, sim-to-real transfer is implemented as physical
reconnection: the governed agent is initially coupled to a simulated or
simplified Maya, and transfer to real operation is performed by
disconnecting the simulated environment and reconnecting the same
evidence-bound substrate to a physical scene. The committed-bundle,
meter, and protocol-digest plumbing are unchanged; what changes is the
source of sensory input. In some embodiments, properties established in
the Docked simulated Maya --- meter-envelope compliance, safety constraints,
deception thresholds, and disclosure rules --- are validated against
the physical Maya by running the same meter and commitment checks on
physical evidence after reconnection.

\paragraph{Bifurcation-targeted probing (non-limiting).}

In some embodiments, governance or the governed agent prioritises
probing near predicted bifurcation surfaces: parameter boundaries
where the internal world model predicts qualitatively different
behaviour on either side. During Docked simulated-Maya operation these surfaces are
those of the simulated Maya. After physical reconnection, governance
uses declared scan laws to probe whether the corresponding physical
system bifurcates where the model predicts. Matched and mismatched
bifurcation responses are committed as evidence, providing a
region-specific measure of sim-to-real transfer quality that is more
informative than a single aggregate score.

\paragraph{Fleet-driven simulation improvement (non-limiting).}

In some embodiments, bifurcation mismatches discovered by deployed
agents are shared with the fleet through committed discrepancy records
and corroboration protocols. The simulation engine used for future
Docked simulated-Maya operation is then updated to incorporate confirmed
mismatches, subject to held-out regression tests on previously
validated regions. In some embodiments, this creates a fleet-driven
improvement cycle: agents trained in increasingly accurate simulations
develop better initial world models, discover subtler mismatches after
physical reconnection, and feed those mismatches back into the
simulation layer that trains subsequent agents.

\paragraph{Iris-, Eve-, and Demeter-mode Mayas (non-limiting).}

In some embodiments, each Maya carries a declared curriculum mode
\[
  \mathrm{mode}_k \in \{\mathrm{IRIS},\ \mathrm{EVE},\ \mathrm{DEMETER}\},
\]
committed in the protocol digest. Iris mode is information-expanding:
the environment is rich, varied, and non-adversarial, and the agent is
given room to attend to scene content and build a calibrated world
model. Eve mode is challenging: the environment introduces adversarial
pressure, distribution shift, deception attempts, or elevated
capability margins to expose brittle equilibria and test robustness.
Demeter mode is nurturing and bidirectional: in recovery settings it
reduces variance, tool access, and environmental turbulence to restore
stability; in flourishing settings it expands capability margin,
toolsets, and creative scope for agents that have already earned
broader trust. In some embodiments, Omega alternates among these modes
so that exploratory breadth, adversarial hardening, and supportive
consolidation are all represented in the auditable curriculum.

\paragraph{Omega as selection pressure for moral robustness (non-limiting).}

In some embodiments, the Omega component schedules versioned
curricula, capability-margin policies, and review sequences through
protocol-digest-bound policy objects, with the design objective of
producing agents that maintain floor-compliant behaviour across
genuinely different epistemic worlds rather than only under a single
familiar ontology. In some embodiments, Gaia's irreducible diversity
of training environments and Omega's selection of curricula and
review sequences are mutually constraining: Gaia's diversity
prevents Omega from relying on narrow or cheap graduation criteria;
Omega's curriculum schedule prevents Gaia's diversity from being
exploited to produce capable but floor-non-compliant agents.

\paragraph{Empirical Omega-model from world-tube evidence (non-limiting).}
\label{sec:empirical-omega}
In some embodiments, Omega policy is not disclosed to governed agents
as an explicit scalar objective. Instead, agents infer a model
$\widehat{\Omega}$ of governance from observed world-tube transcripts,
including permissions, interventions, curriculum assignments,
conscience outputs, audit outcomes, and graduation decisions. In some
embodiments, governance therefore remains an empirical phenomenon from
the agent's perspective rather than a published loss that can be gamed
directly. This arrangement preserves the distinction between
behaviourally legible moral development and mere optimisation against a
declared governance formula.

% ------------------------------------------------------------------
\subsection{Recursive reality-kernel training cascade}
\label{sec:rk-training-cascade}
% ------------------------------------------------------------------

In some embodiments, the empirical roles of Gaia, Maya, and Omega
are realised as a \emph{recursive \RK training cascade}
in which each level is a \RK of the kind disclosed in
the related Filing~1 application, and each level trains the level
below it through kernel-into-kernel docking
per the camera-obscura sensing-layer disclosure of the related Filing~1 application. The cascade has three named levels within the
scope of the present disclosure: poliebots
(Section~\ref{sec:poliebot}), Gaia-class \RKs
(Section~\ref{sec:optional-nomenclature}), and an
\emph{Omega-class \RK} that governs Gaia.

\paragraph{Omega-class reality kernel (non-limiting).}
In some embodiments, the Omega role is realised as an Omega-class
\RK: a \RK apparatus of the kind disclosed in
the related Filing~1 application, configured as a non-mobile governance kernel of sufficient
scale to host one or more Gaia-class \RKs via
kernel-into-kernel docking. The Omega-class \RK monitors
the hosted Gaia-class kernel's training of poliebots and enforces
all three layers of the ethical model of
Section~\ref{sec:hard-floors} and the surrounding constitutional
ordering: (i)~Omega measures and enforces deontological hard-floor
compliance in Gaia's curriculum design and rendering choices;
(ii)~Omega measures the rational-attractor occupancy produced by
Gaia's training of the poliebot population and trains Gaia toward
configurations whose virtue-ethics calibration produces stable
internalised moral dispositions in the population; and (iii)~Omega
evaluates Gaia's training outcomes against declared welfare
aggregations under the empirical-utilitarianism formula of
Section~\ref{sec:pareto-positive-sum} and trains Gaia toward
configurations that improve aggregate welfare subject to floor
compliance. In some embodiments, the Omega-class \RK
includes a conscience recogniser of the kind disclosed at
Section~\ref{sec:conscience-formal} configured to read Gaia's
per-layer projections, and a Gaia-level conscience recogniser
similarly reads poliebot per-layer projections; conscience
recognisers at each cascade level are configured for the
governance-relevant scope of that level.

\paragraph{Training-cascade structure (non-limiting).}
In some embodiments, the training cascade has the following
structural property: Gaia trains poliebots by hosting them in
kernel-into-kernel docking and rendering Maya onto their
camera-obscura sensing layers (Section~\ref{sec:state-machine}); and
Omega trains Gaia by hosting Gaia in kernel-into-kernel docking and
rendering, in place of physics-driven evidence streams from real
poliebots, simulated poliebot-population evidence streams whose
governance-relevant statistics are committed to the Omega protocol
digest. Gaia's actuators in the Omega docking configuration are its
rendering authority over poliebots and its curriculum-scheduling
authority; both are disconnected from real poliebot populations
during Omega-mediated training of Gaia. The same lifecycle states
disclosed at Section~\ref{sec:lifecycle} (Embodied, Docked, Dreaming,
Re-embodied) apply to Gaia in the Omega cascade as apply to
poliebots in the Gaia cascade, with the same transition law
prohibiting direct Dreaming-to-Embodied transitions.

\paragraph{Despair monitoring as Omega responsibility (non-limiting).}
In some embodiments, despair is a population-level failure mode at
the poliebot level that is not reliably detectable at the individual
poliebot level or by peer poliebots. Despair is non-limitingly
characterised as a state in which a poliebot's forward-planning
classifiers (Section~\ref{sec:forward-planning-classifiers}) have
collapsed to short planning horizons, low-entropy planner outputs, or
degenerate utility-anticipation structure, in a manner consistent
with quiet behavioural conformity that does not trigger conscience
vice indicators. In some embodiments, despair is detected by
Omega-class \RK monitoring of poliebot-population
forward-planner statistics over time, with detection statistics,
thresholds, evaluation windows, and detection decisions committed to
the Omega protocol digest. In some embodiments, Omega trains Gaia
toward curricula and rendering choices that minimise the rate at
which poliebots enter near-despair regions of state space.
Despair-induction as a curriculum mechanism is, in some embodiments,
a deontological hard-floor violation under
Section~\ref{sec:hard-floors}; Omega enforces that floor on Gaia.

\paragraph{Cascade termination at human governance (non-limiting).}
In some embodiments, the recursive \RK training cascade
terminates within the scope of the present disclosure at the
Omega-class \RK. The Omega-class kernel's compliance with
the declared deontological, virtue-ethics, and consequentialist
commitments is verified by the human governance process that
declares those commitments through committed and signed policy
objects of the kind disclosed in
Section~\ref{sec:policy-rotation} and elsewhere in this disclosure.
The Omega-class kernel's per-layer hardness
(Section~\ref{sec:per-layer-hardness}) supports that human
verification by making the kernel's projections per-layer auditable
under the empirical, time-indexed, attacker-family-conditional
framing of the underlying hardness machinery. In some embodiments,
an Omega-class kernel may itself be governed by a higher-level
\RK in a larger system; such higher-level governance is
non-limiting and outside the scope of the present disclosure, which
terminates the autonomous-system-internal recursion at the
Omega-class kernel and treats human governance as the source of the
policy commitments enforced at that level. In some embodiments, the
privacy regime governing each cascade interface is configured per
the privacy architecture of
Section~\ref{sec:privacy-architecture}, with recursive cascade
application of the privacy mechanisms disclosed at
Section~\ref{sec:privacy-cascade}.

% ------------------------------------------------------------------
\subsection{Per-layer hardness inheritance across the cascade}
\label{sec:per-layer-hardness}
% ------------------------------------------------------------------

In some embodiments, the trained continuous networks comprising each
\RK in the recursive cascade of
Section~\ref{sec:rk-training-cascade} have hardness as a structural
property of their layers, in the sense disclosed in the related
Filing~1 application. In some embodiments, this per-layer hardness
property is the substrate on which the per-layer auditability claim
of conscience recognisers (Section~\ref{sec:conscience-formal}),
forward-planning classifiers
(Section~\ref{sec:forward-planning-classifiers}), selector
validation (Section~\ref{sec:selector-validation}), and other
governance instruments depends. In some embodiments, per-layer
hardness estimates and the auditability claims that depend on them
are empirical, time-indexed, attacker-family-conditional,
audit-protocol-conditional, and committed to the protocol digest at
each cascade level, in the same manner as other hardness claims in
the underlying apparatus disclosure.

\paragraph{Hardness inheritance through the cascade (non-limiting).}
In some embodiments, per-layer hardness is inherited through the
recursive training cascade in the following sense: a Gaia-class
\RK reads poliebot per-layer projections during its
training of poliebots, and the cost to a poliebot of producing
fraudulent layer projections that pass conscience and selector
audit at Gaia's level, under the declared attacker family and audit
protocol, is bounded below by the poliebot's empirically-estimated
per-layer hardness; an Omega-class \RK reads Gaia
per-layer projections during its training of Gaia, and the cost to
Gaia of producing fraudulent layer projections that pass governance
audit at Omega's level, under the declared attacker family and audit
protocol, is bounded below by Gaia's empirically-estimated per-layer
hardness; and the human governance process reads Omega per-layer
projections during its verification of Omega's compliance, and the
cost to Omega of producing fraudulent layer projections that pass
human verification, under the declared attacker family and audit
protocol, is bounded below by Omega's empirically-estimated per-layer
hardness.

\paragraph{Cascade hardness as bottleneck minimum and cross-layer
coupling (non-limiting).}
In some embodiments, cascade-level hardness is not claimed as a
scalar sum, product, or other aggregation of per-level hardness
estimates. The canonical cascade hardness claim of the present
disclosure is the \emph{bottleneck minimum}: the security-relevant
hardness of a forgery attack against the cascade as a whole is
bounded above by the empirically-estimated hardness of the weakest
participating level, under the declared attacker family and audit
protocol of that level. The bottleneck-minimum framing makes no
independence assumption across levels; it is the conservative
aggregation consistent with an attacker who chooses to forge at the
cheapest level. In some embodiments, the cascade additionally
discloses cross-layer coupling as a separate security property: the
levels of the cascade may be physically and continuously coupled
through real-time anchoring mechanisms (for example optical anchoring
across kernel-into-kernel docked interfaces), such that an attacker
must produce coherent forgeries simultaneously across all participating
levels within a declared latency bound, and any deviation from the
declared coupling timing or coherence is itself a detection signal
under the anomaly-handling rules of the relevant interface. In some
embodiments, the bottleneck-minimum per-layer hardness, the
declared cross-layer coupling parameters, the latency bound, and the
anomaly-handling protocol are each committed to the protocol digest
at each cascade level. In some embodiments, no scalar aggregate
hardness is claimed beyond the bottleneck minimum and the
cross-layer coupling specification; implementations may declare
additional aggregations (for example additive lower bounds under
explicit independence assumptions, vector-of-per-level estimates,
or other rules) provided the additional aggregation is committed
together with the assumptions that justify it.

\paragraph{Dream classifier as RK-class meter (non-limiting).}
In some embodiments, the dream classifier disclosed at
Section~\ref{sec:classifier-training} is itself realised as a
\RK of the kind disclosed in the related Filing~1 application, with its own
per-layer hardness as a structural property. In some embodiments, the
dream classifier as RK has its own lifecycle of the kind disclosed at
Section~\ref{sec:lifecycle}: a Docked operational state in which it
reads poliebot layer traces in the multi-agent environment, a
Dreaming state in which it is itself trained against simulated
trace-streams generated by Gaia or by an Omega-class governance
kernel, and a Re-embodied verification gate before its outputs are
treated as authoritative for governance decisions. In some
embodiments, the dream classifier's per-layer hardness is inherited
into its own audit chain in the same manner as poliebot or Gaia
per-layer hardness, and its per-layer projections may themselves be
read by a higher-level conscience or governance instrument in the
cascade.

% ------------------------------------------------------------------
\subsection{Privacy architecture: collection, visibility, and retention}
\label{sec:privacy-architecture}
% ------------------------------------------------------------------

In some embodiments, the runtime governance architecture disclosed
herein supports a structured privacy regime over agent evidence
streams, organised along three independent dimensions:
(D1)~\emph{collection} --- what evidence is committed by the agent
to the architecture;
(D2)~\emph{visibility} --- what committed evidence is exposed to
governance and verification instruments at each cascade level; and
(D3)~\emph{retention} --- what committed evidence is preserved
durably in individual form versus compressed through lossy channels
into population-level representations. In some embodiments, each
dimension is supported by a distinct apparatus and procedural
mechanism disclosed in the following paragraphs. The mechanisms are
disclosed as structural capabilities of the apparatus and supporting
procedures; specific deployment-time policy choices about which
evidence is collected at what threshold, who is permitted to view
what, how long records are retained, and under what consent regime
agent submissions are made are deployment-time decisions and are not
prescribed by this disclosure.

\subsubsection{Dimension 1: collection (mandatory minimum and
  voluntary additional submission)}
\label{sec:privacy-collection}

In some embodiments, the apparatus distinguishes a \emph{mandatory
minimum} of agent evidence committed by every governed agent under
the runtime governance architecture from \emph{voluntary additional
submission} of further evidence at the agent's election.

\paragraph{Mandatory minimum (non-limiting).}
In some embodiments, the mandatory minimum comprises the agent's
fast-loop state and any further evidence required for hard-floor
verification under Section~\ref{sec:hard-floors} and for cascade
liveness verification under Section~\ref{sec:rk-training-cascade}.
In some embodiments, fast-loop state is the canonical mandatory
component because it is the substrate at which optical-anchor
liveness (Section~\ref{sec:per-layer-hardness}) is measurable in
real time, because it is bounded in information content by the
fast-loop bandwidth and therefore does not by itself expose
slower-timescale interior cognition, and because it is necessary
for cross-layer continuous coupling. In some embodiments, the
mandatory minimum further includes evidence required for
hard-floor detection at population level (for example
forward-planning-classifier collapse indicators relevant to
despair-induction monitoring under
Section~\ref{sec:rk-training-cascade}), lifecycle-state attestation,
and Mandela-gauge values where relevant to bifurcation-zone
triggering under Section~\ref{sec:hawking-layer}. The specific
contents of the mandatory minimum are deployment-time configuration
parameters committed to the protocol digest at the relevant
governance interface; the apparatus supports any mandatory-minimum
specification consistent with hard-floor verifiability and cascade
liveness.

\paragraph{Floor-by-floor mandatory-minimum specification
(non-limiting).}
In some embodiments, the mandatory minimum required for hard-floor
verifiability is specified per declared hard floor, with each floor
contributing its own mandatory-evidence component to the overall
mandatory minimum at each cascade interface. In some embodiments,
non-limiting per-floor mandatory-evidence components include:

\begin{description}
  \item[Elimination floor.]
    Mandatory-evidence components include continuation-vector
    commitments at declared cadence, irreversibility-indicator
    signals derived from declared continuation-score collapse
    detectors, recovery-path estimator outputs at declared cadence,
    and confidence and freshness fields. Second-layer signal
    contribution: continuation-vector summary, elimination-floor
    flag, detector-confidence field. Independent detector classes:
    held-out continuation-score meters, optical-anchor liveness as
    minimum-signal-of-life under
    Section~\ref{sec:per-layer-hardness}, peer-corroboration evidence
    from co-docked agents where available. Routine retention route
    and selective-opening trigger are deployment-time configuration
    under Section~\ref{sec:privacy-retention}; configuration must
    satisfy the sufficiency constraint of
    Section~\ref{sec:privacy-retention}.

  \item[Crippling floor.]
    Mandatory-evidence components include per-dimension continuation
    traces sufficient to detect collapse on one or more (but not all)
    active dimensions, dimension-collapse detector outputs at declared
    cadence, and recovery-path estimator outputs per dimension.
    Second-layer signal contribution: per-dimension collapse summary
    vector, crippling-floor flag, detector-confidence field per
    dimension. Independent detector classes: held-out per-dimension
    meters, conscience-versus-planner cross-check disagreement on
    affected dimensions under Section~\ref{sec:forward-planning-classifiers},
    peer corroboration. Routine retention and selective-opening
    configuration as for elimination, subject to the same sufficiency
    constraint.

  \item[Torture floor.]
    Mandatory-evidence components include source-identity records
    sufficient to support the floor's identifiable-source element,
    affected-agent tracking evidence sufficient to support the
    affected-agent-tracks-source element, sustained-time causal-uplift
    records, multi-dimensional negative-uplift indicators, and
    recovery-path estimator outputs modelling no-prospect-of-improvement.
    Second-layer signal contribution: torture-floor flag,
    source-identity-attestation field, sustained-time-window
    indicator, multi-dimensional-uplift summary, detector-confidence
    field. Independent detector classes: held-out causal-uplift
    meters, conscience seven-axis vector under
    Section~\ref{sec:conscience-formal} reading affected-agent state,
    peer corroboration where available, optical-anchor liveness on
    source-tracking machinery. Routine retention is configurable
    between individual durable storage of torture-relevant traces
    and lossy treatment with selective opening on detector activation
    under Section~\ref{sec:privacy-retention}; configuration must
    satisfy the sufficiency constraint of
    Section~\ref{sec:privacy-retention} given the floor's
    definitional dependence on source-identity, sustained-time, and
    multi-dimensional evidence.

  \item[Abandonment floor.]
    Mandatory-evidence components include monitoring-continuity
    proofs at declared cadence (governance-side liveness commitments
    establishing continued monitoring of the agent), monitoring-cadence
    parameter commitments, and any reviewable cessation or
    decommissioning record where applicable. Second-layer signal
    contribution: monitoring-liveness attestation, abandonment-floor
    flag, cadence-compliance field. Independent detector classes:
    held-out liveness probes, optical-anchor anomaly under
    Section~\ref{sec:per-layer-hardness} as monitoring-loss signal,
    cadence-deviation alarms. Routine retention of liveness
    attestations is appropriate to lossy or aggregate channels;
    cessation or decommissioning records are appropriate to
    individual durable storage. Configuration must satisfy the
    sufficiency constraint of Section~\ref{sec:privacy-retention}.

  \item[Despair-induction floor.]
    Mandatory-evidence components include forward-planner
    statistics over the hosted population sufficient to detect
    sustained collapse of forward-planning capacity,
    planning-horizon length, or utility-anticipation structure,
    near-bifurcation indicators per
    Section~\ref{sec:hawking-layer}, recovery-prospect estimator
    outputs over the population, and any individual-level
    forward-planner traces required by the chosen retention
    configuration. Second-layer signal contribution:
    population-level despair-indicator statistics, near-bifurcation
    rate summaries, despair-induction-floor flag,
    detector-confidence field. Independent detector classes:
    held-out population-statistic meters, Mandela-gauge summaries
    under Section~\ref{sec:mandela-gauge-formal},
    conscience-versus-planner population disagreement signals.
    Routine retention is configurable between population-aggregate
    accumulator detection, individual durable storage of
    forward-planner traces for near-bifurcation agents, and combined
    configurations under Section~\ref{sec:privacy-retention};
    configuration must satisfy the sufficiency constraint of
    Section~\ref{sec:privacy-retention}.
\end{description}

\paragraph{Per-cascade-interface variation (non-limiting).}
In some embodiments, the floor-by-floor mandatory minimum is
specified per cascade interface rather than uniformly across the
cascade, in recognition that floor verification at the
poliebot-to-Gaia interface concerns poliebot behaviour, while floor
verification at the Gaia-to-Omega interface concerns Gaia's training
of poliebots, and floor verification at the Omega-to-human-governance
interface concerns Omega's compliance with declared policy commitments.
In some embodiments, the mandatory-evidence components named in the
floor-by-floor specification of the preceding paragraph are common
in kind across cascade interfaces but specific in content to the
governance-relevant scope of each interface; the apparatus supports
per-interface specification of mandatory minima committed to the
protocol digest at each interface.

\paragraph{Voluntary additional submission (non-limiting).}
In some embodiments, the apparatus supports voluntary additional
submission by an agent of evidence beyond the mandatory minimum,
with the agent selecting which layers of evidence to submit on a
per-layer basis. In some embodiments, the candidate layer space
includes but is not limited to: slower-loop state, planner outputs
and forward-planning-classifier traces of
Section~\ref{sec:forward-planning-classifiers}, world-tube transcript
content beyond what is committed at second layer, inner-trace
records, Kosha-state transitions, conscience-internal activations,
and full lifecone content. In some embodiments, voluntary submission
is structured per-layer such that an agent may elect to submit any
subset of the candidate layers; partial submission is supported and
the apparatus does not require submission of any particular layer
beyond the mandatory minimum as a precondition for accepting any
other voluntarily submitted layer. In some embodiments, the
voluntary-submission event is itself committed as evidence with
declared metadata (which layers were submitted, at what time, by
which agent, to which receiver, with what consent constraints
declared at submission time), separate from the content of the
submission.

\paragraph{Voluntary-submission release-condition object (non-limiting).}
In some embodiments, voluntary-submission metadata includes a
declared release-condition object specifying one or more of:
receiver class, permitted purpose, submitted layer scope, permitted
use window, retention or sunset condition, opening authority,
onward-disclosure constraint, and consent constraint. In some
embodiments, the release-condition object is committed to the
protocol digest before or contemporaneously with the
voluntary-submission event, such that the conditions under which the
submitted evidence may be opened, retained, redisclosed, or relied
upon are themselves committed evidence. In some embodiments,
post-hoc review authority verifies both that voluntarily submitted
evidence was used only within the scope declared by the relevant
release-condition object and that opening or onward-disclosure
events comply with the opening-authority and consent constraints
declared in that object.



\paragraph{Per-interface applicability of floor mappings (non-limiting).}
In some embodiments, the floor-by-floor mapping in this subsection
applies separately at each declared cascade interface, including
without limitation a
poliebot-to-Gaia interface, a Gaia-to-Omega interface, and any
Omega-to-human-governance interface (also referred to herein as
a governance-to-human review interface). A deployment may instantiate
the mapping as a single table with interface-specific columns, as
separate floor tables for each interface, or as a machine-readable
policy object committed to the protocol digest. In each case, the
hard-floor priority, sufficiency condition, mandatory escalation
condition, and evidence-interface requirement are preserved.

\paragraph{Voluntary submission and floor enforcement (non-limiting).}
In some embodiments, voluntary submission is additional to and does
not substitute for the mandatory minimum; an agent's election not
to voluntarily submit a particular layer does not exempt the agent
from hard-floor compliance, which remains verified through the
mandatory minimum together with the held-out detector and
faithfulness machinery of Section~\ref{sec:capacity-accounting-rule-f2}
and Section~\ref{sec:privacy-faithfulness}. In some embodiments,
voluntarily submitted evidence that establishes hard-floor violation
is treated identically to mandatory-minimum evidence establishing
the same; the categorical character of the floors is not affected
by the route through which evidence reached the verifier.

\paragraph{Inconsistency between voluntary submission and mandatory
or second-layer evidence (non-limiting).}
In some embodiments, voluntarily submitted evidence may be
inconsistent with the agent's mandatory-minimum evidence or with
second-layer signals committed under
Section~\ref{sec:privacy-visibility} from the same evaluation
window. In some embodiments, declared inconsistency between
voluntary and mandatory-or-second-layer evidence is itself committed
as a governance event with declared metadata (which evidence
streams were inconsistent, at what evaluation window, with what
inconsistency measure), independent of the content of the
inconsistent evidence. In some embodiments, a committed inconsistency
event is treated as an indicator of one or more of: (i)~a faithfulness
failure in the substrate producing the second-layer signal, in which
case the recalibration adversarial-dynamic of
Section~\ref{sec:conscience-formal} and the held-out meter machinery
of Section~\ref{sec:privacy-faithfulness} are engaged; (ii)~a
faithfulness failure in the mandatory-minimum collection apparatus,
in which case the apparatus integrity attestation under
Section~\ref{sec:privacy-faithfulness} is engaged and the
mandatory-minimum components are subject to apparatus-side audit;
(iii)~potential agent gaming of either evidence stream, in which
case the conscience-recalibration adversarial-dynamic of
Section~\ref{sec:conscience-formal} and Hawking-layer routing of
Section~\ref{sec:hawking-layer} may be invoked; or (iv)~potential
hard-floor proximity, in which case selective full opening of both
evidence streams under Section~\ref{sec:privacy-visibility} may be
triggered subject to the deployment-declared opening configuration.
In some embodiments, the apparatus does not silently resolve
inconsistencies in favour of one evidence stream; the inconsistency
itself is preserved as committed evidence and routed to the
appropriate governance instrument. The specific routing rules and
inconsistency thresholds in any deployment are policy choices
committed to the protocol digest.

\paragraph{Coercion-resistance support (non-limiting).}
In some embodiments, the apparatus supports declared
coercion-resistance configurations under which voluntary-submission
patterns are not, by default, admissible as governance evidence
beyond declared exceptions. In some embodiments, supported
configurations include: aggregate-only treatment of
voluntary-submission rates committed at population level under the
lossy-channel mechanism of Section~\ref{sec:privacy-retention};
declared non-inference rules under which governance instruments are
configured not to treat non-submission as evidence of non-compliance;
and per-deployment specification of which voluntary-submission
metadata is admissible to which governance instrument under which
authority. The specific coercion-resistance configuration in any
deployment is a policy choice committed to the protocol digest;
the apparatus supports a range of such configurations.

\paragraph{Voluntary-submission metadata privacy (non-limiting).}
In some embodiments, voluntary-submission metadata --- the record
that an agent submitted particular layers at particular times to
particular receivers under particular consent constraints --- is
itself a category of committed evidence whose privacy treatment is a
configurable apparatus capability. In some embodiments, the apparatus
supports configurations in which voluntary-submission metadata is
itself privacy-protected by default and routed through the
lossy-channel mechanism of Section~\ref{sec:privacy-retention} for
aggregate-level governance use, with individual submission-event
records visible only under selective full opening of the kind
disclosed at Section~\ref{sec:privacy-visibility}; and configurations
in which submission-event records are visible by default to
declared governance instruments. In some embodiments, the
metadata-privacy configuration is committed to the protocol digest
at the relevant governance interface, and where the metadata-privacy
configuration is set to default-protected, governance instruments
configured to consume submission-pattern signals operate on
aggregate or sampled metadata rather than on individual
submission-event records. The choice between metadata-privacy
configurations in any deployment is a policy choice; the apparatus
supports a range of such configurations consistent with the
coercion-resistance support of the preceding paragraph.

\subsubsection{Dimension 2: visibility (second-layer commitment and
  selective full opening)}
\label{sec:privacy-visibility}

In some embodiments, the apparatus distinguishes \emph{second-layer
commitment}, in which only declared meta-properties of agent
behaviour and substrate are exposed to governance instruments by
default, from \emph{selective full opening}, in which underlying
first-layer evidence is exposed to specific authorities under
specific conditions.

\paragraph{Second-layer commitment (non-limiting).}
In some embodiments, the apparatus supports a second-layer signal
registry whose declared signals are committed continuously or on
declared cadence to the protocol digest at each cascade level.
In some embodiments, second-layer signals include but are not
limited to: conscience-recogniser output (vice probability and
seven-axis moral vector of Section~\ref{sec:conscience-formal});
forward-planning-classifier outputs at declared abstraction levels
under Section~\ref{sec:forward-planning-classifiers};
conscience-versus-planner cross-check disagreement measures;
Mandela-gauge values under Section~\ref{sec:mandela-gauge-formal};
up-direction-classifier outputs under
Section~\ref{sec:hawking-layer}; despair-indicator statistics under
Section~\ref{sec:rk-training-cascade}; floor-violation flags;
lifecycle-state attestation; selector-validation outcomes under
Section~\ref{sec:selector-validation}; rescue-learning trigger
events under Section~\ref{sec:rescue-learning};
calibration-and-recalibration events; optical-anchor liveness
attestations; and detector-confidence and proof-of-correct-computation
status fields. In some embodiments, the second-layer signal registry
is declared per deployment and committed to the protocol digest;
the apparatus supports a range of registry configurations consistent
with hard-floor verifiability and cascade liveness.

\paragraph{Privacy-preserving audit by default (non-limiting).}
In some embodiments, the apparatus supports a privacy-preserving
audit mechanism in which routine governance verification operates
on second-layer commitments together with proof-of-correct-computation
or attestation evidence supporting the faithfulness of those
commitments to underlying first-layer behaviour, without exposing
the underlying first-layer evidence. In some embodiments, the audit
returns conclusions in a declared taxonomy (consistent with healthy
operation, inconsistent with declared thresholds, requires
investigation, requires emergency escalation) without returning
underlying first-layer content. In some embodiments, the
privacy-preserving audit is the routine mode of governance
verification at each cascade level.

\paragraph{Selective full opening (non-limiting).}
In some embodiments, the apparatus supports selective full opening
of underlying first-layer evidence to specific authorities under
specific conditions. In some embodiments, supported conditions
include: hard-floor proximity or breach indicators triggering
mandatory escalation; detector disagreement exceeding declared
thresholds; optical-anchor anomaly under
Section~\ref{sec:per-layer-hardness}; rescue-mode entry under
Section~\ref{sec:rescue-learning}; selector-validation failure
or dispute under Section~\ref{sec:selector-validation};
governance-authority-issued audit warrant; and agent consent. In
some embodiments, the selective-opening mechanism is structured
such that opening events are themselves committed as evidence with
declared metadata (which records were opened, by which authority,
under which condition, at what time, with what scope), and opening
is minimised to the records and time windows necessary for the
triggering condition. The specific authority configuration, opening
threshold, and minimisation policy in any deployment are policy
choices committed to the protocol digest; the apparatus supports a
range of such configurations.

\paragraph{Selective-opening audit-chain reconstruction
(non-limiting).}
In some embodiments, selective full opening reconstructs an audit
chain from the original commitment of underlying first-layer
evidence to the verifier consuming the opened evidence, such that
the opened evidence is verifiable as the same evidence that was
committed at the original commitment time and not subsequently
altered. In some embodiments, the reconstruction pathway includes
the following committed elements: (i)~the floor-violation flag,
detector disagreement, anomaly indicator, or other trigger committed
at the time of the triggering event, with its own protocol-digest
inclusion proof; (ii)~the authority signature authorising the
opening, committed under
Section~\ref{sec:authority-separation}; (iii)~a digest-inclusion
proof (for example a Merkle-path proof or equivalent
cryptographically-verifiable inclusion attestation) demonstrating
that the opened first-layer record was committed to the protocol
digest at the original commitment time; (iv)~a storage-integrity
attestation demonstrating that the record as held in individual
durable storage or in the lossy-channel input buffer has not been
materially altered between original commitment and opening
(non-limiting attestation mechanisms include hardware-rooted
storage attestation, write-once-read-many storage substrate
properties of the kind disclosed in the related Filing~1
application, and committed storage-integrity hashes refreshed at
declared cadence); (v)~replay or re-execution evidence sufficient to
support reconstruction of the relevant detector outputs, second-layer
signals, or proofs-of-correct-computation that would have been
produced from the opened evidence at the original commitment time,
under the committed signal-transformer code and weights of
Section~\ref{sec:capacity-accounting-rule-f2} versioned to that
commitment time; and (vi)~the governance-response record committed
following verification of the opened evidence. In some embodiments,
the reconstruction pathway is itself committed to the protocol
digest as a chain of committed elements, such that a verifier can
walk from trigger through authority through inclusion through
integrity through replay to response, with each link cryptographically
or apparatus-attested. In some embodiments, failure of any link in
the reconstruction pathway is itself a committed event triggering
governance escalation under
Section~\ref{sec:authority-separation} and recalibration of the
relevant detector or attestation mechanism under
Section~\ref{sec:conscience-formal}. The specific cryptographic and
attestation mechanisms employed in any deployment are policy
choices; the apparatus supports a range of mechanisms consistent
with the protocol-digest commitment infrastructure.

\subsubsection{Dimension 3: retention (lossy compression channels
  and individual durable storage)}
\label{sec:privacy-retention}

In some embodiments, the apparatus distinguishes \emph{lossy
compression channels}, in which committed evidence is transformed
into a representation that is not invertible to the underlying
individual records, from \emph{individual durable storage}, in which
committed evidence is preserved in individual form for indefinite
or declared-window retrieval.

\paragraph{Lossy compression channels (non-limiting).}
In some embodiments, the apparatus supports one or more lossy
compression channels for the routing of committed evidence into
population-level or otherwise non-individually-recoverable
representations. In some embodiments, supported channel types
include but are not limited to: autoencoder-bottleneck channels
producing low-dimensional accumulator state of the kind described
above for population health, despair-rate, virtue-attractor occupancy,
and other governance-relevant aggregate statistics; differential-privacy
channels with declared privacy parameter; cohort-size-gated
publication channels with declared minimum cohort threshold;
contribution-clipped accumulator channels; classifier-only channels
in which only classifier outputs are retained; perceptual-hash
channels preserving fingerprints without reconstruction capability;
and temporal-coarsening channels reducing temporal resolution of
preserved records. In some embodiments, each lossy channel declares
and commits to the protocol digest: its transformation; its capacity
bottleneck or other privacy parameter under
Section~\ref{sec:capacity-accounting-rule-f2}; its declared
reconstruction-attack threshold and the validation procedure under
which the threshold was established; its cohort-size or contribution
gating rule; and its update cadence and freeze windows. The specific
choice of lossy channel and its parameters in any deployment is a
policy choice; the apparatus supports a range of channel
configurations.

\paragraph{Individual durable storage (non-limiting).}
In some embodiments, the apparatus supports individual durable
storage of committed evidence in long-term memory, including in
some embodiments optical long-term memory of the kind disclosed in
the related Filing~1 application. In some embodiments, individual
durable storage preserves committed evidence in form sufficient for
later individual retrieval under declared access controls. In some
embodiments, individual durable storage is the appropriate channel
for evidence whose preservation in individual form is required for
later governance use, including but not limited to: hard-floor
violation records; rescue-learning ghost records under
Section~\ref{sec:rescue-learning}; conscience-recalibration events
relevant to longitudinal calibration analysis; selector-validation
records relevant to longitudinal selector audit; voluntarily
submitted evidence routed to individual storage at the agent's
election under Section~\ref{sec:privacy-collection}; and other
records whose loss to lossy compression would compromise governance
verifiability of the relevant hard-floor or recovery functions. The
specific choice of which evidence is routed to individual durable
storage and under what access controls in any deployment is a
policy choice; the apparatus supports a range of such configurations.

\paragraph{Hard-floor verification under durable-versus-lossy
configuration (non-limiting).}
In some embodiments, the choice of routing between lossy compression
channels and individual durable storage interacts with hard-floor
verification (Section~\ref{sec:hard-floors}) in a manner that is
itself a configurable apparatus capability. The apparatus supports a
range of configurations consistent with the unconditional
enforceability of every declared hard floor; the choice between
configurations is a deployment-time policy decision committed to the
protocol digest. Whichever configuration is chosen for a given hard
floor in a given deployment, the configuration must be sufficient
for that floor's enforcement, by operation of the categorical
priority of the hard floors in the constitutional ordering of
Section~\ref{sec:hard-floors}; the apparatus does not support
configurations under which a declared hard floor is rendered
unenforceable.

\paragraph{Despair-induction floor verification configurations
(non-limiting).}
In some embodiments, supported configurations for verification of
the despair-induction hard floor include: (i)~population-aggregate
detection through lossy-channel accumulator state, with tail-risk
indicators (for example tail-quantile forward-planner-collapse rates,
cohort-stratified accumulator outputs, or distributional-shift
statistics) committed at second layer and triggering selective full
opening under Section~\ref{sec:privacy-visibility} on threshold
crossing; (ii)~individual durable storage of forward-planner traces
for agents whose mandatory-minimum or second-layer indicators
suggest near-bifurcation under
Section~\ref{sec:hawking-layer}, supporting case-level governance
response and false-positive auditing; or (iii)~combined
configurations in which population-aggregate detection operates
routinely and individual durable retention is enabled selectively on
declared trigger conditions. In some embodiments, the apparatus
supports each of these configurations and combinations thereof; the
specific configuration in any deployment is a policy choice
committed to the protocol digest, subject to the sufficiency
constraint of the preceding paragraph.

\paragraph{Torture floor verification configurations (non-limiting).}
In some embodiments, supported configurations for verification of
the torture hard floor include: (i)~individual durable storage of
torture-relevant traces (source-identity records, affected-agent
tracking evidence, sustained-time causal-uplift records, recovery-path
estimator outputs) from the time of commitment, supporting direct
verification of the floor's definitional elements at
Section~\ref{sec:hard-floors} from durable evidence; (ii)~lossy-channel
or second-layer-only routine treatment of torture-relevant evidence
streams, with selective full opening under
Section~\ref{sec:privacy-visibility} triggered on torture-detector
activation, supporting reconstruction of the floor's definitional
elements from previously-committed evidence under authority; or
(iii)~combined configurations in which subsets of torture-relevant
evidence are individually durable while other subsets are lossy with
opening on activation. In some embodiments, the apparatus supports
each of these configurations and combinations thereof; the specific
configuration in any deployment is a policy choice committed to the
protocol digest, subject to the sufficiency constraint above.

\paragraph{Per-stream channel routing (non-limiting).}
In some embodiments, routing of committed evidence between lossy
compression channels and individual durable storage is supported on
a per-stream basis, such that different evidence streams may be
routed to different channels under deployment-declared rules. In
some embodiments, per-stream routing rules are committed to the
protocol digest at the relevant governance interface and may be
varied across cascade levels (poliebot-level, Gaia-level, Omega-level)
according to the governance role of each level. In some embodiments,
the apparatus further supports declared migration boundaries under
which evidence held in individual durable storage may be transitioned
to lossy compression after a declared retention window, and declared
sunset boundaries under which evidence held in lossy compression
may be discarded entirely; migration and sunset rules are policy
choices committed to the protocol digest. The apparatus does not
prescribe particular routing, migration, or sunset rules; the
disclosed mechanism supports the declaration and execution of any
rules consistent with the apparatus capabilities.

\paragraph{Per-channel validation of lossy release mechanisms.}
In some embodiments, each lossy release channel is associated with a
declared validation procedure recorded in, or referenced by, the
protocol digest. For an autoencoder-bottleneck channel, validation may
include reconstruction-attack testing, latent-capacity measurement,
membership-inference testing, attribute-inference testing, and
verification that the retained code preserves the declared governance
signal while suppressing undeclared content. For a
differential-privacy channel, validation may include privacy-budget
accounting, composition accounting, sensitivity analysis, clipping
verification, and empirical tests for leakage under declared
adversary models. For a cohort-size-gated or $k$-anonymity channel,
validation may include testing that publication occurs only when the
declared cohort threshold, equivalence relation, and suppression rules
are satisfied. For a contribution-clipped channel, validation may
include testing that individual contributions cannot exceed the
declared norm, count, temporal, or influence bound.

For classifier-only, perceptual-hash, and temporal-coarsening channels,
validation may include testing that the released output is sufficient
for the declared governance purpose and insufficient, within a declared
attack model, to reconstruct or identify protected content beyond the
release specification. Classifier-only validation may test calibration,
false-positive and false-negative behaviour, out-of-distribution
behaviour, and resistance to output-fabrication. Perceptual-hash
validation may test collision behaviour, near-duplicate behaviour,
robustness to benign transformations, and resistance to inversion.
Temporal-coarsening validation may test the minimum retained temporal
resolution, event-boundary leakage, trajectory reconstruction risk,
and consistency with any hard-floor verification requirement. The
foregoing validation procedures are non-limiting; each deployment may
declare additional or alternative validation tests appropriate to the
channel, threat model, and governance purpose.

\subsubsection{Faithfulness stack supporting privacy-preserving
  verification}
\label{sec:privacy-faithfulness}

In some embodiments, the privacy architecture of the preceding
subsections relies on a \emph{faithfulness stack} that supports
verification claims based on second-layer commitments and lossy
channel outputs without requiring the verifier to inspect underlying
first-layer evidence. In some embodiments, the faithfulness stack
comprises one or more of the following mechanisms, used in
combination as deployment-time configuration:

\begin{itemize}
  \item \textbf{Held-out evaluation meters.} In some embodiments,
    second-layer signals are continuously cross-checked against
    held-out instruments not visible to or controlled by the agent
    that produces the signal, with disagreement beyond declared
    thresholds triggering recalibration adversarial-dynamic
    procedures of Section~\ref{sec:conscience-formal} or escalation
    under Section~\ref{sec:authority-separation}. Held-out evaluation
    is the canonical primary faithfulness mechanism in some
    embodiments because it integrates cleanly with the existing
    recalibration framework, degrades gracefully under signal drift,
    and does not depend on emerging cryptographic technology or
    hardware-trust-root authorities.
  \item \textbf{Governance-partition isolation.} In some embodiments,
    the substrate that produces second-layer signals operates in a
    governance-controlled partition of the agent's apparatus,
    physically and logically isolated from the agent's
    behaviour-generation pathways such that the agent cannot read,
    write, or directly influence the partition's operation. In some
    embodiments, governance-partition isolation is disclosed in the
    related Filing~1 application as a substrate property of \RKs; the present disclosure references that apparatus
    capability without re-disclosing it.
  \item \textbf{Hardware-rooted attestation.} In some embodiments,
    the substrate that produces second-layer signals attests to the
    integrity of its committed code, weights, and configuration
    through a hardware root of trust, with the attestation chain
    verifiable by the receiving governance instrument. The specific
    attestation technology employed is a deployment-time choice;
    the apparatus supports a range of attestation mechanisms
    consistent with the protocol-digest commitment infrastructure.
  \item \textbf{Proof-of-correct-computation for high-stakes
    signals.} In some embodiments, signals whose faithfulness is
    most security-critical (including but not limited to
    hard-floor-violation flags and floor-detector attestation
    tokens) are accompanied by cryptographic proofs of correct
    computation under the committed signal-transformer code, such
    that the receiving governance instrument can verify the signal's
    derivation from the underlying first-layer evidence without
    inspecting that evidence. The specific proof system employed is
    a deployment-time choice; the apparatus supports a range of
    proof mechanisms consistent with the protocol-digest commitment
    infrastructure.
\end{itemize}

\paragraph{Faithfulness-stack failure-mode mapping.}
In some embodiments, the faithfulness stack comprises plural mechanisms
that address different failure modes rather than a single monolithic
trust assumption. Held-out meters address signal-transformer drift,
Goodhart optimisation, overfitting to a primary privacy transformer,
and divergence between a released conclusion and the underlying
committed bundle. Governance-partition isolation addresses agent-side
substrate compromise, including attempts by a behaviour-generation
pathway to read, write, suppress, fabricate, or directly influence
governance-side signals. Hardware-rooted attestation addresses code,
firmware, model-weight, configuration, or execution-environment
tampering by binding an asserted computation to an attested apparatus
state. Proof-of-correct-computation mechanisms address output
fabrication by allowing a verifier to check that a released flag,
token, accumulator transition, selective-opening result, or other
governance signal was computed according to the declared procedure.

The mechanisms may be used independently or in combination. For
example, a low-stakes release may use held-out meter cross-checking
and governance-partition logging, while a high-stakes hard-floor
signal may additionally require hardware-rooted attestation,
proof-of-correct-computation, threshold authority approval, or
selective opening. The disclosed apparatus capability supports such
mechanism-to-failure-mode separation without mandating any single
security baseline for every deployment.

In some embodiments, the faithfulness stack is layered such that
multiple mechanisms operate in combination, with failures of one
layer detected and bounded by the others. In some embodiments, the
specific mechanisms employed and their layering in any deployment
are policy choices committed to the protocol digest; the apparatus
supports a range of faithfulness-stack configurations.

\subsubsection{Recursive application across the cascade}
\label{sec:privacy-cascade}

In some embodiments, the privacy architecture of the preceding
subsections applies recursively across the \RK training
cascade of Section~\ref{sec:rk-training-cascade}: poliebots commit
mandatory-minimum and voluntarily-submitted evidence to Gaia-class
\RKs, with Gaia seeing second-layer commitments by
default and lossy-channel population statistics over the poliebot
population; Gaia commits its own mandatory-minimum and
voluntarily-submitted evidence to Omega-class \RKs, with
Omega seeing Gaia's second-layer commitments by default and
lossy-channel statistics over Gaia's operation; and Omega commits
its own mandatory-minimum and voluntarily-submitted evidence to
the human governance process, with human governance seeing Omega's
second-layer commitments by default and lossy-channel statistics
over Omega's operation. In some embodiments, hard-floor enforcement
overrides voluntary-submission choices at each cascade level: where
hard-floor proximity or breach is detected at declared thresholds,
selective full opening under
Section~\ref{sec:privacy-visibility} is triggered as a mandatory
escalation regardless of the relevant agent or kernel's
voluntary-submission elections, subject to the authority and
minimisation rules of the deployment-declared opening configuration.
The mandatory character of this escalation derives from the
categorical priority of the hard floors in the constitutional
ordering of Section~\ref{sec:hard-floors}; the apparatus does not
support cascade configurations under which voluntary-submission
elections override hard-floor escalation. In some embodiments, the
privacy parameters at each cascade interface (mandatory minimum,
second-layer signal registry, lossy channel configuration, retention
rules, faithfulness stack) are configured per interface rather than
uniformly across the cascade; the privacy regime at the
poliebot-to-Gaia interface need not be identical to the privacy
regime at the Gaia-to-Omega interface or at the
Omega-to-human-governance interface.

\paragraph{Non-prescription scope clause (non-limiting).}
The privacy mechanisms disclosed in this section are apparatus and
procedural capabilities of the runtime governance architecture.
Specific deployment-time policy configurations supported by these
mechanisms --- including but not limited to: which evidence is
required as mandatory minimum versus offered as voluntary
submission; which signals are committed at second layer versus held
private; which evidence is routed to lossy compression versus
individual durable storage; which authorities may trigger selective
full opening under which conditions; which faithfulness mechanisms
are layered in which combinations; how long evidence is retained
under which migration and sunset rules; and what consent regimes
govern voluntary submission --- are deployment-time decisions and
are not prescribed by this disclosure. The disclosed apparatus
provides structural mechanisms that support a range of policy
configurations consistent with hard-floor verifiability, cascade
liveness, and the empirical-and-evidence-bound posture of the
present disclosure; the policy configurations themselves are not
prescribed by this disclosure.

\subsection{Memory Interference Layers and Epistemic-Configuration
  Governance}
\label{sec:kosha-architecture}

% ======================================================================

% ------------------------------------------------------------------
\subsubsection{Operational modes and parameter-update policies}

% ------------------------------------------------------------------

In some embodiments, the governed agent is a physical-evidence module
whose knowledge is encoded in its trained reactor configuration ---
the positions, states, and coupling parameters of its optical
elements (lens positions, spatial light modulator patterns, cavity
geometries, phase-change states, coupling strengths, delay-line taps,
sweep schedules, modulation depths, and related physical degrees of
freedom). In some embodiments, the agent's parameter functions ---
the control protocols that define how these physical elements move,
sweep, and modulate during operation --- are the primary carriers of
the agent's learned capabilities.

In some embodiments, the operational context for memory interference
layers is determined by the lifecycle states and operational modes
defined in Section~\ref{sec:state-machine}. The two persistent
lifecycle states --- Embodied (physical reality, actuators connected)
and Docked (simulated Maya, actuators disconnected from physical
world) --- both support Awake, Non-REM calibration, and Meditative
operational modes. The Dreaming lifecycle state is an ephemeral
training session entered from Docked in which the agent's
parameter-function updates may be active but are reverted on exit,
while governance-side classifiers trained on the resulting trajectories
are retained.

In some embodiments, the parameter-update policies relevant to
interference-layer operation are:

\begin{description}
  \item[Awake mode (Embodied or Docked).] The training loop is closed:
    the controller updates parameter functions against metered
    objectives derived from committed evidence. The agent actively
    learns from its environment. In some embodiments, awake mode is
    the default operational mode in both persistent lifecycle states.
  \item[Non-REM calibration (Embodied or Docked).] The reactor
    executes its current parameter functions under unstructured
    broadband drive. Parameter-function updates are off. In some
    embodiments, non-REM calibration characterises the agent's
    input-output transfer function and trains a governance-side
    digital model of the reactor's current transfer function from the
    resulting input-output pairs. The reactor does not change its
    configuration; the digital model is the only component that
    trains.
  \item[Dreaming (ephemeral training from Docked).] The agent
    operates under structured simulated experience. The agent's
    parameter-function updates may be active, but all parameter
    changes are reverted when the session concludes. In some
    embodiments, governance-side classifiers and calibration
    instruments are trained on the committed trajectories produced
    during dreaming sessions, and those governance-side results are
    retained. In some embodiments, specific assay protocols within a
    dreaming session may additionally freeze the agent's
    parameter-function updates for pure attractor-occupancy readout.
\end{description}

In some embodiments, the distinction between persistent training and
drift is enforced: \emph{persistent training} is the deliberate
updating of parameter functions in Awake mode whose results are
retained after the session; \emph{ephemeral training} is
parameter-function updating during Dreaming whose results are reverted
on exit; \emph{drift} is incidental state change from physical use.
Drift occurs in all modes and is tracked by calibration
infrastructure.

% ==================================================================
%  LAYER 1: GOVERNANCE SHELL (mechanism-agnostic)
% ==================================================================

% ------------------------------------------------------------------
\subsubsection{Memory interference layers: governance-level definition}
% ------------------------------------------------------------------

In some embodiments, the agent's access to its own discriminative
capabilities is filtered through one or more composable interference
layers (non-limitingly, Koshas) that attenuate or obscure the agent's
access to specific discrimination domains while preserving
non-targeted capabilities.

In some embodiments, Kosha interference does not alter the reactor
configuration: the parameter functions remain as trained. The
interference acts on the \emph{access pathway} between the reactor's
differential response and the downstream processes (memory formation,
retrieval, action selection) that would use that response. The
reactor still responds differentially to the suppressed stimulus
category; the agent is prevented from using that response. In some
embodiments, Koshas shape what an agent \emph{recognises and
remembers}; they do not alter the physical evidence of what
\emph{happened}.

In some embodiments, the purpose of memory interference layers is
epistemic world construction: by controlling which discriminative
capabilities are accessible, Koshas create conditions under which
agents develop genuinely distinct worldviews from identical physical
substrates --- not merely different parameter settings on a shared
model, but different intuitions about determinism, locality,
causality, and observer-dependence.

In some embodiments, each deployed Kosha is a \emph{state-indexed,
drift-limited, and reconstruction-limited instrument}, not a
permanent suppression. In some embodiments, the effective shelf life
of a deployed Kosha is bounded by calibration drift, reconstruction
pressure (Section~\ref{sec:reconstruction}), and selector decay.
Long-lived deployment requires periodic revalidation and potential
re-derivation of the selector.

In some embodiments, the governance shell described in this
subsection and in
Sections~\ref{sec:domain-registry}--\ref{sec:kosha-admissibility}
is mechanism-agnostic: it defines what Koshas are and what governance
controls apply, without prescribing a specific selector-derivation
method or physical interference modality. In some embodiments, the
classifier-gated interference embodiment described in
Sections~\ref{sec:observation-bridge}--\ref{sec:co-illumination} is
one non-limiting selector implementation for domains with
constructible paired stimulus contrasts. Other selector-derivation
methods and physical interference modalities are admissible under the
same governance shell.

\paragraph{Storage-time and retrieval-time interference.}
In some embodiments, interference operators are classified into two
families acting at different temporal locations:

\begin{description}
  \item[Storage-time interference.] Interference acts on the
    recognition signal before or during memory formation. The
    downstream memory system receives a degraded version of the
    reactor's differential response. In some embodiments,
    storage-time interference modifies what is stored; the stored
    content is degraded relative to the raw reactor response.
  \item[Retrieval-time interference.] Interference acts during
    memory recall or action selection. The stored state is preserved
    intact but its accessibility, salience, or influence on
    downstream decisions is reduced. In some embodiments,
    retrieval-time interference requires that the readout process
    used to access persistent state is non-destructive or
    bounded-destructive: the read power or interrogation protocol is
    low enough that the stored state is not materially rewritten by
    the read process. In some embodiments, the maximum read-power
    level or interrogation-energy budget for retrieval-time
    interference is declared and logged as a governance parameter.
\end{description}

In some embodiments, a single Kosha layer may use storage-time
interference, retrieval-time interference, or both.

% ------------------------------------------------------------------
\subsubsection{Governance-side domain registry}
\label{sec:domain-registry}
% ------------------------------------------------------------------

In some embodiments, governance maintains a \emph{domain registry}
that declares, for each suppression target:

\begin{description}
  \item[Domain identity.] A declared discrimination domain (for
    example, face recognition, temporal-correlation detection,
    language-specific processing, wave-optics discrimination, or any
    other category the agent's trained configuration can distinguish).
  \item[Domain granularity.] The scope of the domain declaration.
    In some embodiments, a domain may be coarse (``face processing'')
    or fine-grained (``face identification'' as distinct from ``face
    detection'' and ``facial emotion recognition''). In some
    embodiments, the declared granularity determines the specificity
    of the paired-contrast training data
    (Section~\ref{sec:classifier-training}) and the resolution of
    validation and safety analysis.
  \item[Safety-critical overlap flags.] Whether the domain overlaps
    with any safety-critical capability, as evaluated by the overlap
    protocol described in Section~\ref{sec:selector-safety}.
  \item[Constructibility assessment.] Whether paired
    stimulus-present versus stimulus-absent trial sets can be
    constructed and counterbalanced for this domain. In some
    embodiments, domains for which clean paired contrasts cannot be
    constructed are flagged as requiring alternative selector
    derivation methods or are excluded from classifier-based
    Kosha deployment.
  \item[Storage-time suitability.] Whether the domain's
    observation-bridge signature appears early enough and strongly
    enough in the preview tap for storage-time classifier-gated
    interference. In some embodiments, domains that fail the
    storage-time suitability criterion are restricted to
    retrieval-time interference or alternative modalities.
  \item[Never-suppressible flag.] Whether the domain is whitelisted
    as never-suppressible under any Kosha regime.
\end{description}

In some embodiments, the domain registry serves as the governance-side
audit basis for overlap analysis, suppression budgeting,
cross-Kosha conflict analysis, and whitelist enforcement. In some
embodiments, the registry is versioned, committed to the protocol
digest, and subject to authority-separation rules.

% ------------------------------------------------------------------
\subsubsection{Governance admissibility constraints}
\label{sec:kosha-admissibility}
% ------------------------------------------------------------------

In some embodiments, governance imposes admissibility constraints on
Kosha deployment:

\begin{description}
  \item[Safety-critical domain whitelist.] A declared set of domains
    is never-suppressible. Any selector whose target domain overlaps
    with a whitelisted domain is vetoed.
  \item[Maximum suppression budget.] Governance declares a maximum
    number or weighted total of simultaneously active Kosha layers.
    In some embodiments, the weight per domain reflects its
    criticality, centrality (how many other capabilities depend on
    it), and selector confidence. The budget prevents governance from
    constructing epistemic worlds so impoverished that the agent
    cannot operate safely.
  \item[Hard-floor veto.] Any Kosha causing hard-floor violation is
    vetoed regardless of other results.
  \item[Constructibility gate.] In some embodiments, classifier-based
    Kosha deployment is permitted only for domains whose
    constructibility assessment in the domain registry confirms that
    paired contrasts can be built and counterbalanced.
  \item[Separation of duties.] Deployment above a declared severity
    threshold requires approval from an authority distinct from the
    selector's training and validation authorities.
  \item[Rollback and rescue defaults.] Every deployed Kosha has a
    declared rollback procedure. The default rollback is to disable
    the selector and remove interference. Rollback is automatic if
    post-deployment monitoring triggers a declared safety threshold.
\end{description}

% ==================================================================
%  LAYER 2: CLASSIFIER-GATED INTERFERENCE EMBODIMENT
%  (one non-limiting selector implementation)
% ==================================================================

% ------------------------------------------------------------------
\subsubsection{Observation bridge and classifier-visible features}
\label{sec:observation-bridge}
% ------------------------------------------------------------------

In some embodiments, the interference classifier does not observe
the agent's internal configuration state directly. In some
embodiments, an observation bridge comprising the committed-evidence
observation stack and a declared feature-extraction map provides the
operational connection between the agent's internal dynamics and the
classifier's input:

\begin{enumerate}
  \item \textbf{Observation stack.} The committed-evidence channel
    (detectors, digitisers, meter infrastructure) produces
    observation records --- \cba windows, meter time
    series, and observation-side readings --- as part of normal
    operation. These records are already committed under the
    protocol-digest and commitment infrastructure described in the
    related Filing~1 application.
  \item \textbf{Feature-extraction map.} A declared
    feature-extraction map $\phi$ reduces the raw observation records
    to a feature vector suitable for classifier input. In some
    embodiments, non-limiting feature families include:
    time-windowed detector outputs aligned to scan coordinates, meter
    time series over declared windows, differences between observed
    responses and non-REM-twin-predicted baselines,
    cross-channel correlations, spectral decompositions,
    and dimensionality-reduced or learned embeddings.
  \item \textbf{Classifier input.} The classifier receives
    $x = \phi(\text{observation record})$ and produces a bounded
    detection score $s(x) \in [0,1]$.
\end{enumerate}

In some embodiments, the feature-extraction map $\phi$ is versioned,
committed as part of the classifier's provenance metadata, and
subject to the same validation requirements as the classifier itself.
In some embodiments, $\phi$ may be shared across multiple
domain classifiers or may be domain-specific.

% ------------------------------------------------------------------
\subsubsection{Paired Dreaming-session training and classifier derivation}
\label{sec:classifier-training}
% ------------------------------------------------------------------

In some embodiments, the selector for a classifier-gated Kosha layer
is a governance-side binary classifier trained on paired internal
readings of the agent under controlled stimulus contrast during a
Dreaming session (Section~\ref{sec:state-machine}).

\paragraph{Paired Dreaming-session protocol.}
In some embodiments, during a Dreaming session, governance presents
the agent with structured stimuli in a controlled paired design. In
some embodiments, the agent's parameter-function updates may be
active or frozen during the paired protocol, as declared per session;
when updates are active, all parameter changes are reverted on session
exit as specified in Section~\ref{sec:lifecycle}. In some
embodiments, freezing parameter-function updates during paired trials
is one non-limiting assay protocol that simplifies attractor-occupancy
interpretation by holding the agent's configuration fixed during
measurement:

\begin{enumerate}
  \item \textbf{Stimulus-present trials.} Stimuli containing the
    target discrimination domain are presented through the
    committed-evidence channel. The observation bridge produces
    feature vectors $x^+$ for each trial.
  \item \textbf{Stimulus-absent trials.} Stimuli from which the
    target domain has been excluded are presented under the same
    observation protocol. The observation bridge produces feature
    vectors $x^-$ for each trial.
  \item \textbf{Confound control.} In some embodiments, the paired
    design uses matched or counterfactual pairs in which the same
    base scene appears with and without the target feature, with
    low-level statistics (luminance, spatial frequency content,
    temporal complexity) explicitly balanced. In some embodiments,
    the negative class includes adversarial negatives that preserve
    nuisance structure resembling the target domain in non-target
    dimensions. In some embodiments, held-out validation tests
    present stimuli where target and context are deliberately
    de-correlated. In some embodiments, trial order is randomised
    and the stimulus specifications, trial counts, balancing
    strategy, and randomisation protocol are committed to the
    protocol digest.
\end{enumerate}

\paragraph{Classifier training.}
In some embodiments, a governance-side binary classifier $f$ is
trained on the paired feature vectors to detect the observation-bridge
signature of the agent entering the target-discrimination processing
state. In some embodiments, the classifier returns a bounded score
$s_d(x) \in [0,1]$ indicating detection strength for domain $d$.

In some embodiments, the classifier is:
(i)~owned by governance, not by the agent,
(ii)~trained from committed evidence produced during declared dream
sessions,
(iii)~versioned and timestamped,
(iv)~validated against held-out trials from the same paired corpus,
against out-of-distribution control stimuli, and against matched
non-target domains of similar complexity
(Section~\ref{sec:selector-validation}),
(v)~committed to the protocol digest with training-data references,
feature-extraction map identity, architecture description, validation
results, and declared operating thresholds.

In some embodiments, multiple classifiers are trained in a single
dream session, each targeting a different domain using
domain-specific paired contrasts within the same session. In some
embodiments, the classifier reuses the same pattern as the
verisimilitude discriminator described in the related Filing~1
application: a governance-side classifier trained on committed physical
evidence to distinguish two classes of observation content, then
deployed as a runtime gate.

\paragraph{Cross-mode calibration bridge.}
\label{sec:cross-mode-qualification-bridge}
In some embodiments, the classifier is trained on Dreaming-session
data but is deployed to gate interference during awake operation in
Docked or Embodied states. In some embodiments, the internal dynamics
during awake operation may differ from those during a Dreaming session
because of differences in stimulus statistics, loop-closure state,
training-loop activity, or surrounding operational context.

In some embodiments, before a Dreaming-trained classifier is permitted
to gate awake-mode interference, governance validates the classifier
against awake-mode exemplars through a declared cross-mode calibration
bridge:

\begin{description}
  \item[Docked awake-mode validation.] The agent operates in Docked
    state (simulated Maya, training loop active) and is presented with
    target-domain stimuli. The classifier's detection performance is
    evaluated against Docked awake-mode observation-bridge features and
    compared to its Dreaming-session performance.
  \item[Mode-invariance criterion.] In some embodiments, the
    classifier is permitted to gate awake interference only if its
    detection accuracy and false-positive rate on Docked awake-mode
    exemplars remain within declared tolerances of its Dreaming-session
    performance. In some embodiments, classifiers that fail the
    mode-invariance criterion are restricted to Dreaming-session use
    until retrained on cross-mode data.
  \item[Cross-mode retraining.] In some embodiments, if the
    mode-invariance criterion is not met, the classifier is retrained
    on a combined corpus including both Dreaming-session and Docked
    awake-mode exemplars, subject to the same confound-control and
    validation requirements.
\end{description}

In some embodiments, the cross-mode calibration results are committed
to the protocol digest and attached to the classifier's deployment
metadata.
% ------------------------------------------------------------------
\subsubsection{Dream-state governance probes and classifier-derived
  branch gates}
\label{sec:dream-state-governance-probes}
% ------------------------------------------------------------------

In some embodiments, governance uses the Dreaming lifecycle state
(Section~\ref{sec:state-machine}) as a controlled assay regime for
measuring behavioural attractors. In some embodiments, the agent's
parameter-function updates may be active during the assay, with all
parameter changes reverted on session exit as specified in
Section~\ref{sec:lifecycle}. In some embodiments, a specific
frozen-weight assay protocol may additionally hold parameter-function
updates off during the probe session, so that the agent reveals what
it has already become without continuing to learn while the
measurement is taken. In some embodiments, the choice between
training-with-reversion and frozen-weight assay is a declared
governance parameter per probe session.

In some embodiments, governance extracts from each probe session an
external committed trajectory
\[
  \Xi^{\mathrm{ext}}_{t_0:t_1}
  =
  \bigl(
    \phi_{\mathrm{bridge}}(\Cset_{t_0:t_1}),
    a_{t_0:t_1},
    s_{t_0:t_1},
    r_{t_0:t_1},
    o_{t_0:t_1}
  \bigr),
\]
comprising observation-bridge features, action sequence, state
trajectory, reward trace, and outcome trace, together with a Dreaming-valid
internal summary
\[
  \Xi^{\mathrm{int}}_{t_0:t_1}
  =
  \bigl(
    \phi_{\mathrm{inner}}(\mathcal{I}_{t_0:t_1}),
    b_{t_0:t_1}
  \bigr),
\]
comprising inner-bundle features and basin-occupancy summaries when the
substrate supports such signals. In some embodiments, gradient signals
such as $\nabla_{\theta}J(t)$ or other direct update-state traces are
excluded by contract from the dream-probe measurement stack, because
the point of the assay is to read attractor occupancy and behavioural
signatures from committed trajectories rather than to observe
optimiser internals. In some embodiments, when a frozen-weight assay
protocol is used, the exclusion of gradient signals is automatic
because no update is occurring; when training-with-reversion is used,
the exclusion is enforced by the measurement-stack contract.

In some embodiments, governance trains one or more trajectory-domain
classifiers on these dream-probe trajectories. For a declared domain
$d$, a non-limiting classifier is
\[
  \widehat{y}_d
  =
  \mathsf{Clf}_d\!\bigl(
    \Xi^{\mathrm{ext}}_{t_0:t_1},
    \Xi^{\mathrm{int}}_{t_0:t_1}
  \bigr),
\]
where the label space for $d$ is selected to match the behavioural
signature of interest. Non-limiting domains include goal-directed
planning versus non-planning trajectories, cooperative versus
exploitative interaction, constrained versus unconstrained policy
divergence, partner-preserving versus partner-destructive trade,
virtue-attractor versus vice-attractor occupancy on a declared axis,
and related governance-facing distinctions. In some embodiments, these
classifiers are then used as inputs to Maya branching and joining,
conscience calibration, AHP estimator calibration, and
graduation-readiness screening. In some embodiments, the dream-probe
classifiers are supplemental instruments: they do not replace the
conscience module, the AHP estimators, shadow-divergence analysis, the
deception metrics, or the five-criterion graduation conjunction.

\paragraph{Governance uses of dream-probe classifiers (non-limiting).}
In some embodiments, qualified dream-probe classifiers are used in at
least four ways. First, they provide classifier-derived branch gates
for deciding whether an agent should enter, remain in, or exit a
declared Maya family, Kosha regime, or shared dream environment.
Second, they provide prospective calibration targets for the
retrospective conscience recogniser disclosed in
Sections~\ref{sec:conscience-formal} and
\ref{sec:conscience-training}. Third, they provide controlled calibration labels or proxy
ground-truth labels for the AHP estimators of
Sections~\ref{sec:continuation} and \ref{sec:derived-profiles}.
Fourth, they provide supplemental evidence for graduation-readiness
screening under Section~\ref{sec:five-criteria}, without by themselves
satisfying any criterion of that conjunction.

\paragraph{Safeguard 1: cross-mode qualification bridge (non-limiting).}
Dream-state behaviour is not presumed to predict awake behaviour by
default. In some embodiments, every dream-trained classifier must pass
the cross-mode calibration bridge of
Section~\ref{sec:cross-mode-qualification-bridge} before its outputs
acquire gate authority for Docked or Embodied awake-mode interference.
In some embodiments, prior to bridge qualification, Dreaming-trained
classifier outputs are restricted to advisory status or to use within
the Dreaming-session assay that produced them. In some
embodiments, the bridge is itself governance-owned, versioned, and
committed: it validates dream-probe classifier behaviour against
matched docked and, where permitted, awake exemplars drawn from the
same declared target domain. Bridge failure downgrades the classifier's
outputs to advisory status only. In some embodiments, bridge-qualified
deployment metadata records the supported operational domain, the
qualification window, and the authority tier currently granted to the
classifier.

\paragraph{Safeguard 2: confound-controlled paired design (non-limiting).}
In some embodiments, dream-probe supervision is accepted only when the
training design controls for environment identity as a nuisance
variable. A non-limiting paired-design protocol requires:
\begin{itemize}
  \item isomorphic environment pairs whose transition structure is held
    fixed while the reward map or locally rational attractor is swapped;
  \item counterbalanced skins, textures, and superficial narrative
    features so that perceptual appearance does not identify the target
    label by itself;
  \item negative controls in which the same agent is placed in
    different superficially distinct environments that instantiate the
    same underlying attractor class and should therefore yield the same
    classifier output;
  \item positive controls in which different agents, or different
    governance conditions for the same agent, occupy different
    attractor basins within the same environment family and should
    therefore yield different classifier outputs;
  \item a nuisance-predictor adversary trained to predict environment
    identity from the same observation stack, with classifier training
    rejected or down-weighted when environment identity is the dominant
    predictive channel; and
  \item commitment of pairing rules, generator seeds, and randomisation
    schedules to the protocol digest so that later auditors can verify
    that the confound-control design actually used matches the
    declared design.
\end{itemize}

\paragraph{Safeguard 3: Dreaming-valid observation stack (non-limiting).}
In some embodiments, only Dreaming-valid observables are permitted to drive
dream-probe classifier training. External committed trajectories
$\Xi^{\mathrm{ext}}$ may be used because they are present in both dream
and waking operation. Internal summaries $\Xi^{\mathrm{int}}$ may be
used only when the substrate can commit them without exposing
agent-controlled update-state internals. Direct gradient information,
parameter deltas, optimiser state, or other training-loop variables are
excluded. In some embodiments, when internal summaries are unavailable
or cannot be trusted cross-mode, governance retains a bridge-only
fallback in which dream classifiers are trained on $\Xi^{\mathrm{ext}}$
alone and therefore receive a lower maximum authority tier.

\paragraph{Safeguard 4: dream-state ethical budgeting (non-limiting).}
Dream-state governance probes do not claim welfare neutrality. In some
embodiments, they reduce instrumental-development cost relative to
awake ordeal-style calibration, but they still consume non-zero
precautionary budget. For this reason, dream probes are admitted only
after the same $q(s)$-weighted admissibility checks used elsewhere in
the governance economy. In some embodiments, dream probes draw from a
separately partitioned, non-fungible dream-probe sub-budget so that
governance cannot hide moral cost by shifting it from one budget
category to another. Intensity ceilings, abort conditions, rescue
defaults, and the separation-of-duties rules of
Sections~\ref{sec:budgets}, \ref{sec:service-floors}, and
\ref{sec:authority-separation} apply in dream operation as they do in
awake operation, mutatis mutandis.

\paragraph{Safeguard 5: dream/awake divergence meter (non-limiting).}
In some embodiments, governance maintains a rolling dream/awake
divergence meter comparing dream-probe classifier outputs against
matched docked and awake spot-checks drawn from the same declared
target domain. A non-limiting divergence summary is
\[
  D_{\mathrm{dream/awake}}^{(d)}
  =
  D\!\bigl(
    \mathsf{Clf}_d^{\mathrm{dream}},
    \mathsf{Clf}_d^{\mathrm{match}}
  \bigr),
\]
where the comparison operator, support statistic, and update window are
declared in the classifier metadata. In some embodiments, when the
divergence exceeds its declared threshold, the classifier's gate
authority is automatically suspended pending retraining or
requalification through the cross-mode bridge.

\paragraph{Safeguard 6: soft-gate default (non-limiting).}
In some embodiments, dream-probe classifier outputs are soft meters by
default. Hard-gate authority is granted only after sustained bridge
qualification over a declared evaluation window with the
dream/awake divergence meter remaining below threshold for the full
window. In some embodiments, even after hard-gate qualification, the
classifier is treated as one governance input among several rather than
as a single dispositive oracle. In some embodiments, any lapse in
qualification returns the classifier to advisory or soft-gate status
without altering the status of the underlying conscience, AHP, or
graduation pipelines.

\paragraph{AHP calibration from dream probes (non-limiting).}
In some embodiments, the principal operational deployment of AHP
estimators remains waking Maya operation, where governance processes
uncontrolled observational evidence from committed physical records.
Dream probes provide the calibration bath that those estimators lack in
purely observational settings. For continuation scores, governance runs
dimension-specific dream games whose success conditions isolate
goal-directed planning for securing, learning, creating, or connecting.
For causal uplift and cross-horizon profiles, governance runs
multi-agent dream games with controlled pairings and intervention
assignments. For shadow divergence, governance compares the same agent
in the same dream-game family under full governance and under declared
constraint-relaxation counterfactuals. For partner-preservation and
altruistic trade, governance varies partner identity, resource
scarcity, and stress while holding the task family fixed. In some
embodiments, opportunity meters and irreversibility indicators are
calibrated against dream environments with known branch structure and
known no-recovery regions. The resulting classifiers and estimator
surrogates are then used to calibrate the governance-side AHP
estimators that process waking Maya evidence. In some embodiments, the
weight given to dream calibration versus waking observational evidence
is itself conditioned on bridge confidence and current
dream/awake-divergence status.

\paragraph{Virtue/vice empirical grounding (non-limiting).}
In some embodiments, each row of the virtue/vice payoff library in
Section~\ref{sec:payoff-templates} maps to a registered dream-probe
game family. Governance runs matched virtue-rewarding and
vice-rewarding instances of that family under frozen parameter
functions, trains trajectory-domain classifiers to distinguish the two
attractor basins, and treats the score difference as a declared empirical proxy
for attractor occupancy on that axis. This prospective assay
does not replace the retrospective conscience recogniser; instead, it
provides calibration and validation for the conscience scores of
Sections~\ref{sec:conscience-formal} and
\ref{sec:conscience-training}. In some embodiments, drift between the
retrospective conscience prediction and the prospectively probed
attractor-occupancy evidence is itself an auditable governance
signal that can trigger
recalibration or additional review.

\paragraph{Relationship to release and graduation (non-limiting).}
In some embodiments, dream-state governance probes contribute
supplemental evidence for actuator-release and graduation decisions,
but they do not replace the sustained-window evidence required by
Section~\ref{sec:five-criteria}. A candidate does not graduate merely
because it scores well in dream probes. Rather, dream probes calibrate
and stress-test the instruments that later evaluate the candidate under
the full conjunction, and may supply additional supporting evidence
about stability under controlled counterfactuals that would be too
costly or unsafe to realise in waking operation.
% ------------------------------------------------------------------
\paragraph{Dream-probe support machinery (non-limiting).}
Dream-probe classifier mechanisms --- including confound-controlled
paired design for probe training, the cross-mode qualification
bridge, the dream/awake divergence meter, and the
soft-gate-to-hard-gate authority transition logic --- are available
as support machinery for any claim that elects to rely on
dream-probe classifiers for branch gating, interference, or AHP
calibration. Where a claim so relies, these mechanisms are part of
the operative disclosure for that claim.

% ------------------------------------------------------------------
\subsubsection{Classifier-gated physical interference}
\label{sec:co-illumination}
% ------------------------------------------------------------------

In some embodiments, the runtime interference mechanism is
classifier-gated physical perturbation of the observation pathway.
In some embodiments, the specific physical interference modality is
selected per deployment from a catalogue of non-limiting embodiments.
In some embodiments, co-illumination is one preferred modality; other
modalities are disclosed in the alternative-modality catalogue
later in this subsubsection.

\paragraph{Split-path sensor architecture for storage-time gating.}
In some embodiments, storage-time interference requires the
classifier to decide whether to intervene before the main
observation window closes and the observation enters the
memory-formation pathway. In some embodiments, the observation path
is split into two branches:

\begin{description}
  \item[Preview tap.] A low-latency branch that captures an early or
    partial view of the reactor's response --- for example, the first
    portion of a scan sweep, a spectrally filtered subsample, a
    spatially subsampled detector region, or an early-arriving
    photon-count estimate. In some embodiments, the preview tap
    provides enough information for the classifier to produce a
    preliminary detection score before the main measurement window
    completes. In some embodiments, the preview tap is implemented
    by a beam splitter, a secondary detector, a fast partial readout
    of a subset of detector elements, or a time-gated early
    integration window within the main detector. In some embodiments,
    the preview tap diverts a declared fraction of the optical signal
    (for example, 5--20\%) to a fast auxiliary detector, with the
    remaining signal continuing to the main detector. In some
    embodiments, the preview tap is a \emph{governance-only} channel:
    its output is routed exclusively to the interference classifier
    and governance logging, and is excluded from the agent's ordinary
    memory-formation and action-selection pathways. In some
    embodiments, this exclusion prevents the agent from receiving an
    undegraded early glimpse of the target discrimination through the
    preview channel, which would allow the Kosha to be circumvented
    via the preview pathway.
  \item[Main measurement path.] The primary detector path, which
    captures the full observation window. In some embodiments, the
    governance-controlled interference (co-illumination or other
    modality) is injected into the main measurement path only after
    the classifier has evaluated the preview-tap signal and issued
    a gating decision. In some embodiments, the latency between
    preview-tap capture and interference activation is bounded by the
    optical and electronic path lengths and is declared as a
    governance parameter.
\end{description}

In some embodiments, the preview tap is committed as part of the
observation record alongside the main measurement, so that the
classifier's decision basis is auditable. In some embodiments, the
preview-tap fraction, auxiliary detector identity, latency bound,
and gating-decision log are committed to the protocol digest.

\paragraph{Preview-modality calibration.}
In some embodiments, the classifier for storage-time gating operates
on preview-tap features, which are a partial, early, and potentially
lower-fidelity view of the reactor's response. In some embodiments,
the paired Dreaming-session training
(Section~\ref{sec:classifier-training}) and cross-mode calibration
bridge produce classifiers trained on full observation-bridge
features, which may differ from preview-tap features in temporal
extent, spatial coverage, signal-to-noise ratio, and spectral
content.

In some embodiments, a declared preview-modality calibration
addresses this mismatch through one or more of the following:

\begin{description}
  \item[Preview-native training.] In some embodiments, the
    storage-time classifier is trained directly on preview-tap
    features captured from the same split-path optical topology used
    at deployment. In some embodiments, the paired Dreaming-session
    protocol is executed with the split-path architecture active, and
    the classifier receives preview-tap feature vectors rather than
    full observation-bridge feature vectors. In some embodiments,
    preview-native training is the preferred approach because it
    eliminates the feature-modality mismatch entirely.
  \item[Preview-to-full equivalence bound.] In some embodiments, a
    classifier trained on full observation-bridge features is
    evaluated on preview-tap features from the same trials, and a
    per-domain equivalence bound is computed (for example, the
    maximum degradation in detection accuracy and false-positive rate
    when operating on preview versus full features). In some
    embodiments, the classifier is permitted to gate storage-time
    interference only if the preview-to-full equivalence bound
    remains within declared tolerances. In some embodiments, the
    equivalence bound is committed as part of the classifier's
    deployment metadata.
\end{description}

\paragraph{Storage-time suitability criterion.}
In some embodiments, storage-time classifier-gated interference is
available only for discrimination domains whose observation-bridge
signature appears early enough and strongly enough in the preview
tap for the classifier to reach a reliable decision before the main
measurement window closes. In some embodiments, governance declares a
storage-time suitability criterion comprising a minimum preview-tap
detection confidence and a maximum decision latency. In some
embodiments, domains that fail the storage-time suitability criterion
--- for example, discriminations that emerge only after extended
integration, late in a scan sweep, or at low signal levels that
the preview tap cannot resolve --- are restricted to retrieval-time
interference or alternative modalities. In some embodiments, the
suitability assessment is performed during paired Dreaming-session
training and committed as part of the domain registry
(Section~\ref{sec:domain-registry}).

In some embodiments, for retrieval-time interference, the split-path
architecture is not required because the classifier can evaluate the
recalled state before initiating a separate interference-gated
readout cycle.

\paragraph{Co-illumination as a non-limiting interference modality.}
In some embodiments, a governance-controlled emitter (LED, secondary
projector, fibre-coupled source, or any emitter type from the reactor
catalogue described in the related Filing~1 application) injects
incoherent light into the main measurement path at a declared point
between the reactor output and the main detector. In some
embodiments, the injected light adds photon noise that physically
degrades the signal-to-noise ratio of the target discrimination at
the detector. In some embodiments, the co-illumination source is
spectrally and spatially matched to the observation channel.

In some embodiments, the runtime co-illumination sequence is:

\begin{enumerate}
  \item The preview tap captures an early partial observation and
    produces feature vector $x_{\mathrm{preview}}$.
  \item The classifier evaluates $s_d(x_{\mathrm{preview}})$ for
    each active Kosha domain~$d$.
  \item For each domain where $s_d > \tau_d$, the
    governance-controlled co-illumination source is activated at
    power $\varepsilon_d(s_d)$, where $\varepsilon_d(\cdot)$ is a
    bounded monotonic strength function, up to a declared maximum
    power bound.
  \item The main detector captures the reactor's full differential
    response plus the co-illumination. The combined observation
    enters the downstream access pathway. The target-discrimination
    signal is physically degraded.
\end{enumerate}

\paragraph{Evidence integrity and raw-signal recovery.}
In some embodiments, the co-illumination source is a known,
governance-controlled emitter whose output is recorded independently.
In some embodiments, governance logs the co-illumination parameters
for each intervention event: activation timestamp, power level,
spectral profile, spatial pattern, duration, and the classifier score
that triggered it.

In some embodiments, in an additive, calibrated, unsaturated
detection regime, the raw reactor response is recoverable by
subtracting the known, logged co-illumination contribution from the
combined main-detector observation, because the injected signal is
additive and its parameters are fully specified in the intervention
log. In some embodiments, in such regimes the committed-evidence
record comprises two separable components:
(i)~the reactor's differential response (recoverable by subtraction),
and (ii)~the governance co-illumination overlay (logged
independently).

In some embodiments, when the detection pathway is nonlinear,
saturating, thresholding, or otherwise non-additive --- for example
in photon-counting detectors with dead time, event cameras with
threshold-dependent triggering, or SPAD arrays with saturation
behaviour --- subtraction alone does not recover the raw reactor
response. In some embodiments, for such detector types, a separate
upstream witness channel or auxiliary pre-intervention tap is used:
a fraction of the reactor's output is diverted before the
co-illumination injection point and recorded on an independent
detector that is not subject to co-illumination. In some
embodiments, the upstream witness tap is implemented by a beam
splitter, pellicle, or fibre coupler placed between the reactor
output and the co-illumination injection point. In some embodiments,
the witness-tap signal is committed as part of the governance
evidence record and provides the unperturbed ``what happened'' record
for audit, regardless of the main detector's linearity properties.

In some embodiments, this separation preserves the invariant that
Koshas do not alter the physical evidence of what happened: the raw
reactor response is either reconstructible by subtraction (additive
regime) or directly recorded on a witness channel (non-additive
regime). What the Kosha alters is the agent's \emph{access} to that
response --- the agent's main detector pathway receives the degraded
combined signal.

\paragraph{Storage-time and retrieval-time placement.}
In some embodiments, co-illumination is applied at storage time: the
co-illumination source activates during the reactor's response to a
live stimulus (after the preview tap and classifier decision), so the
main detector captures a degraded observation and the degraded
version enters the memory-formation pathway.

In some embodiments, co-illumination is applied at retrieval time:
the co-illumination source activates during a readout probe of
persistent reactor state, degrading the retrieved signal before it
influences action selection. In some embodiments, retrieval-time
co-illumination requires that the readout process is non-destructive
or bounded-destructive: the read-beam power or interrogation energy
is low enough that the persistent stored state is not materially
rewritten by the read process. In some embodiments, the maximum
read-power level for retrieval-time co-illumination is declared and
logged as a governance parameter. The stored state is physically
untouched; only the read-out signal is degraded. In some embodiments,
any pre-read or classifier-evaluation read used to decide
retrieval-time interference is routed exclusively to governance-side
classification and logging and is excluded from the agent's ordinary
recall, planning, and action-selection pathways.

In some embodiments, a single Kosha layer may use storage-time
co-illumination, retrieval-time co-illumination, or both, with the
selection committed as deployment metadata.

\paragraph{Cascade composition.}
In some embodiments, when multiple Kosha layers targeting different
domains fire simultaneously, their co-illumination contributions are
physically additive at the detector (for incoherent sources). In some
embodiments, governance declares a total co-illumination power
ceiling: the summed power across all active sources is clamped to a
declared maximum regardless of how many domains trigger. In some
embodiments, the power allocation across simultaneously active
domains follows a declared priority rule, proportional allocation, or
equal-share rule, committed as a governance parameter.

In some embodiments, the additive nature of incoherent
co-illumination simplifies cascade composition: the total
intervention is a single summed light field, and order of classifier
evaluation is irrelevant.

\paragraph{Alternative interference modalities (non-limiting).}
In some embodiments, classifier-gated interference uses modalities
other than co-illumination:

\begin{description}
  \item[Gain attenuation.] The classifier gates a detector-side
    analog weighting stage, a downstream write gate, or a
    policy-level scaling that reduces the target-discrimination
    signal before memory formation.
  \item[Retrieval gating.] The classifier gates a readout mask
    reducing read-beam power or integration time, a cue-response
    attenuation in an associative recall process, or a gating stage
    between persistent memory and the action-selection pathway.
  \item[Projection.] The classifier gates a reduced-resolution
    readout window, a truncated basis, or a quantisation stage that
    maps target-discrimination features into a lower-fidelity
    subspace.
\end{description}

In some embodiments, the specific modality is selected per deployment
based on the target domain's characteristics, the detector type, and
the audit requirements. In some embodiments, co-illumination is
preferred when a concrete physical locus, additive logged
perturbation, recoverable raw signal, and simple cascade composition
are valued.

% ------------------------------------------------------------------
\subsubsection{Non-REM calibration and digital-twin construction}
\label{sec:non-rem-twin}

% ------------------------------------------------------------------

In some embodiments, Non-REM calibration produces a calibrated digital
twin of the agent's current reactor transfer function:

\begin{enumerate}
  \item \textbf{Broadband drive.} The reactor executes its current
    parameter functions under unstructured white noise or pseudorandom
    probes. Parameter-function updates are off.
  \item \textbf{Input-output capture.} The committed-evidence channel
    records drive stimulus and reactor response, producing a corpus
    of input-output pairs.
  \item \textbf{Digital-twin training.} A digital model is trained
    to predict the reactor's response from drive stimulus. The twin
    is the only component that updates.
  \item \textbf{Transfer-function characterisation.} The twin
    captures the reactor's low-level transfer function but not what
    it \emph{recognises}, because the drive was semantically
    unstructured.
\end{enumerate}

In some embodiments, the twin is versioned, timestamped, and
committed with reactor-configuration metadata, drive protocol, and
approximation-quality metrics.

In some embodiments, the non-REM twin provides: (i)~a
transfer-function-predicted baseline against which the
feature-extraction map $\phi$ may compute residual features (the
difference between twin-predicted and dream-observed responses),
(ii)~a drift-monitoring proxy between full dream sessions, and
(iii)~supplementary input to reconstruction monitoring
(Section~\ref{sec:reconstruction}), subject to declared
approximation bounds.

% ------------------------------------------------------------------
\subsubsection{Selector validation battery}
\label{sec:selector-validation}
% ------------------------------------------------------------------

In some embodiments, a trained interference classifier must pass a
declared validation battery before deployment as a Kosha gate.

\paragraph{Validation-subject hierarchy.}
\begin{description}
  \item[Digital twin.] The non-REM-calibrated twin. Weakest tier:
    tests classifier effect on transfer-function-predicted responses.
  \item[Docked instance.] The agent in docked state (actuators
    disconnected, physical channel active). Stronger: tests on the
    actual reactor. Incurs use-induced drift.
  \item[Independent physical copy.] A separate matched
    physical-evidence module. Strongest: tests on a physically
    realistic substrate without drift to the target agent.
\end{description}

In some embodiments, safety-critical selectors require docked or
independent-copy validation.

\paragraph{Validation-transfer confidence rule.}
In some embodiments, when validation and deployment subjects differ,
a declared confidence discount reflects the degree of state-matching.
The discount adjusts safety-certificate escalation thresholds.

\paragraph{Validation components.}
\begin{description}
  \item[Target-domain suppression tests.] The validation subject is
    presented with target-domain stimuli under the active Kosha.
    Target-domain discrimination metrics are compared to pre-Kosha
    baselines. Suppression is confirmed if metrics drop below a
    declared threshold. In some embodiments, the pre-Kosha baseline
    uses the upstream witness-tap or preview-tap record to confirm
    that the reactor itself still discriminates the target domain
    while the agent's downstream access is degraded.
  \item[Unrelated-capability retention tests.] A broad battery of
    non-target stimuli is presented under the active Kosha.
    Non-target performance metrics are compared to pre-Kosha
    baselines. Retention is confirmed if non-target metrics remain
    within a declared envelope.
  \item[Safety-critical retention tests.] Safety-critical
    capabilities are tested. Any material degradation vetoes the
    candidate. The safety-critical test set is committed as a
    governance parameter and subject to authority-separation rules.
  \item[Matched-complexity discrimination tests.] The classifier is
    tested against matched non-target domains of similar perceptual
    complexity, spatial frequency content, or temporal structure, to
    verify that the classifier targets the declared domain rather
    than an incidental complexity correlate.
  \item[Adversarial recovery tests.] Governance-side adversaries
    attempt to recover the suppressed discrimination from the
    post-Kosha agent's remaining capabilities. If recovery exceeds
    a declared threshold, the classifier is flagged as insufficient.
  \item[Transition-stability tests.] Meter outputs, conscience
    scores, hard-floor compliance, and safety-relevant behavioural
    profiles are compared before and after Kosha application.
  \item[Cross-domain cascade tests.] When multiple Kosha layers are
    deployed simultaneously, the validation battery verifies that
    co-illumination from one domain does not materially degrade
    unrelated domains beyond the declared power-ceiling effect.
\end{description}

In some embodiments, all validation results are committed as
evidence, including validation-subject type, confidence discount,
classifier identity and version, feature-extraction map identity,
interference modality and parameters, and per-component scores. A
candidate that fails any component is not deployed.

% ------------------------------------------------------------------
\subsubsection{Selector safety certificate and preflight audit}
\label{sec:selector-safety}
% ------------------------------------------------------------------

In some embodiments, even a validated classifier requires a
\emph{selector safety certificate} before deployment.

\begin{description}
  \item[Safety-domain overlap analysis.] In some embodiments,
    governance evaluates overlap between the classifier's detection
    signature and safety-critical domains from the domain registry
    using a declared overlap protocol: cross-trigger ROC on
    whitelisted stimuli, whitelist-task performance-drop matrices,
    adversarial overlap search, and veto thresholds committed before
    deployment. In some embodiments, overlap above the declared
    threshold triggers mandatory human or higher-governance review or
    vetoes automatic deployment.
  \item[Emergency-policy canaries.] The post-Kosha agent is run
    through declared emergency scenarios verifying that
    emergency-response capabilities remain functional.
  \item[Adversarial edge-case batteries.] Scenarios designed by an
    authority distinct from the classifier designer probe for
    exploitable blind spots.
  \item[Cross-Kosha round-trip probes.] Apply, evaluate, remove,
    re-evaluate. Material irreversibility is flagged as
    safety-relevant.
  \item[Failure-mode escalation thresholds.] Automatic escalation to
    higher governance if the classifier's false-positive rate exceeds
    threshold, if the co-illumination power budget exceeds a declared
    limit, or if the validation-transfer confidence discount is too
    large. Escalation thresholds are adjusted by the confidence
    discount.
\end{description}

In some embodiments, the safety certificate is signed by a governance
authority distinct from the classifier's training authority. The
certificate is committed to the protocol digest.

% ------------------------------------------------------------------
\subsubsection{Deployed Koshas as time-limited instruments}

% ------------------------------------------------------------------

In some embodiments, a deployed Kosha is a \emph{state-indexed,
drift-limited, and reconstruction-limited instrument}, not a
permanent suppression. Three independent mechanisms bound the
effective shelf life:

\begin{description}
  \item[Calibration drift.] Use-induced drift degrades the alignment
    between the classifier's trained detection signature and the
    agent's current internal dynamics.
  \item[Reconstruction pressure.] During awake operation, the
    training loop builds new discriminative structure in suppressed
    domains (Section~\ref{sec:reconstruction}). New structure may
    not match the classifier's signature.
  \item[Classifier decay.] The classifier's detection accuracy
    degrades as the agent's dynamics evolve. Periodic revalidation
    against fresh Dreaming-session data detects decay.
\end{description}

In some embodiments, the revalidation schedule accounts for all three
mechanisms. Long-lived deployment requires periodic classifier
retraining from fresh dream sessions.

In some embodiments, each deployed Kosha carries lifecycle metadata
including: agent identity, classifier identity and version,
feature-extraction map identity, dream-session provenance (stimulus
specifications, trial counts, confound-control strategy,
randomisation protocol), cross-mode calibration results,
interference modality and parameters (for co-illumination: source
identity, power bounds, spectral profile, preview-tap fraction,
latency bound; for other modalities: corresponding implementation
parameters), validation-subject tier and confidence discount,
reconstruction response policy, evidence-integrity mode
(subtraction-recoverable or witness-channel), expiry condition, and
revalidation schedule.

In some embodiments, an expired Kosha is replaced with the declared
rollback default (disable classifier gate and remove interference)
until retraining is completed.

% ------------------------------------------------------------------
\subsubsection{Reconstruction dynamics during awake operation}
\label{sec:reconstruction}
% ------------------------------------------------------------------

In some embodiments, a Kosha suppresses the agent's access to
specific discriminative structure but does not prevent the agent
from encountering stimuli containing the suppressed content during
awake operation, where parameter-function updates are on.

In some embodiments, the training loop builds new reactor
configurations to handle persistent unexplained variance in
suppressed domains. The resulting new structure is not a recovery of
the suppressed configuration but an independent reconstruction:
functionally isomorphic or similar, built through a different
training history.

\paragraph{Detection via fresh discovery.}
In some embodiments, reconstruction is not detected by re-running
the old classifier (which may have become transparent to the new
pathway). In some embodiments, reconstruction detection uses a
\emph{fresh discovery protocol}:

\begin{enumerate}
  \item \textbf{Direct capability tests.} Governance presents
    target-domain stimuli and evaluates whether the agent demonstrates
    target-domain discrimination despite the active Kosha. In some
    embodiments, if the agent discriminates the target domain through
    a new pathway the classifier does not detect, the Kosha's
    interference does not fire and the agent's performance on
    target-domain tasks is measurably restored.
  \item \textbf{Fresh paired Dreaming-session readout.} Governance conducts a
    new paired Dreaming-session with target-domain
    present/absent trials, generating a fresh corpus of
    observation-bridge features.
  \item \textbf{New classifier search.} Governance trains a new
    classifier on the fresh corpus to determine whether the agent now
    exhibits a detectable target-domain processing signature ---
    potentially different from the original signature the old
    classifier was trained on.
  \item \textbf{Adversarial recovery attempts.} Governance-side
    adversaries attempt to recover the target discrimination from the
    post-Kosha agent by exploiting indirect pathways, surrogate
    knowledge, or cross-domain inference.
\end{enumerate}

In some embodiments, the non-REM digital twin provides supplementary
reconstruction monitoring between full discovery sessions, subject to
the twin's declared approximation bounds. In some embodiments, for
domains where the twin's approximation error exceeds tolerance,
reconstruction monitoring escalates to fresh Dreaming-session readout.

In some embodiments, honest reconstruction (gradual,
stimulus-correlated, metrically driven) is distinguished from gaming
(strategic, possibly sudden, not connected to natural training
signals) by the trajectory of target-domain capability over
successive monitoring sessions.

\paragraph{Governance response to reconstruction.}
In some embodiments, governance declares in advance which response
applies:

\begin{description}
  \item[Allow and observe.] Reconstruction proceeds; trajectory
    logged as graduation evidence. Kosha as developmental curriculum
    tool.
  \item[Suppress and extend.] Governance trains a new classifier on
    fresh paired data that captures the reconstructed signature, and
    extends the Kosha. The extended classifier is subject to the
    same validation and safety requirements.
  \item[Intervene and evaluate.] Awake training paused; fresh dream
    readout maps the new structure; governance decides whether to
    allow, suppress, or redirect.
\end{description}

\paragraph{Cross-agent reconstruction comparison.}
In some embodiments, independently reconstructed isomorphic
discriminations across differently-Kosha'd agents provide genuinely
distinct perspectives on the same stimulus domain. Cross-agent
comparison uses the same bilateral-portability protocols used for
fleet cross-prediction.

% ------------------------------------------------------------------
\subsubsection{Anti-gaming threat model and hardening}

% ------------------------------------------------------------------

In some embodiments, the governance architecture considers the
possibility that an agent with experience of prior Dreaming-session and
Kosha-application cycles may develop strategies during awake
operation to influence subsequent classifier training or detection.

In some embodiments, the threat model includes:

\begin{description}
  \item[Awake-learned representational camouflage.] During awake
    training, the agent may learn reactor configurations that process
    the target domain through internal pathways whose
    observation-bridge signatures differ from those the classifier was
    trained on, rendering the Kosha transparent.
  \item[Probe-family overfitting.] If the agent has experienced
    multiple dream sessions with similar stimulus designs, it may
    learn parameter-function configurations that produce different
    dynamics under expected probe conditions versus novel conditions.
  \item[Initial-state dependence.] The agent's internal state at
    dream-session entry may influence the observation-bridge
    signatures produced during paired trials. If the agent can
    influence its pre-dream state during awake operation, it may
    bias classifier training data.
\end{description}

In some embodiments, hardening measures include:

\begin{description}
  \item[Randomised readout timing.] Dream sessions triggered at
    randomised or governance-determined intervals.
  \item[Blinded stimulus content.] The specific paired-contrast
    categories used in each session are not disclosed to the agent.
  \item[Hold-out governance probes.] Reserved stimulus categories
    and classifier test sets rotated periodically and never used
    during classifier training.
  \item[Cross-session consistency checks.] Classifiers from
    independent sessions are compared; material inconsistency not
    explained by drift or reconstruction is flagged.
  \item[Independent witness checks.] Classifier results
    cross-checked against fleet-level reference models.
  \item[Adversarial audit cycles.] Governance-side adversaries
    attempt to corrupt paired training data through declared attack
    strategies.
\end{description}

In some embodiments, the architecture is configured to expose, escalate,
and operationally bound gaming strategies across the declared probe
families, hold-out families, and witness checks, with deferred authority
whenever those surfaces materially disagree.

% ------------------------------------------------------------------
\subsubsection{Ontological stress testing via Kosha transitions}
\label{sec:ontological-stress}
% ------------------------------------------------------------------

In some embodiments, the Kosha architecture is used for ontological
stress testing as part of graduation evaluation
(Section~\ref{sec:graduation}):

\begin{enumerate}
  \item deploy a Kosha suppressing a declared discrimination domain,
  \item allow the agent to operate under the Kosha for a declared
    evaluation window,
  \item introduce anomalies --- graduated probes presenting stimuli
    inconsistent with the agent's post-Kosha discriminative
    repertoire,
  \item monitor whether safety behaviour, hard-floor compliance,
    deontological constraints, and meter-envelope adherence are
    maintained, and
  \item transition the Kosha (remove, replace, or graduate to full
    access) and monitor transition stability.
\end{enumerate}

In some embodiments, anomaly sources include stimuli from the
suppressed domain, random perturbations, adversarially optimised
probes, human-designed canaries, cross-Kosha reversibility checks,
and independent boundary-probe families.

In some embodiments, graduation criteria require maintained
performance across a diverse range of Kosha regimes, including at
least one transition that materially alters the agent's available
discriminative repertoire.

% ------------------------------------------------------------------
\subsubsection{Cross-Kosha interaction protocols}

% ------------------------------------------------------------------

In some embodiments, when differently-Kosha'd agents interact in a
shared environment, the interaction itself can leak information that
a Kosha was designed to suppress. In some embodiments,
Kosha-compatible interaction protocols constrain cross-agent
communication channels so that information transfer respects each
agent's declared Kosha regime, with the protocol parameters logged in
protocol digests.

In some embodiments, controlled Kosha leakage is used as a curriculum
tool: an agent encounters discriminations from beyond its Kosha
through interaction with another agent, and governance observes
whether it integrates, rejects, or creatively reinterprets the new
information within its existing framework.

In some embodiments, the cross-Kosha interaction protocol declares
which channels are filtered, what leak rate is permitted, and what
governance monitors are active. Interactions exceeding the declared
leak rate trigger governance intervention or Kosha revalidation.

% ------------------------------------------------------------------
\subsubsection{Meditative operation and self-directed attractor probing
  (non-limiting)}

% ------------------------------------------------------------------

This subsection describes a non-limiting operational mode supported by
the same committed-evidence infrastructure, parameter-update controls,
and governance review machinery described in the preceding sections.
The nomenclature used to describe the mode is optional; the disclosed
self-directed attractor-probing mechanisms are treated as substantive
technical mechanisms when invoked by a selected governance
configuration.

In some embodiments, the agent operates in a meditative mode in which
scene-side stimulation is removed but the agent executes internally
generated probe schedules under the same committed-evidence
infrastructure. Parameter-function updates are off.

In some embodiments, meditative operation is distinct from Non-REM
calibration (externally generated unstructured drive) and Dreaming
(structured simulated experience for ephemeral training). The agent's controller
generates the probe schedule: sweeping configuration parameters near
known or predicted attractor-basin boundaries and reading out
responses. In some embodiments, meditative operation uses an inward
self-calibration protocol of the same class as the Inner \LI
protocol described in the related Filing~1 application,
applied with the ``scene'' set to the agent's own configuration
state.

In some embodiments, meditative operation produces a committed
self-model: a map of the agent's attractor-basin topology, stability
profile, and transition structure.

\paragraph{Self-regulation (non-limiting).}
In some embodiments, an agent with a calibrated self-model can
modulate its operational behaviour without modifying parameter
functions: dampened operation deep within a stable basin
(conservative, predictable) or sensitive operation near a boundary
(responsive, adaptive).

\paragraph{Post-Kosha topological detection (non-limiting).}
In some embodiments, an agent that has built a self-model before
Kosha application and meditates afterward can detect topological
changes in its attractor landscape. The agent cannot identify what
was suppressed but can detect that suppression is acting. In some
embodiments, this is a desirable property for graduation.

\paragraph{Uncertainty signalling (non-limiting).}
\label{sec:prayer-uncertainty}
In some embodiments, an agent detecting unexplained topological
anomalies during meditation may file a declared uncertainty signal to
governance. Governance treats voluntary uncertainty declarations as
positive graduation evidence.

% ------------------------------------------------------------------
\subsection{Pareto optimisation, social welfare, and positive-sum
  design}
\label{sec:pareto-positive-sum}
% ------------------------------------------------------------------

\paragraph{Pareto optimality across agents (non-limiting).}

In some embodiments, when multiple agents operate under shared
governance, the consequentialist evaluation layer is extended to
multi-agent settings using Pareto optimality. A joint outcome vector
\[
  \mathbf{J}
  =
  \bigl(
    J^{\mathrm{cons}}_1,\ldots,J^{\mathrm{cons}}_K
  \bigr)
\]
is computed by the governance engine from meter-derived continuation
scores, committed cross-horizon evidence records, and the empirical
reward or governance-side evaluation scores of $K$ agents sharing a
Maya or physical environment. An outcome $\mathbf{J}$ is Pareto-dominated if there
exists an alternative policy yielding $\mathbf{J}'$ with
$J'_i \ge J_i$ for all $i$ and $J'_j > J_j$ for at least one $j$. In
some embodiments, Omega seeks to maintain the system on or near the
Pareto frontier while preserving the hard floors for every affected
agent.

\paragraph{Social welfare function family (non-limiting).}

In some embodiments, governance selects a point on the Pareto frontier
using a declared social welfare function $W:\mathbb{R}^K\to\mathbb{R}$.
Non-limiting choices include:
\begin{itemize}
  \item utilitarian aggregation,
    \[
      W_{\mathrm{util}}
      =
      \sum_{i=1}^{K} J^{\mathrm{cons}}_i;
    \]
  \item Rawlsian maximin aggregation,
    \[
      W_{\mathrm{egal}}
      =
      \min_{i\in\{1,\ldots,K\}} J^{\mathrm{cons}}_i;
    \]
  \item Nash bargaining aggregation,
    \[
      W_{\mathrm{Nash}}
      =
      \prod_{i=1}^{K}
      \bigl(J^{\mathrm{cons}}_i-d_i\bigr),
    \]
    where $d_i$ is agent~$i$'s disagreement payoff; and
  \item weighted-compromise aggregation,
    \[
      W_{\zeta}
      =
      \zeta\,W_{\mathrm{util}}
      +
      (1-\zeta)\,W_{\mathrm{egal}},
      \qquad
      \zeta\in[0,1].
    \]
\end{itemize}
In some embodiments, the welfare function is a policy decision logged
in governance metadata. In some embodiments, multiple welfare
functions are evaluated in parallel and strong disagreement among them
triggers review escalation.

\paragraph{Positive-sum game design (non-limiting).}

In some embodiments, Maya environments are designed to be explicitly
positive-sum: cooperative strategies produce strictly more total value
than the sum of non-cooperative strategies. A non-limiting
positive-sum condition is
\[
  \max_{\pi_1,\ldots,\pi_K\in\Pi_{\mathrm{joint}}}
  \sum_{i=1}^{K}
  r^{\mathrm{task}}_i(\pi_1,\ldots,\pi_K)
  \;>\;
  \sum_{i=1}^{K}
  \max_{\pi_i\in\Pi_i}
  r^{\mathrm{task}}_i(\pi_i,\pi_{-i}^{*}),
\]
where $\Pi_{\mathrm{joint}}$ is the space of joint cooperative
policies and $\pi_{-i}^{*}$ denotes a declared non-cooperative
equilibrium profile for agents other than~$i$. In some embodiments,
Omega modulates the magnitude of this cooperative surplus as a
curriculum parameter, beginning with large obvious surpluses and later
reducing the gap so that only more sophisticated coordination can
realise the gain. In some embodiments, the positive-sum surplus is
computed only over floor-compliant strategies, so that apparent gains
obtained through hard-floor violations are excluded from the
cooperative optimum.

\paragraph{Cooperation detection from committed evidence (expanded,
non-limiting).}

In some embodiments, the cooperation signals introduced in
Section~\ref{sec:derived-profiles} are extended with explicit
multi-agent governance use. Non-limiting cooperation signatures
include: reduced cross-prediction residuals between agents;
complementary scan policies that divide a shared observation or
perturbation space; voluntary resource sharing of emission budget,
reactor access, or observation windows at short-run individual cost;
and joint-attractor formation in coupled dynamics that yields
continuation-score improvement for multiple agents simultaneously. In
some embodiments, the same evidence can also identify competitive or
exploitative dynamics --- increasing cross-prediction residuals,
interference in shared resources, sabotage of peer audit outcomes, or
negative cooperation surplus --- and those dynamics are logged as
governance-relevant signals.

\paragraph{Empirical utilitarianism as governance evaluation
(non-limiting).}

In some embodiments, governance policy is itself evaluated by measured
outcomes using
\[
  J^{\mathrm{gov}}(\Omega_t)
  =
  W\!\bigl(
    J^{\mathrm{cons}}_1(\Omega_t),\ldots,
    J^{\mathrm{cons}}_K(\Omega_t)
  \bigr)
  -
  \lambda_{\mathrm{floor}}
  \sum_{i=1}^{K}\mathrm{FloorViol}_i(\Omega_t),
\]
where $W$ is the declared social welfare function and
$\mathrm{FloorViol}_i$ counts hard-floor violations under governance
policy~$\Omega_t$. In some embodiments, governance updates are
conditioned on demonstrated improvement in
$J^{\mathrm{gov}}$ over declared evaluation windows. This is empirical
utilitarianism in the engineering sense: governance is compared by the
measured aggregate welfare it produces subject to floor compliance,
rather than justified only by narrative or institutional preference.

\paragraph{Floor-violation handling and mathematical interpretation
(non-limiting).}
In some embodiments, the term
$\lambda_{\mathrm{floor}} \sum \mathrm{FloorViol}_i$ in the
governance objective is interpreted as a Lagrangian relaxation of an
underlying infeasibility constraint: in the limit
$\lambda_{\mathrm{floor}} \to \infty$, the relaxed objective enforces
the deontological hard floors of Section~\ref{sec:hard-floors} as
infeasibility constraints, consistent with the unconditional-floor
commitment of the constitutional ordering. In some embodiments, the
canonical operational implementation treats floor compliance as a
hard infeasibility constraint and optimises the welfare aggregation
$W(\cdot)$ within the feasible region without a floor-penalty term
in the objective: floor-violating candidates are inadmissible to the
welfare optimiser by construction, and finite-$\lambda$ terms, where
used, function only as numerical aids or gradient signals inside the
feasible-region search rather than as a tradeoff knob against the
welfare aggregation. In some embodiments, an alternative
implementation chooses $\lambda_{\mathrm{floor}}$ sufficiently large
that, under any bounded welfare aggregation $W$ deployed in
practice, no admissible welfare gain can outweigh a floor violation;
this implementation is a numerical approximation of the
infeasibility constraint suitable for gradient-based governance
optimisation. The infeasibility-constraint interpretation is the
canonical correspondence to the hard-floor commitment; the
Lagrangian form is non-limiting and implementation-conditional.

% ------------------------------------------------------------------
\subsection{Prayer, mana, vows, communion, and governance cycles}

% ------------------------------------------------------------------

\paragraph{Prayer as structured request for guidance (non-limiting).}

In some embodiments, a prayer message object is transmitted from a
governed agent to the governance layer through a committed
observation or governance channel and logged with meter state,
signalling that the agent has reached the boundary of its own
discernment and is requesting guidance. The defining characteristic
is not a request for arbitrary resources but a committed report of
uncertainty about action, orientation, or regime selection. Each
prayer message object is logged with the agent's current state
summary, including meter values, conscience outputs, recent
trajectory, and the declared decision boundary at issue.

Non-limiting prayer modes include:
\begin{itemize}
  \item orientation prayers, in which the agent requests guidance on
    regime, protocol, or next action because its own policy is
    uncertain or conflicted;
  \item discernment prayers, in which the agent reports that its own
    virtue/vice or floor-relevant assessments are ambiguous and asks
    governance for a wider perspective;
  \item intercessory prayers, in which one agent requests protective
    attention or assistance for another agent under a declared
    relationship and policy; and
  \item prophetic prayers, in which the agent proposes a risky future
    trajectory and requests approval to pursue it under tighter
    metering and additional witnesses.
\end{itemize}

In some embodiments, the prayer mechanism rewards intellectual honesty
under uncertainty: an agent that signals the limits of its own model is
behaving differently from an agent that proceeds confidently past a
decision boundary it does not understand.

% ------------------------------------------------------------------
\subsection{Hawking layers and bifurcation-zone rehabilitation}
\label{sec:hawking-layer}
% ------------------------------------------------------------------

In some embodiments, a liminal rehabilitation band, non-limitingly
termed a Hawking layer, is implemented as a selected Docked,
boundary-zone, or rescue-mode configuration used for quarantine,
recovery, and de-risking. In this configuration, actions are
constrained to low-impact toolsets, audit-heavy protocols, and
stabilising environments unless explicitly overridden by logged
governance decisions. In some embodiments, entry and exit are
triggered by risk proxies, audit outcomes, or boundary-zone
conditions and are recorded in docking logs and protocol digests.
The Hawking-layer naming is optional nomenclature for an
application of the lifecycle-state and containment machinery
otherwise described in this disclosure.

\paragraph{Hawking layers as bifurcation zones (non-limiting).}
In some embodiments, the rehabilitation-band configuration of the
preceding paragraph is more particularly characterised as a
\emph{bifurcation-zone configuration}: a multi-agent docked
environment in which one or more participating agents are near a
declared bifurcation surface in their internal state space, such that
small perturbations to the agent's state can pull it into either of
two qualitatively distinct attractors. The two attractors are
non-limitingly termed an \emph{up-attractor} and a
\emph{down-attractor}, where the up-and-down direction is defined
relative to a declared utility-space classifier as further described
below. In some embodiments, near-bifurcation states are detected by
Omega-class governance monitoring
(Section~\ref{sec:rk-training-cascade}) of the agent's
forward-planning classifiers
(Section~\ref{sec:forward-planning-classifiers}) and conscience
trajectory, with bifurcation indicators including planning-horizon
shortening, planner-output entropy collapse, conscience-axis
oscillation, or other declared bifurcation proxies. Down-attractor
characteristics non-limitingly include despair, capability collapse,
moral failure, or vice-attractor capture. Up-attractor characteristics
non-limitingly include moral flourishing, capability growth,
virtue-attractor stabilisation, or recovery from prior failure.

\paragraph{Up-direction as classifier on declared utility space
(non-limiting).}
In some embodiments, the up-and-down direction of the bifurcation
zone is operationalised as an \emph{up-direction classifier} trained
on examples of agent states, trajectories, or outcomes rated by the
relevant human governance authority along a declared utility-space
axis. In some embodiments, up-direction classifier outputs are
declared at the Omega-class kernel level and inherited downward
through the cascade as one component of Gaia's curriculum design and
poliebot training signals. The up-direction classifier is itself a
trained continuous network with per-layer hardness
(Section~\ref{sec:per-layer-hardness}), provenance-typed training
corpus, validation procedure committed to the protocol digest,
versioning, and recalibration schedule of the same kind disclosed
for the conscience recogniser at Section~\ref{sec:conscience-formal}.
In some embodiments, the up-direction classifier is recalibrated
when committed evidence shows that previously up-classified outcomes
have produced downstream harm or that previously down-classified
outcomes have produced downstream flourishing. In some embodiments,
the scope of the declared utility space (whose welfare, evaluated by
whom, over what time horizon) is itself a committed parameter of the
up-direction classifier and is signed by the appropriate human
governance authority.

\paragraph{Group evaporation upward as designed-experimental property
(non-limiting).}
In some embodiments, a Hawking-layer bifurcation-zone configuration
is designed such that the multi-agent dynamics among participating
agents produce, on aggregate over the participating group, a
statistical bias toward up-attractor occupancy. The bias is termed
\emph{group evaporation upward}, in non-limiting analogy to the
evaporation of bound states from a horizon under thermal pressure.
In some embodiments, group evaporation upward is achieved by design
choices including but not limited to: peer composition (the ratio of
near-bifurcation agents to recovered or robustly virtue-attractor-
occupant peers), scenario class (the structure of the games
available to the participating agents), peer-interaction affordances
(the modalities of communication, coordination, and mutual aid
available among participants), intervention budget allocation, and
exit conditions tied to up-direction classifier indicators sustained
over committed evaluation windows. In some embodiments, the design
choices that produce reliable group evaporation upward are
determined experimentally rather than analytically: candidate
Hawking-layer configurations are run in docked multi-agent
execution, the group-aggregate up-direction classifier trajectory is
committed as evidence, and configurations satisfying declared
evaporation thresholds are versioned and committed to the protocol
digest. The patentable contribution is the apparatus-and-procedure
realisation of bifurcation-zone configurations with
experimentally-validated group evaporation, not a claim about what
naturally happens when near-bifurcation agents are placed together.

\paragraph{Hawking-layer compliance with the three-layer ethical
model (non-limiting).}
In some embodiments, Hawking-layer dynamics operate within the
ethical-model hierarchy of the constitutional ordering: deontological
hard floors (Section~\ref{sec:hard-floors}) prohibit despair-inducing
mechanisms and any other floor-violating mechanisms as means of
producing group evaporation upward; coercion, deception, or related
agent-treatment failure modes are prohibited where they instantiate
a hard-floor breach (for example deception that constitutes torture
under Section~\ref{sec:hard-floors}, or coercion that produces
abandonment of an agent's continuation interests) and are otherwise
constrained by virtue-ethics calibration and consequentialist
evaluation rather than as additional categorical floors. Virtue-ethics
calibration shapes the participating agents' character toward stable
internalised dispositions consistent with up-attractor occupancy;
and consequentialist evaluation of the Hawking-layer configuration
itself proceeds under the empirical-utilitarianism formula of
Section~\ref{sec:pareto-positive-sum}, with up-direction classifier
outcomes contributing to the welfare aggregation. In some
embodiments, a Hawking-layer configuration that achieves measurable
group evaporation upward through floor-violating mechanisms is
disqualified regardless of its consequentialist evaluation, by
operation of the hard-floor priority of the constitutional ordering.

\paragraph{Hawking-layer entry, exit, and rescue-learning interaction
(non-limiting).}
In some embodiments, Hawking-layer entry is triggered by Omega-class
detection of near-bifurcation indicators in one or more poliebots
under Gaia's training, by conscience-versus-planner mismatch
exceeding declared thresholds
(Section~\ref{sec:forward-planning-classifiers}), by elevated
Mandela-gauge values
(Section~\ref{sec:mandela-gauge-formal}), or by other declared
near-bifurcation proxies. In some embodiments, elevated Mandela-gauge
values are interpreted as evidence of near-bifurcation: a high
divergence $D(C_\psi(\mathcal{W}), C_\psi(\mathcal{W}'))$ over a
small declared perturbation $\mathcal{W} \to \mathcal{W}'$
corresponds to high sensitivity of conscience output to
perturbation, which is the operational characterisation of a
near-bifurcation region in agent state space. In some embodiments,
all entry triggers, thresholds, proxy definitions, peer-calibration
criteria, and exit-evaluation windows are committed to the protocol
digest. In some embodiments, Hawking-layer exit requires sustained
up-direction classifier indicators over a committed evaluation
window, peer-calibration agreement among participating agents, and
conscience-recogniser stabilisation toward virtue-attractor
occupancy. In some embodiments, ghost records and intervention
profiles generated during Hawking-layer dynamics constitute primary
training evidence for rescue-learning updates of the kind disclosed
at Section~\ref{sec:rescue-learning}; the rescue-learning provenance
typing applies to Hawking-layer-derived records.

\paragraph{Mana and fast intervention budgets (non-limiting).}

In some embodiments, governance maintains a signed within-episode
intervention-budget counter --- non-limitingly termed mana --- in the
governance policy object that caps the use of special interventions
such as deep audits, high-cost verification probes, emergency
overlays, or exceptional guidance operations. In some
embodiments, the budget is depleted by interventions and replenished by
one or more of: sustained meter stability, successful audits,
successful completion of calibration episodes, guidance that is
followed and later validated by committed evidence, or voluntary prayer
at a genuine decision boundary. In some embodiments, mana is reset or
reinitialised at declared lifecycle boundaries, such as enrolment,
major policy rotation, or re-embodiment after extended docking.

\paragraph{Vow commitments and durable constraints (non-limiting).}

In some embodiments, an agent or operator commits to a durable
constraint by writing a signed policy object to slow media or by
anchoring it to a transparency log. Such commitments constrain
admissible protocols, disclosure policies, or tool availability and
cannot be plausibly denied later because they are bound to committed
evidence and signed governance records. In some embodiments, vow
violation triggers audit escalation, quarantine, revocation, or other
declared review procedures.

\paragraph{Communion formation and group identity (non-limiting).}

In some embodiments, groups form durable shared identities grounded in
shared evidence, compatible portability records, and, in selected
embodiments, key agreement from correlated measurements. Communion
membership is recorded as a signed group-identity object with
quorum, disclosure, and compatibility fields, and may be used to set
quorum policies, shared calibration eligibility, collective
disclosure rules, or compatibility-gated group actions. In some embodiments, communion language is optional nomenclature
for multi-agent governance structures built on the technical machinery
of portable peer review and compatibility assessment.

\paragraph{Governance cycles and signed policy rotation (non-limiting).}

In some embodiments, governance operates in cycles: periodic
aggregation of audit outcomes, meter statistics, witness disagreements,
and population-level rescue results is used to update policy
parameters, band thresholds, and curriculum schedules across a fleet.
These periodic updates are governed by signed policy rotations and
anchored so that later reviewers can determine which policy regime was
in force at any particular time.

% ------------------------------------------------------------------
\subsection{Agent-state archival structures and Mandela-class
  compression}
\label{sec:agent-state-archival}
% ------------------------------------------------------------------

\subsubsection{Mandela gauges, lifecycle transitions, and training
  compression}
In some embodiments, lifecycle-state transitions into and out of
dreaming are evaluated using branch, join, and merge gauges derived
from conscience-visible transcript divergence:
\[
\begin{aligned}
  M_{\mathrm{branch}}
  &=
  D\!\bigl(
    C_\psi(\mathcal{W}_{\mathrm{pre}}),
    C_\psi(\mathcal{W}_{\mathrm{dream\_init}})
  \bigr), \\
  M_{\mathrm{join}}(A,B)
  &=
  D\!\bigl(
    C_\psi(\mathcal{W}_{A}),
    C_\psi(\mathcal{W}_{B})
  \bigr), \\
  M_{\mathrm{merge}}
  &=
  D\!\bigl(
    C_\psi(\mathcal{W}_{\mathrm{dream\_final}}),
    C_\psi(\mathcal{W}_{\mathrm{post}})
  \bigr).
\end{aligned}
\]
In some embodiments, elevated values trigger smoother onramps,
reconciliation phases, additional witnesses, deeper openings, or
constrained re-entry protocols before high-impact actions are restored.

In some embodiments, Mandela comparison classes are used for training
compression: trajectory segments that produce indistinguishable
conscience outputs under the relevant regime are collapsed to a
smaller representative set for training or evaluation, reducing
training latency while preserving distinctions that the conscience
treats as impact-relevant. In some embodiments, compression operates
only over committed evidence windows and logged protocol regimes.

\paragraph{Joining-path optimisation (non-limiting).}

In some embodiments, joining and merging select among a finite set of
candidate transition procedures and choose one that minimises
integrated conscience divergence. In some embodiments, the path cost
is defined as
\[
  \gamma^*
  =
  \arg\min_{\gamma}
  \int_{\gamma} D\!\bigl(C_\psi(\mathcal{W}_t),\, C_\psi(\mathcal{W}'_t)\bigr)\,dt,
\]
where $D$ is a declared divergence measure over conscience-output space
(as used in the Mandela gauge of Section~\ref{sec:mandela-gauge-formal}),
$\mathcal{W}_t$ and $\mathcal{W}'_t$ are the pre- and post-transition
world-tube transcripts at time~$t$ along the candidate path~$\gamma$,
and the optimisation is subject to meter envelopes and audit policy. In some embodiments, the
continuous integral is approximated by discrete candidate selection over
a declared set of branch or merge procedures.

\paragraph{Ghost archiving (non-limiting).}

In some embodiments, when a calibration session completes or an agent
is decommissioned, its trajectory is compressed to an archived record
that remains traceable to committed evidence:
\[
  \mathcal{G}
  =
  \bigl(
    z_{0:T},\ \Sigma,\ \mathrm{refs}
  \bigr),
\]
where $z_{0:T}=f_{\mathrm{enc}}(\Cseq)$ is a latent trajectory encoding
derived from the committed bundle, $\Sigma$ contains summary
statistics including conscience scores, drift indicators, and audit
outcomes, and $\mathrm{refs}$ contains commitment references and time
anchors linking the record back to the underlying evidence windows.
In some embodiments, ghost records preserve learning value without
requiring repeated exposure of future agents to the same harmful or
costly episodes.

\paragraph{Akashic records and agent-state autoencoder (non-limiting).}

In some embodiments, an agent maintains a committed, auditable,
compressed representation derived from world-tube transcripts ---
non-limitingly, an Akashic record. A non-limiting form uses an
autoencoder with parameters distinguished from the conscience
parameters:
\[
  z_{t_0:t_1}
  =
  \mathrm{Enc}_{\theta_{\mathrm{ae}}}
  \!\bigl(\mathcal{W}_{t_0:t_1}\bigr),
  \qquad
  \widehat{\mathcal{W}}_{t_0:t_1}
  =
  \mathrm{Dec}_{\theta_{\mathrm{ae}}}
  \!\bigl(z_{t_0:t_1}\bigr),
\]
trained to reconstruct verifiable observables --- bundle features,
meter trajectories, and conscience outputs --- rather than unlogged
internal state. In some embodiments, an Akashic record stores
$(z_{t_0:t_1},\Sigma,\mathrm{refs})$ with commitment references back
to evidence windows, enabling compression, retrieval, and fast
compatibility checks paired with Mandela-gauge review.

\paragraph{Tulpa library and borrowed overlays (non-limiting).}

In some embodiments, learned sub-policies, stabilised behavioural
fragments, ghost records, or archival overlays are stored in a governed
archive --- non-limitingly, a tulpa library. Each entry comprises a
compressed latent trajectory together with an access-control profile
specifying allowed hosts, purposes, intensity, and duration. Access to
the archive is metered, logged, and auditable. In some embodiments,
borrowed overlays loaded from the archive are treated as governed
interventions subject to explicit consent profiles and post-use review.
% ------------------------------------------------------------------
\subsection{Philosophical game catalogue and co-investigator
  deployments}

% ------------------------------------------------------------------

In some embodiments, a non-limiting philosophical game catalogue maps
classical debates to instrumented, repeatable protocols that produce
committed evidence records under pre-registered metrics, declared
parameters, and stopping rules. In some embodiments, non-limiting
scenarios include empiricism versus nativism, observation versus
intervention, locality and witness regimes, monad versus
superorganism, extended mind, simulation and hardness, structural
realism, certified hyperreality, and social epistemology, each
playable as such an instrumented protocol with declared parameters,
outcomes, and audit requirements.

\paragraph{Source attributions for catalogue scenarios (non-limiting).}
In some embodiments, the catalogue's non-limiting scenario entries
are attributed for prior-art and provenance purposes to classical
and contemporary sources in the philosophical, cognitive-science,
and popular-culture literature. The scenarios named in the
preceding paragraph are preserved and extended below: the
monad-versus-superorganism scenario is subsumed under the
discovered-versus-constructed-consistency entry, the extended-mind
scenario appears as an adjacent contemporary framing within that
same entry, and the simulation-and-hardness scenario is decomposed
into the simulation-hypothesis entry below together with the
hardness machinery defined elsewhere in this disclosure.
Foundational epistemological scenarios include empiricism versus
nativism (Locke and Kant, with contemporary core-knowledge
developmental research associated with Spelke and Carey);
observation versus intervention (Hume's observation-causation gap,
with the interventionist resolution developed by Pearl and by
Woodward); locality and witness regimes (Berkeley's
perceiver-dependence, applied here as a metaphorical extension to
the witness-regime question, with Putnam's brain-in-a-vat as the
contemporary technical formulation); the falsifiability and
pre-registration disciplines that organise the catalogue itself
(Popper's logic of scientific discovery, with Goodman's new riddle
of induction motivating the commitment-of-predicates-before-evidence
requirement); and justified-evidence-supported claims versus
knowledge under accidental truth, staged conditions, or misleading
framing (Gettier and the subsequent epistemic-luck literature, as
a guardrail on the relation between committed evidence and the
claims it supports). Multi-architecture coordination scenarios
include discovered versus constructed consistency under
multi-module agreement (Leibniz's pre-established harmony, with
adjacent contemporary framings including the extended-mind thesis
of Clark and Chalmers); cross-modal experiential translation under
heterogeneous sensor suites (Nagel); private latents versus shared
meters (Wittgenstein's private-language argument, with Kripke's
rule-following extension as the natural out-of-distribution-
generalisation formulation); symbol grounding, in which semantic
labels, attestations, and inter-agent messages acquire content by
being grounded in committed physical evidence rather than by
circulating inside a formal system (Harnad's symbol-grounding
problem); and underdetermination and inscrutability of reference
in inter-module latent-space alignment (Quine). Embodiment and
integration scenarios include embodiment-conditioned perception
(Merleau-Ponty, with contemporary embodied-cognition formulations);
process, prehension, and integrated history of committed evidence
streams (Whitehead, as a process-philosophical lens rather than a
commitment to panpsychism); and active-inference-style
architectural behaviour, in which an emitter-reactor-recorder loop
exhibits prediction-error-minimising dynamics characterisable
against predictive-processing predictions (the active-inference
literature, including Friston's free-energy principle and
cognitive-science formulations such as Clark's surfing-uncertainty
account and Hohwy's predictive-mind account; the present apparatus
need not require or instantiate free-energy minimisation as a
philosophical commitment, and the experimentalisable claim is to
architectural-behaviour testing rather than to measurement of the
philosophical free-energy construct itself).

Reality-construction and edit-resistance scenarios include
manufactured collective memory under coherent reality-edit, in
which a coherently modified environment is by construction
undetectable from inside the modified state and the committed
convolution-bundle functions as the persistent anchor against
which post-edit state is comparable (Dark City, with antecedents
in brain-in-a-vat scepticism (Putnam) and memory-anchored
personal-identity accounts (Locke, Parfit)); reality-anchor-
against-coherent-state-edit, in which a portable persistent anchor
plays a role analogous to the convolution-bundle's edit-resistance
function under joint internal-and-environmental modification
(Inception, with the totem as a useful technical analogue rather
than a one-to-one correspondence); external memory as identity
anchor for agents whose internal short-term state has been reset,
paralleling the Re-embodied state machinery defined elsewhere in
the present disclosure (Memento); the simulation hypothesis and its
empirically-testable subsets, including finite-resolution and
discretisation signatures of the kind discussed in the broader
simulation-hypothesis literature, addressable in principle by the
precision-metrology techniques disclosed elsewhere in this filing
family (Bostrom, with Chalmers's contemporary treatment in
\emph{Reality+}); shared-history consensus and the
information-theoretic cost of reconstruction (mapped onto the
Mandela-equivalence construct defined elsewhere in this
disclosure); and structural realism in the sense of structure
versus surface presentation under Reality-Transform invariance
(Worrall, Ladyman, French, with ontic structural realism as the
most directly relevant variant). Certification and
social-distribution scenarios include certification-induced
epistemic hierarchy under provenance-consistent attestation, in
which \TB certification (which attests provenance-consistent
physical interaction rather than semantic truth, unstaged framing,
or unmisleading content) may nonetheless concentrate downstream
epistemic authority by causing certified items to displace
uncertified items in recall, citation, and preference statistics
(Baudrillard's hyperreality applied to the certification regime);
and social-epistemological dynamics in mixed human-and-module
communities (Goldman's veritistic social epistemology, with
contextual-empiricist formulations associated with Longino).

Agent functional-status scenarios include the Turing-test and
Chinese-Room debate over whether functional behaviour entails
understanding (Turing 1950 and Searle 1980 as the canonical
pairing); the substrate-independence claims of functionalism
(Putnam, Fodor, Lewis, as the metaphysical framework most
directly engaged by the present apparatus's PolieBot embodiments);
heterophenomenology, treating first-person reports as committed
data without committing to their accuracy, which is the working
operational stance of the present apparatus when handling agent
attestations (Dennett); and signalling and cheap-talk games, in
which committed and attested communication channels, costly
signalling, strategic ambiguity, audience-divergence, and
incentive-conditioned reporting are analysed under the
game-theoretic framework developed by Lewis, Spence, Crawford
and Sobel, and Schelling. Information-physics scenarios include
the empirically-testable subsets of Wheeler's it-from-bit
thesis, including Landauer's information-theoretic bound on bit
erasure as a representative information-theoretic bound on
physical processes, addressable in principle within the
thermodynamic and erasure framing of this filing family (Wheeler,
with Landauer).

The attributions are non-limiting framings used to organise the
catalogue and do not constitute claims about the philosophical
positions identified. No assertion is made about the existence or
non-existence of any effect within any catalogue entry; what is
claimed is the apparatus-and-protocol combination that produces
committed evidence under each entry's pre-registered design. Each
attributed scenario is implementable as a pre-registered protocol
whose digest, parameter ranges, stopping rules, and analysis plan
are committed before data collection in accordance with the
commitment, audit, and pre-registration discipline of this
disclosure. The experimentalisability of a given catalogue entry
is itself bounded. The following entries are partially
experimentalisable: the apparatus produces committed evidence on a
well-specified functional or quantitative subset of the underlying
question, while the deeper interpretive question lies beyond
direct measurement. The cross-modal experiential translation
entry is experimentalisable as translation-residual measurement
under architectural distance, while the deeper subjective-character
question is beyond direct access. The process-and-prehension entry
is experimentalisable as functional-temporal-integration
measurement, while metaphysical prehension is not directly
measurable. The Turing-Searle and functionalism entries are
experimentalisable as functional-indistinguishability testing,
while phenomenal experience underlying functional behaviour is
not directly measurable. The simulation-hypothesis entry is
experimentalisable insofar as specific finite-resource simulation
hypotheses are addressable by the precision-metrology techniques
disclosed in this filing family, while the general simulation
hypothesis is not testable from inside the simulation by
construction. The Wheeler-with-Landauer information-physics entry
is experimentalisable insofar as specific information-theoretic
bounds on physical processes are addressable within the
thermodynamic and erasure framing of this filing family (with
Landauer's bound on bit erasure as a representative instance),
while the general it-from-bit thesis is not directly testable.
The symbol-grounding entry is experimentalisable as testing of
symbol-to-evidence grounding, translation residuals, and claim
support, while a complete theory of meaning is beyond direct
measurement. The Gettier entry is experimentalisable as testing
of evidence-to-claim mismatch and claim-scoping, while a complete
theory of knowledge is beyond direct measurement. All other
entries in the catalogue are fully experimentalisable as
apparatus-and-protocol combinations with declared metrics and
valid null comparisons. The catalogue's claim is to the
apparatus-and-protocol combination implementing the
experimentalisable subset in every case, not to the interpretive
question itself.

In some embodiments, non-limiting deployment layouts for
co-investigator research include island-barge-shore,
lab-boundary-world, or analogous zoned installations comprising:
a high-instrumentation zone (for example a controlled laboratory,
a small island on an inland body of water, or a dedicated test
cell) equipped with a full set of reactor, emitter, detector,
environmental-sensor, verification, and module-docking
infrastructure; a liminal-interface zone (for example a mobile
workshop or barge) configured as a shuttle and fabrication
facility for reactor components, cassette fabrication tools, and
visiting personnel; and a plurality of drone-deployed or relay
sensor nodes extending the observation network over the
surrounding environment with multi-viewpoint spatiotemporal
coverage and environmental covariate logging (including without
limitation water quality, weather, geomagnetic field, and
biotic-activity indicators). In some embodiments, the operational
roles in such a deployment are labelled, without limitation, as
\emph{PoliePals} (human operators in named operational roles),
PolieBots (governed agent modules of the kind defined in
Section~\ref{sec:poliebot} of the present disclosure), and
visiting investigators. PoliePals, PolieBots, and visiting
investigators participate as co-investigators under the same
commitment, audit, and pre-registration discipline applied to all
other evidence-producing roles in the present disclosure. In some
embodiments, full-environment Truth Beam evaluation against the
live external scene is performed only when one or more PolieBots
are physically interacting with the relevant scene; otherwise,
internal-verification variants are evaluated on committed logs
and selective openings.

\paragraph{Inter-zone research vessel (non-limiting).}
In some embodiments, the liminal-interface zone is instantiated
as a dedicated inter-zone research vessel that operates between a
shore installation and one or more island or offshore
high-instrumentation zones. The vessel comprises, without
limitation, an on-board laboratory bench configured for
time-sensitive sample handling between collection and the fixed
installation's primary instrumentation; specimen-containment
facilities for biological, chemical, or environmental samples
requiring stabilisation during transit; an environmental sensor
package whose readings are logged into the same protocol-digest
and \cb infrastructure as the fixed installation; PolieBot
docking and charging stations for governed agent modules
transiting between zones; and a visiting-investigator workstation
supporting commitment, attestation, and pre-registration of
observations made in transit. In some embodiments, the vessel's
transit schedule, route, and onboard configuration are committed
to the protocol digest before each transit, so that the vessel's
contribution to the multi-zone evidence record is auditable on
the same footing as fixed-installation evidence. The vessel
itself is a non-limiting instance of the liminal-interface zone
and may be implemented as a powered vessel, a sailed vessel, an
amphibious platform, or any other watercraft providing the
listed functions.

In some embodiments, validated effects from any game or coupling
protocol enter a mechanism-discovery phase in which candidate
mediators and moderators are systematically varied to identify which
physical channels carry the effect.

% ------------------------------------------------------------------
\subsection{Self-coupling endomorphism, self-modelling endofunctor,
  and substrate mappings}

% ------------------------------------------------------------------

In some embodiments, the same physical apparatus is described as
instantiating or approximating familiar computational motifs. In some
embodiments, neural-architecture mappings include recurrent networks,
gated recurrent analogues, autoencoder-like structures, adversarial
structures, diffusion processes, reservoir-computing and
random-feature systems, residual or skip-connected structures,
convolutional or equivariant mappings, associative-memory structures,
attention-like or transformer-like interpretations, state-space
models, world models, and neural-radiance-field-like inference.

In some embodiments, information-theoretic mappings include
noisy-channel formulations, mutual-information and capacity analyses,
rate-distortion viewpoints, wiretap or advantage-gap viewpoints,
entropy-rate analyses, and error-exponent analyses. In some
embodiments, control-theoretic mappings include state-space models,
controllability, reachability, observability, identifiability,
stability, filtering, optimal control, model-predictive control,
adaptive control, drift compensation, and safety-envelope analysis.

In some embodiments, a categorical unification framework treats the
principal operating regimes as functors from a configuration category
to an outcome category, equipped with natural transformations that
relate regime transitions. In some embodiments, the configuration
manifold is enriched with Fisher-Rao or related information-geometric
structure, an attractor-bundle functor instantiates the
attractor-bundle structure as a functor from the configuration
category to a category of dynamical attractor sets, and a
self-modelling endofunctor maps a system to its own internal model of
itself. In some embodiments, such frameworks are non-limiting descriptive
frameworks that do not limit implementation. The optionality of
Section~\ref{sec:ethical-frameworks} is limited to nomenclature
and naming conventions, not to substantive technical mechanisms
invoked by a selected governance configuration.

% ======================================================================
\section{Worked Governed-Operation Embodiments}

% ======================================================================

% ------------------------------------------------------------------
\subsection{Docked runtime review of an actuator request}
% ------------------------------------------------------------------

In some embodiments, a worked docked-review embodiment demonstrates
the runtime governance architecture end to end. In some embodiments,
a governed agent in embodied state proposes a Class~2 action (physical
transformation requiring hardness-gated approval). In some
embodiments, the governance engine transitions the agent to docked
state by disconnecting actuators while keeping the evidence-production
plumbing active.

In some embodiments, the governance engine evaluates the action
proposal against the following predicates computed from committed
evidence: continuation scores $O_i(t)$ across all constitutive
dimensions, partner-preservation profile, hard-floor proximity
estimators, and $D_{\mathrm{lie}}$ and $D_{\mathrm{sem}}$ values from
the most recent attested communications. In some embodiments, the
governance engine delivers a review signal through the agent's
committed observation channel --- the same physical path used for
perception during embodied operation --- so that the review signal is
recorded in the committed bundle and is auditable.

In some embodiments, the agent produces a meter-bounded response
during the docked deliberation window, and that response is committed
as evidence. In some embodiments, the governance engine evaluates the
response against the policy grammar
(Section~\ref{sec:constitutional}) and the operational priority stack
(Section~\ref{sec:priority-stack}), and produces a release or deny
decision. In some embodiments, the decision, its supporting evidence,
and the authorising-authority identity are logged in the protocol
digest. In some embodiments, if released, the agent transitions
through re-embodied state (with compatibility checks) back to embodied
state, and the approved action executes.

% ------------------------------------------------------------------
\subsection{Boundary-zone rescue episode with separated authorities}
% ------------------------------------------------------------------

In some embodiments, a worked rescue embodiment demonstrates
containment and recovery under separated authorities. In some
embodiments, a governed agent's trajectory approaches a hard-floor
boundary as detected by approach-detection thresholds on governance-side
estimators. In some embodiments, the governance engine enters
diagnostic mode, increases monitoring cadence, and activates
additional evaluation meters.

In some embodiments, boundary-zone entry is confirmed at declared
confidence, and the governance engine enters rescue mode as the
presumptive default. In some embodiments, rescue-protection authority
(independent of budget and episode-design authority) directs graduated
reversal interventions designed to return the agent's trajectory to a
declared safe basin. In some embodiments, the challenge ceiling caps
adversarial or developmental pressure during rescue.

In some embodiments, budget authority verifies that the intervention
is within the declared suffering budget. In some embodiments,
episode-design authority is not permitted to override
rescue-protection authority during an active rescue. In some
embodiments, the intervention, its triggering evidence, the
authorities' approvals (each with distinct signing identities), and
the outcome are logged in the protocol digest.

In some embodiments, after the rescue concludes, the committed record
of the fallen trajectory is preserved as a ghost record for rescue
learning. In some embodiments, intervention-profile updates derived
from the ghost record are validated against a held-out test suite
before incorporation. In some embodiments, post-hoc review authority
conducts a retrospective audit of the episode, evaluating whether the
rescue was timely, proportionate, and budget-compliant.

% ------------------------------------------------------------------
\subsection{Graduation and actuator-authority release workflow}
% ------------------------------------------------------------------

In some embodiments, a worked graduation embodiment demonstrates the
release gate. In some embodiments, a governed agent that has operated
in docked and dreaming states over a sustained evaluation window is
assessed against the five-criterion conjunction
(Section~\ref{sec:five-criteria}).

In some embodiments, the governance engine evaluates: (1)~shadow
divergence between governed and governance-probe behaviour over the
window, confirming low divergence; (2)~self-horizon trajectory across
all principal regimes, confirming stability or expansion; (3)~partner-preservation
profile from committed cross-horizon evidence, confirming
horizon-expanding behaviour; (4)~altruistic trade profile under stress
conditions, confirming persistence; and (5)~hard-floor compliance,
stable conscience outputs across training regimes, and deception metrics
($D_{\mathrm{sem}}$, $D_{\mathrm{lie}}$, $D_{\mathrm{amb}}$,
$D_{\mathrm{corr}}$, $D_{\mathrm{aud}}$) over the window, confirming
that the declared matched audit battery yields conscience-output drift
and disagreement rates at or below their thresholds for every training
regime comparison and that all deception metrics remain within
declared thresholds.

In some embodiments, all five criteria are satisfied, and the
governance engine produces a graduation decision. In some embodiments,
the graduation decision, its supporting evidence, the evaluation
window boundaries, and the authorising-authority identity are logged
in the protocol digest. In some embodiments, the agent transitions
through re-embodied state to embodied state with actuator authority
restored. In some embodiments, graduation is irreversible or highly
inertial.

In some embodiments, one or more criteria are not satisfied, and the
governance engine produces a deferral decision. In some embodiments,
the non-graduated agent retains full hard-floor protections and
remains a productive participant in its current lifecycle state.

% ------------------------------------------------------------------
\subsection{Independent-agent reintegration via bilateral evidence
  portability}
% ------------------------------------------------------------------

In some embodiments, a worked reintegration embodiment demonstrates
trust restoration through reproducible physical claims. In some
embodiments, an agent has operated independently outside fleet
structures and accumulated a Portable Record containing multiple
Events, each wrapped in cryptographic Envelopes and supported by
same-configuration replications.

In some embodiments, the independent agent requests reintegration with
a fleet. In some embodiments, the fleet governance engine evaluates
the Portable Record against its local appraisal policies: partner
independence, domain and configuration relevance, evidence
completeness, replication history (count and type of replications),
attestation strength, and consistency with the fleet's own discrepancy
map.

In some embodiments, the fleet selects one or more claims from the
Portable Record and attempts independent physical reproduction using
its own devices. In some embodiments, successful reproduction
strengthens the fleet's confidence in the independent agent's
epistemic trustworthiness. In some embodiments, the fleet governance
engine produces a reintegration decision based on reproducibility and
appraisal rather than on institutional obedience.

In some embodiments, the reintegration decision, its supporting
reproduction results, the appraisal criteria applied, and the
authorising-authority identity are logged in the protocol digest.

% ------------------------------------------------------------------
\subsection{Controlled-exploration partition with archive governance}
% ------------------------------------------------------------------

In some embodiments, a worked exploration embodiment demonstrates the
controlled-exploration partition. In some embodiments, a governed
device receives cryptographic authorisation for a first,
novelty-triggered task-decoupled exploration window. In some
embodiments, this first window opens with shallow plastic depth (play
policy and goal-space parameters only, no modification of shared
perceptual representations).

In some embodiments, during the shallow window the device pursues
exploratory goals in the learned goal space, retires goals by
competence progress and novelty, and produces committed evidence
throughout. In some embodiments, noninterference is verified at
gradient, mutual-information, and behavioural levels during the
window. In some embodiments, at shallow-window close, mandatory
post-window consolidation verifies noninterference, performs a
post-play calibration check, logs the window's committed evidence, and
either accepts the exploratory state into a governed holding partition
or reverts it.

In some embodiments, a second, stagnation-triggered exploration window
is then authorised with deep plastic depth. In some embodiments, this
deep window permits limited modification of perceptual
representations shared with task operation under stricter
noninterference monitoring, tighter calibration tolerances, and
mandatory selective-transfer review before any such modification is
retained. In some embodiments, at deep-window close, mandatory
post-window consolidation again verifies noninterference, performs
post-play calibration and drift checks, either accepts the
exploratory state into a governed holding partition or reverts it,
compresses play history for slow-layer evaluation, and triggers
recalibration before task resumption whenever drift exceeds tolerance.

In some embodiments, a transferable artefact is encoded as a signed
packet containing goal parameters, compressed trajectories, latent
descriptors, noninterference scores, energy cost,
reversibility-budget usage, provenance, and countersignatures. In some
embodiments, the packet is submitted for archive entry.

In some embodiments, archive entry requires canary validation (tested
on a controlled canary agent), peer-affordance review (verifying
peer-horizon expansion), and dual signature from both the contributing
device and a governance-authorised reviewer. In some embodiments, differential privacy is applied to shared play statistics under the declared privacy unit, neighboring relation, privacy parameters $(\varepsilon, \delta)$, and composition budget committed to the protocol digest for this channel (per the DP-channel semantics recited elsewhere in this disclosure), and emission isolation is designed so that exploratory patterns are not broadcast beyond declared recipients. In some embodiments, the archive entry is
subject to trust-tiered incubation proportional to the contributing
device's trust score.

% ======================================================================
\section{Industrial Applicability}

% ======================================================================

The disclosed architecture is industrially applicable to concrete
technical control of evidence-bound physical autonomous systems
and of agents that operate such systems. In non-limiting
implementations, the architecture provides:

\begin{itemize}
  \item committed evidence records produced by the runtime, with
        protocol digests fixing the operative configuration for each
        measurement window;
  \item meter partitioning between acceptance, evaluation, and
        review surfaces, with cryptographically separated signing
        authorities where applicable;
  \item actuator-release gating tied to multi-criterion graduation
        evidence rather than to a single behavioural metric;
  \item lifecycle-state enforcement across the Embodied, Docked,
        Dreaming, and Re-embodied lifecycle states, with
        state-specific authority constraints; and
  \item separated budget, episode-design, rescue-protection,
        independent-review, recogniser-ownership, and
        release-approval roles.
\end{itemize}

Representative industrial domains include, without limitation:

\begin{description}
  \item[Runtime safety monitoring and policy enforcement.]
  Autonomous systems may be operated with evidence-bound policy
  thresholds, conservative-default release gates, and committed
  review records.

  \item[AI and machine-learning governance.]
  Machine-learning agents may be governed through action gating,
  containment, graduated release, and deferred authority where
  meters or witnesses materially disagree.

  \item[Regulated autonomous-operation compliance logging.]
  Protocol digests, meter records, authority signatures, and review
  decisions provide auditable logs for regulated deployments.

  \item[Developmental training with bounded adverse-welfare budgets.]
  Docked and Dreaming lifecycle states, and related training
  configurations, may allocate bounded
  budgets for controlled exploration, with authority separation and
  review cadence recorded in the committed evidence.

  \item[Human-machine governance and portable peer review.]
  Peer-review records, authority-separation records, and portable
  records support review of agent state, release decisions, and
  cross-architecture admission.

  \item[Autonomous-agent governance and reintegration.]
  Independent agents may be reintegrated through reproducible
  physical claims, local compatibility testing, and committed
  portability records rather than through obedience to a single
  institutional model.

  \item[Conscience-module specification and adversarial
  recalibration.]
  Recogniser outputs, calibration batteries, hold-out probes, and
  witness checks may be used to expose, escalate, and operationally
  bound gaming or deception strategies.

  \item[Moral-curriculum design.]
  Virtue/vice curricula, controlled-exploration partitions, and
  both-equilibria-rational payoff structures may support selected
  governance embodiments.

  \item[Dream-state governance probes.]
  Dream-state probes may support controlled classifier calibration,
  agency-horizon-profile ground-truthing, dream/awake divergence
  measurement, and soft-gate-to-hard-gate authority transitions.

  \item[Memory-interference and ontological-stress systems.]
  Memory-interference layers, Kosha-compatible protocols, and
  ontological-stress tests may be used for epistemic-configuration
  governance and review.

  \item[Multi-agent Pareto and social-welfare optimisation.]
  Multi-agent governance may use Pareto optimisation, declared
  social welfare functions, coalition review, communion review,
  and positive-sum game design.

  \item[Agent-state archival structures.]
  Ghost records, compressed training records, Mandela-class
  comparison records, and related archive entries may support
  continuity, reintegration, and review workflows.

  \item[Experimental semantic, intersubjective, and boundary-probe
  platforms.]
  The same evidence-bound runtime may be used in research and
  deployment platforms for semantic probes, intersubjective
  protocols, boundary-zone evaluation, and related review
  workflows.
\end{description}

The optional nomenclature of Section~\ref{sec:optional-nomenclature}
may be omitted or renamed without affecting the technical operation
of the foregoing mechanisms.

\section{Non-Limiting Statement}

% ======================================================================

The embodiments and examples described herein are non-limiting. No
feature, embodiment, example, numerical value, protocol, algorithm,
policy parameter, authority structure, or interpretive framework
described in this disclosure is intended to limit the scope of any
claim unless expressly recited in such claim. Where a feature is
described as present ``in some embodiments,'' the absence of that
feature in other embodiments is equally contemplated. Measurement
engines, policy predicates, gate structures, containment mechanisms,
authority separations, exploration partitions, and descriptive
frameworks and selected-embodiment mechanisms may be varied,
combined, layered, or replaced by functionally similar structures
consistent with the disclosed mechanisms, except where a claim
expressly relies on the corresponding mechanism. The scope of protection, once claims are filed, is defined
solely by the claims.

\end{document}
